
\documentclass[11pt]{article}
\usepackage[utf8]{inputenc}
\usepackage[T1]{fontenc}
\usepackage[a4paper,margin=1in]{geometry}
\usepackage{lmodern}
\usepackage{microtype}
\usepackage{amsmath,amssymb,mathtools,amsthm}
\usepackage{mathrsfs}
\usepackage{enumitem}
\usepackage{longtable,booktabs,array}
\usepackage{tikz-cd}
\usepackage[hidelinks,hypertexnames=false]{hyperref}
\setlength{\parindent}{0pt}
\setlength{\parskip}{0.55em}
\emergencystretch=3em
\hbadness=10000
\providecommand{\Spec}{\operatorname{Spec}}
\providecommand{\Sch}{\operatorname{Sch}}
\providecommand{\Gal}{\operatorname{Gal}}
\providecommand{\GL}{\operatorname{GL}}
\providecommand{\Aut}{\operatorname{Aut}}
\providecommand{\Isom}{\operatorname{Isom}}
\providecommand{\Pic}{\operatorname{Pic}}
\providecommand{\Tor}{\operatorname{Tor}}
\providecommand{\Ker}{\operatorname{Ker}}
\providecommand{\Im}{\operatorname{Im}}
\providecommand{\cO}{\mathcal O}
\providecommand{\Z}{\mathbb Z}
\providecommand{\Q}{\mathbb Q}
\providecommand{\Ql}{\mathbb Q_\ell}
\providecommand{\Zl}{\mathbb Z_\ell}
\providecommand{\Fp}{\mathbb F_p}
\providecommand{\Gm}{\mathbb G_m}
\providecommand{\A}{\mathbb A}
\providecommand{\Pj}{\mathbb P}
\providecommand{\et}{\mathrm{et}}
\providecommand{\prof}{\operatorname{prof}}
\providecommand{\cd}{\operatorname{cd}}
\providecommand{\lcm}{\operatorname{lcm}}
\providecommand{\rank}{\operatorname{rank}}
\providecommand{\pr}{\operatorname{pr}}
\providecommand{\Sp}{\operatorname{Sp}}
\providecommand{\limproj}{\varprojlim}
\providecommand{\liminj}{\varinjlim}
\providecommand{\eps}{\varepsilon}
\providecommand{\PhiC}{\Phi}
\providecommand{\PsiC}{\Psi}
\newtheorem*{theorem*}{Theorem}
\newtheorem*{proposition*}{Proposition}
\newtheorem*{lemma*}{Lemma}
\newtheorem*{corollary*}{Corollary}

% SGA7 batch 099 macro additions
\providecommand{\Hom}{\operatorname{Hom}}
\providecommand{\ob}{\operatorname{ob}}
\providecommand{\Ens}{\mathrm{Ens}}
\providecommand{\Set}{\mathrm{Set}}
\providecommand{\pro}{\operatorname{pro}}
\providecommand{\Spf}{\operatorname{Spf}}
\providecommand{\Ext}{\operatorname{Ext}}
\providecommand{\Ccat}{\mathcal C}
\providecommand{\Acat}{\mathcal A}
\providecommand{\Fcat}{\mathcal F}
\providecommand{\Gcat}{\mathcal G}
\providecommand{\Mcat}{\mathcal M}
\providecommand{\Vcat}{\mathcal V}
\providecommand{\m}{\mathfrak m}
\providecommand{\ma}{\mathfrak m_\Lambda}
\providecommand{\mh}{\widehat{\mathfrak m}}
\providecommand{\ka}{k[\varepsilon]}
\providecommand{\CL}{\mathcal C_\Lambda}
\providecommand{\CLK}{\mathcal C_{\Lambda,K}}
\providecommand{\D}{\mathrm D}
\providecommand{\Isom}{\operatorname{Isom}}
\providecommand{\id}{\operatorname{id}}
\providecommand{\alg}{\operatorname{alg}}

\begin{document}

\begin{center}
{\Large SGA 7-I: Monodromy Groups in Algebraic Geometry}\\[0.4em]
{\large Source-grounded English translation, original scan pages 1--24}\\[0.4em]
Directed by A. Grothendieck, with the collaboration of M. Raynaud and D. S. Rim.
\end{center}

\section*{Title page}
\begin{center}
Seminar on Algebraic Geometry at Bois-Marie, 1967--1969\\[0.3em]
\textbf{Monodromy Groups in Algebraic Geometry}\\
(SGA 7 I)\\[0.8em]
Directed by A. Grothendieck\\
with the collaboration of M. Raynaud and D. S. Rim.
\end{center}

\section*{Introduction}
Let $S$ be a scheme and let $f:X\to S$ be a morphism of schemes. If $f$ is proper and smooth, and if the prime number $\ell$ is invertible on $S$, the $\ell$-adic cohomology groups
\[
 H^i(X_{\bar s},\Zl)
\]
of the geometric fibres of $f$ form a local $\ell$-adic system on $S$. For $f$ assumed only proper, these groups are the fibres of an $\ell$-adic sheaf $R^if_*\Zl$ on $S$. The theory of vanishing cycles relates the ramification of this sheaf on $S$ to the singularities of $f$.

We shall consider only the case where $S$ is a henselian trait, i.e. the spectrum of a henselian discrete valuation ring. This is in practice the essential case. I refer to (I 2.1) for a heuristic description of the theory. Let $S=\Spec(V)$, let $k(\bar\eta)$ be an algebraic closure of the field of fractions $k(\eta)$ of $V$, and let $k(\bar s)$ be the corresponding algebraic closure of the residue field $k(s)$; thus $k(\bar s)$ is the residue field of the normalization $\bar V$ of $V$ in $k(\bar\eta)$. In the particular case where $X$ is proper and flat over $S$, of relative dimension $n$, and where $f$ has only one point $x\in X_s$ at which it is not smooth, one defines $\Gal(\bar\eta/\eta)$-modules $\Phi^i$, purely local in nature in a neighbourhood of $x$, zero for $i\notin[0,n]$ (I 4.2), and one constructs a long exact sequence of $\Gal(\bar\eta/\eta)$-modules
\[
 \cdots\longrightarrow H^i(X_s,\Zl)\xrightarrow{\operatorname{sp}}
 H^i(X_{\bar\eta},\Zl)\longrightarrow \Phi^i\longrightarrow
 H^{i+1}(X_s,\Zl)\longrightarrow\cdots .
\]
The $\Gal(\bar s/s)$-module $H^i(X_s,\Zl)$ is regarded as a $\Gal(\bar\eta/\eta)$-module with trivial inertia action; $\operatorname{sp}$ is the specialization map. Criteria are also given, at present only in characteristic $0$, for the $\Phi^i$ with $i$ small to vanish; for example, if $X$ is moreover locally a complete intersection, then $\Phi^i=0$ for $i\ne n$ in the situation of I 4.5.

The monodromy theorem asserts that the action of the inertia subgroup $I$ of $\Gal(\bar\eta/\eta)$ is quasi-unipotent: for $T\in I$, there exist integers $N>0$, $M>0$ such that the endomorphism
\[
 (T^M-I)^N
\]
of $H^i(X_{\bar\eta},\Zl)$ is zero. Two proofs of this theorem are given. The first (I 1.2), arithmetic in nature, applies as soon as the cohomology groups considered have reasonable finiteness properties. The key point is that, when $k(s)$ is of finite type over its prime subfield, the action of $I$ is quasi-unipotent for every $\ell$-adic representation of $\Gal(\bar\eta/\eta)$. The second proof, more geometric, requires purity and resolution of singularities; it is at present valid only in characteristic $0$ or for $H^1$. It gives useful information about the exponent of nilpotence $N$ ($N\leq i+1$ for $H^i$), and, as will be seen later, lends itself well, over $\mathbb C$, to comparison with transcendental theory.

Applied to $H^1$ of abelian varieties, i.e. to their Tate modules, these results make it possible to study reduction modulo $p$. Thus two proofs are given of the stable reduction theorem for abelian varieties: after ramification, the connected special fibre of the Neron model is an extension of an abelian variety by a torus. The first proof (I 6) takes up again the arithmetic method of the monodromy theorem. The second (IX 3.6) rests on a much finer analysis of the Neron model and of its polarization properties.

Grothendieck's Exposes I to V were not written up. They have been summarized in Expose I. The results stated there are proved succinctly, but essentially completely. Expose II applies the method of Lefschetz pencils and (I 5.3) to the study of the fundamental group. For $S$ a surface over an algebraically closed field $k$, of characteristic exponent $p$, one shows that the profinite group $\pi_1^{(p')}(S,s)$, completed away from $p$, is of finite pro-$(p')$ presentation. As explained above, Exposes III to V do not exist.

Expose VI contains, with some complements, Schlessinger's deformation theory. Care is taken there to account for infinitesimal automorphisms of the objects being classified, and not to pass too abruptly to the functor of isomorphism classes of objects. A new construction is also given of
\[
 \operatorname{Ext}^1(L_{X/S},-),
\]
where $L_{X/S}$ is the relative cotangent complex. This expose serves in the rest of the seminar especially through the study made there of deformations of ordinary quadratic singularities.

Exposes VII and VIII are devoted to the theory of biextensions. This theory plays an essential role in Expose IX, to express what happens to a polarization when one passes from an abelian variety to its Neron model. Expose IX contains the proof of the stable reduction theorem for abelian varieties and various applications. The sequel to this seminar, SGA 7 II, by P. Deligne and N. Katz, will appear later.

\begin{flushright}
Bures-sur-Yvette, May 1972,\\
P. Deligne
\end{flushright}

\section*{Table of contents of SGA 7 I}
\begin{longtable}{@{}p{0.72\linewidth}r@{}}
Expose I. Summary of the first exposes of A. Grothendieck, written up by P. Deligne & 1\\
Expose II. Finiteness properties of the fundamental group, by Michele Raynaud & 25\\
Expose VI. Formal deformation theory, by D. S. Rim & 32\\
Expose VII. Local monodromy groups, by A. Grothendieck & 133\\
Expose VIII. Complements on biextensions; general properties of biextensions of group schemes, by A. Grothendieck & 218\\
Expose IX. Neron models and monodromy, by A. Grothendieck & 313
\end{longtable}

\section*{Expose I. Summary of the first exposes of A. Grothendieck}
\begin{center}
written up by P. Deligne
\end{center}

\subsection*{Contents}
0. Preliminaries. 1. Arithmetic proof of the monodromy theorem. 2. Vanishing cycles. 3. Geometric proof of the monodromy theorem. 4. Vanishing criteria for sheaves of vanishing cycles. 5. Action of monodromy on the $\pi_1$. 6. Appendix by P. Deligne: arithmetic proof of the stable reduction theorem.

\subsection*{0. Preliminaries}
\subsubsection*{0.0. Terminology}
We recall the following definitions.

(0.0.1) A \emph{trait} is the spectrum of a discrete valuation ring. Its generic and closed points are often denoted by $\eta$ and $s$.

(0.0.2) A ring is \emph{strictly local} if it is henselian local with separably closed residue field. A scheme is strictly local if it is the spectrum of such a ring.

(0.0.3) A geometric point of $S$ in this expose is a morphism $\bar x:\Spec(k)\to S$ with image $x\in S$, where $k(x)\to k$ and $k$ is a separable closure of $k(x)$. For a henselian trait $S=\Spec(V)$, one often writes $\bar\eta$ for a generic geometric point of $S$. The residue field of the valuation ring $\bar V$, the integral closure of $V$ in $k(\bar\eta)$, is a separably closed algebraic extension of $k(s)$ and defines a geometric point $\bar s$ of $S$ localized at $s$.

(0.0.4) The strict localization of $X$ at the geometric point $\bar x$ is as in SGA 4 VIII 4.4.

(0.0.5) A local, henselian local, or strictly local $S$-scheme is said to be essentially of finite type over $S$ if it is the spectrum of a local ring, respectively the henselization of such a local ring, respectively the strict localization, of a scheme of finite type over $S$.

(0.0.6) $F\mapsto F(n)$ denotes the Tate twist.

(0.0.7) An endomorphism $T$ of a finite-dimensional vector space is quasi-unipotent if its eigenvalues are roots of unity, i.e. if there exist $N\geq1$ and $M\geq1$ such that $(T^N-1)^M=0$.

\subsubsection*{0.1. Reminders on $\ell$-adic cohomology}
Let $X$ be a scheme of finite type over an algebraically closed field $k$ of characteristic exponent $p$, and let $\ell$ be a prime prime to $p$.

The $\ell$-adic cohomology groups of $X$ were defined in SGA 4 and SGA 5:
\begin{enumerate}[label=\alph*)]
\item $H^i(X,\Z/\ell^n)$ is the $i$-th cohomology group of the etale site $X_{\et}$ with values in $\Z/\ell^n$;
\item $H^i(X,\Zl)=\limproj_n H^i(X,\Z/\ell^n)$;
\item $H^i(X,\Ql)=H^i(X,\Zl)\otimes_{\Zl}\Ql$.
\end{enumerate}
The $\ell$-adic cohomology groups with proper supports are defined similarly from the groups $H_c^i(X,\Z/\ell^n)$ of SGA 4 XVII.

In each of the following cases one can prove that these groups $H^i(X)$ have reasonable finiteness properties: cohomology with proper support; $X$ proper over $k$; $X$ smooth over $k$ (using duality and Poincare duality); when Hironaka resolution of singularities is available for schemes of finite type over $k$ of dimension at most $\dim X$; and for $i=0,1$.

In all these cases, the groups $H^i(X,\Z/\ell^n)$ are finite, the groups $H^i(X,\Zl)$ are of finite type, the groups $H^i(X,\Ql)$ are finite-dimensional, and, except perhaps in the last case, one has split short exact sequences
\[
0\to H^i(X,\Zl)\otimes \Z/\ell^n\to H^i(X,\Z/\ell^n)\to
\Tor_1\bigl(H^{i+1}(X,\Zl),\Z/\ell^n\bigr)\to0.
\]
Moreover, the proof of the finiteness theorem gives each time a proof of generic constructibility: if $k$ is the field of fractions of an integral noetherian scheme $S$ and if $X$ is the generic fibre of $f:X_1\to S$ of finite type, there exists a nonempty open $U\subset S$ such that, over $U$, the sheaves $R^if_*\Z/\ell^n$ are locally constant and have formation compatible with every base change.

\subsubsection*{0.2. Geometric cohomology}
We no longer suppose $k$ algebraically closed. Let $\bar k$ be an algebraic closure (or separable closure, which amounts to the same thing). We consider the geometric cohomology groups $H^i(X_{\bar k})$, where $X_{\bar k}$ is obtained from $X$ by extending scalars from $k$ to $\bar k$. By transport of structure, $\Gal(\bar k/k)$ acts on these groups. In $\Ql$-cohomology, in the cases (0.1.1) to (0.1.5), one obtains continuous representations of $\Gal(\bar k/k)$ in a linear group $\GL(n,\Ql)$.

\subsubsection*{0.3. Reminders on traits}
Let $S$ be a henselian trait, with notation $\eta,s,V,\bar V,\bar\eta,\bar s$ as in (0.0.1) and (0.0.3). Let $p$ be the characteristic exponent of $k(s)$. The map $\Gal(\bar\eta/\eta)\to\Gal(\bar s/s)$ has kernel the inertia group $I$. The subfield of $k(\bar\eta)$ defined by $I$ is the maximal unramified extension of $k(\eta)$.

The profinite group $I$ has a largest pro-$p$ subgroup $P$. The quotient $I/P$ is the tame inertia group, and canonically
\[
 I/P\simeq \limproj_{(n,p)=1}\mu_n(k(\bar s))\simeq \widehat{\Z}^{(p')}(1)(k(\bar s)).
\]
If $\chi_n:I/P\to\mu_n(k(\bar s))$ is the canonical map and if $u$ is an element of valuation $1/n$, then for $\sigma\in I$,
\[
 \sigma(u)=\chi_n(\sigma)\,u\pmod{\text{elements of valuation } >1/n}.
\]
In summary,
\[
\Gal(\bar\eta/\eta)\supset I\supset P,
\quad \Gal(\bar\eta/\eta)/I\simeq \Gal(\bar s/s),
\quad I/P\simeq \widehat{\Z}^{(p')}(1)(k(\bar s)),
\]
where $P$ is a pro-$p$ group (trivial if $p=1$), and conjugation by $\Gal(\bar\eta/\eta)/I$ on $I/P$ is the natural action of $\Gal(\bar s/s)$ on $\widehat{\Z}^{(p')}(1)(k(\bar s))$.

\subsubsection*{0.4. Abhyankar's lemma}
More generally, let $S$ be strictly local regular, let $D$ be a regular divisor in $S$, and let $p'$ be the characteristic exponent of the generic point of $D$. Let $\bar\eta$ be a geometric generic point of $S$ and let $\bar s$ be the geometric point localized at the closed point $s$ which is deduced from it. By Abhyankar's lemma the fundamental group $\pi_1(S-D,\bar\eta)$ has an analogous devissage:
\[
\pi_1(S-D,\bar\eta)\supset I\supset P,
\quad \pi_1(S-D,\bar\eta)/I\simeq \Gal(\bar s/s),
\quad I/P\simeq \widehat{\Z}^{(p')}(1)(k(\bar s)),
\]
where $P$ has no quotient of order prime to $p'$. The action by inner automorphisms of $\pi_1(S-D,\bar\eta)/I$ on $I/P$ is the natural action of $\Gal(\bar s/s)$ on $\widehat{\Z}^{(p')}(1)(k(\bar s))$.

\subsubsection*{0.5. Neron desingularization}
The method of Neron desingularization gives the following result.

\textbf{Lemma 0.5.1.} Let $S\to S_1$ be a morphism of henselian traits, with $S=\Spec(V)$, $S_1=\Spec(V_1)$, and let $\varpi$ be a uniformizer of $V_1$. Suppose that $\varpi$ is also a uniformizer of $V$, and that the extensions $k(s)/k(s_1)$ and $k(\eta)/k(\eta_1)$ are separable. Then $V$ is the filtered inductive limit of henselian $V_1$-algebras, smooth and essentially of finite type over $V_1$.

To construct the desired $V_1$-algebras, one takes $V_1\subset B\subset V$ with $B$ of finite type, and applies Neron desingularization to $\Spec(B)$.

Let $S=\Spec(V)$ be a henselian trait. If $S$ is of equal characteristic, with uniformizer $\varpi$, (0.5.1) applies to $S_1=\mathbb F_p[\varpi]_{(\varpi)}$ after the usual separability check; hence $S$ is a limit of spectra of henselian algebras, smooth and essentially of finite type over a trait of finite type. If $S$ is complete of mixed characteristic, with perfect residue field $k$, then $V$ is an Eisenstein extension of the Cohen-Witt ring $W(k)$; after a small perturbation of the coefficients one reduces to a trait whose residue field is a radical extension of a field of finite type over $\mathbb F_p$, and again (0.5.1) applies.

\subsection*{1. Arithmetic proof of the monodromy theorem}
The reader will find in [5], Appendix, a proof of the following result of A. Grothendieck.

\textbf{Proposition 1.1.} Let $S$ be a henselian trait, let $\ell$ be a prime prime to the characteristic exponent $p$ of $k(s)$, and let $\rho$ be a continuous representation of $\Gal(\bar\eta/\eta)$ in $\GL(n,\Ql)$. Suppose that no finite extension of $k(s)$ contains all roots of unity of order a power of $\ell$. Then there exists an open subgroup $I_1$ of the inertia group $I$ such that $\rho(\sigma)$ is unipotent for $\sigma\in I_1$.

The preceding condition is automatically satisfied when $k(s)$ is of finite type over the prime field, or purely inseparable over such a field.

\textbf{Monodromy theorem 1.2.} With the preceding notation, let $X$ be a scheme of finite type over $\eta$, and let $H$ be a cohomology space of $X_{\bar\eta}$, with coefficients in $\Ql$, of one of the types (0.1.1)--(0.1.5). If the condition of Proposition 1.1 is satisfied, the action of any element of $I$ on $H$ is quasi-unipotent.

The limit arguments of 0.5 allow one to remove this hypothesis.

\textbf{Variant 1.3.} The condition of Proposition 1.1 is superfluous.

Let $H'$ be a $\Zl$-cohomology space of $X_{\bar\eta}$ which gives rise to $H$, and put $H_1=H'/\text{torsion}$. One has $H=H_1\otimes_{\Zl}\Ql$. After replacing the trait by a finite extension, one may suppose that $\Gal(\bar\eta/\eta)$ acts trivially on $H_1/\ell H_1$, and, in mixed characteristic, that the trait is complete with perfect residue field. Using the presentation of the trait as a limit from 0.5, the representation comes from a constructible constant-twisted sheaf on $S_i-D_i$, and the devissages (0.3.1) and (0.4.1) give commutative diagrams
\[
\begin{tikzcd}
\Gal(\bar\eta/\eta) \arrow[r,two heads] \arrow[d] & I \arrow[r,two heads] \arrow[d] & P \arrow[d]\\
\pi_1(S_i-D_i,\bar\eta) \arrow[r,two heads] & I_i \arrow[r,two heads] & P_i
\end{tikzcd}
\qquad
\begin{tikzcd}
I/P \arrow[d] \arrow[r,"\sim"] & \widehat\Z^{(p')}(1) \arrow[d,equal]\\
I_i/P_i \arrow[r,"\sim"] & \widehat\Z^{(p')}(1).
\end{tikzcd}
\]
Since $\Ker(\GL(H_1)\to \GL(H_1/\ell H_1))$ is a pro-$\ell$ group, $P_i$ acts trivially, and one applies the following variant.

\textbf{Variant 1.4.} With the notation of 0.4, let $\rho$ be a continuous representation of $\Gal(\bar\eta/\eta)$ in $\GL(n,\Ql)$. Suppose that $k(s)$ satisfies the condition of 1.1 and that $\rho(P)$ is finite. Then the action of $I$ is quasi-unipotent.

\subsection*{2. Vanishing cycles}
Let $S$ be a strictly local trait, and keep the notation of 0.3; here $s=\bar s$. Let $f:X\to S$ be a scheme of finite type over $S$. The formalism of vanishing cycles is a tool for studying the relation between the cohomology of $X_s$ and that of $X_{\bar\eta}$, and the inertia action on the latter, from local information on $X$ and on the morphism $f$.

In transcendental theory, if $f:X\to D$ is a proper morphism from an analytic space to a disk, the vanishing cycles appear as follows. After shrinking $D$, $X$ is a topological fibration over the punctured disk $D^*$, so the fibres $X_t$ for $t\in D^*$ are all isomorphic. One represents the special fibre $X_0$ as obtained from the general fibre by contraction of polyhedra:
\[
\begin{tikzcd}[column sep=large,row sep=large]
X_t \arrow[rr,"\text{contraction}"] \arrow[dr] && X_0 \arrow[dl]\\
& D &
\end{tikzcd}
\]
The method of vanishing cycles then studies the difference between the cohomology of $X_t$ and that of $X_0$ by means of the Leray spectral sequence of this contraction.

The geometric theory below is very close to this transcendental theory: the general point $t\in D$ is replaced by $\bar\eta$, the geometric generic point of $S$. Conversely, the transcendental theory can be modelled on the geometric one by replacing $t$ or $\bar\eta$ by the universal covering of $D^*$.

\textbf{Proposition 2.3.} Let $u:X'\to X$ be etale, $x'$ a geometric point of $X'$ over the geometric point $x$ of $X_s$. The homomorphism
\[
 (R^i\Psi)_x\longrightarrow (R^i\Psi)_{x'}
\]
induced by $u$ is an isomorphism.

\textbf{Corollary 2.4.} The formation of the $\Psi^i$ and $\Phi^i$ commutes with strict localization at a point $x$.

Let $A$ be the coefficient ring. If $I$ acts on the geometric generic fibre, one has the exact sequence of sheaves
\[
0\to A_{X_s}\to R\Psi A\to R\Phi A\to 0
\]
in the derived sense; for $f$ proper this gives the long exact sequence
\[
\cdots\to H^i(X_s,A)\to H^i(X_{\bar\eta},A)\to H^i(X_s,R\Phi A)\to H^{i+1}(X_s,A)\to\cdots .
\]

\subsection*{3. Geometric proof of the monodromy theorem}
Let $S$ be strictly local, $X$ a normal excellent scheme regular away from a divisor $D$, and suppose that, etale locally, the special fibre is a divisor with normal crossings
\[
X_s=\sum_{i\in C} a_iD_i.
\]
The proof uses purity to describe the prime-to-$p$ coverings of the complement and then takes limits. If $x$ is a point of the special fibre and if $C$ is the set of components through $x$, set
\[
K: \bigoplus_{i\in C}\Z(1)\xrightarrow{(a_i)}\Z(1),
\]
concentrated in degrees $-1,0$.

\textbf{Theorem 3.3.} Assuming purity, the groups of tame vanishing cycles are described by
\[
 R^*\Psi_x \simeq H^*(K\otimes A)\otimes A[t],\qquad \deg(t)=2,
\]
and the inertia action on $R^1\Psi_x$ is obtained from a transitive action on a finite set of cardinal prime to $p$, determined by the multiplicities $a_i$.

\textbf{Corollary 3.4.} Let $N$ be the least common multiple of the multiplicities of the components of the special fibre. If at most $q'$ components meet at a point, then for $T$ in tame inertia
\[
 (T^N-1)^{q'+1}=0\quad\text{on }R^q\Psi.
\]

\textbf{Theorem 3.5.} Let $C$ be a smooth projective curve over the field of fractions of a discrete valuation ring. The inertia group acts on
\[
 H^1(C_{\bar\eta},\Ql)
\]
by quasi-unipotent automorphisms of level $2$: for some $N$ and all $T\in I$,
\[
 (T^N-1)^2=0.
\]
If $A$ is an abelian variety, the same conclusion holds for $H^1(A_{\bar\eta},\Ql)$ and for the Tate module. More generally, for every scheme of finite type over $\eta$, inertia acts quasi-unipotently on $H^i(X_{\bar\eta},\Ql)$; in characteristic zero, for $X$ proper and smooth one obtains $(T-1)^{i+1}=0$ on $H^i$ after passing to an open subgroup of inertia.

\subsection*{4. Vanishing criteria for sheaves of vanishing cycles}
Let $S$ be a strictly local trait and $f:X\to S$ of finite type. Put $n=\dim X_{\bar\eta}$. If $X'$ is the schematic closure of $X_\eta$ in $X$, then on $X_s-X_s'$ one has $\Phi^0=A$ and the other $\Phi^i$ vanish; on $X'$ the $\Phi^i$ and $\Psi^i$ relative to $X/S$ or $X'/S$ agree, and $\Phi^0=0$. Since $X'$ is flat over $S$, this often reduces the problem to the flat case.

\textbf{Theorem 4.2.} One has $\Phi^i=0$ for $i>n$. More precisely, if $x$ is a point of $X_s$ whose closure has dimension $k$ and $\bar x$ is a geometric point localized at $x$, then
\[
 \Phi^i_{\bar x}=0\qquad (i>n-k).
\]

\textbf{Corollary 4.3.} If $f$ is proper and flat and the non-smooth locus of $f$ in $X_s$ has dimension $<d$, the specialization morphism
\[
H^i(X_s)\longrightarrow H^i(X_{\bar\eta})
\]
is an isomorphism for $i>n+d+1$ and an epimorphism for $i=n+d+1$.

To prove in certain cases that the small $\Phi^i$ vanish, one uses local duality, directly or through the local Lefschetz theorems of SGA 2 XIV. This tool is available only in characteristic $0$, or in equal characteristic $p$ when resolution of singularities is known in the dimensions considered.

\textbf{Theorem 4.5.} Suppose that $S$ has characteristic $0$. Let $f:X\to S$ be of finite type and let $x$ be a closed point of $X_s$. Suppose that $x$ is an isolated non-smooth point in its fibre and that $X$ has relative geometric depth $\ge n$ at $x$, i.e. near $x$, $X$ is identified with a subscheme defined by $k$ equations in a smooth scheme of relative dimension $N$ at $x$, with $n<N-k$. Then
\[
 \Phi^i_{\bar x}=0\qquad (i<n).
\]

\textbf{Corollary 4.6.} For $S$ a henselian trait of characteristic $0$, $X/S$ a relative complete intersection of relative dimension $n$, and $x\in X_s$ an isolated non-smooth point in its fibre, one has
\[
 \Phi^i_{\bar x}=0\qquad (i\ne n).
\]

\textbf{Corollary 4.7.} Under the hypotheses of 4.6, $\Phi^n_{\bar x}$ is a flat $A$-module.

\textbf{Variant 4.8.} If in 4.5 one suppresses the isolatedness condition and lets $Z$ be the non-smooth locus of $f$, then
\[
 \Phi^i=0\qquad (i<n-\dim Z).
\]

\textbf{Proposition 4.10.} Under the general hypotheses of 4.5, suppose that $x$ is an isolated non-smooth point in its fibre; that for every point $y$ of $X_{\bar\eta}$ specializing to $x$, with closure of dimension $d(y)$, one has $\prof_{\et,y}(X_{\bar\eta})\ge n-d(y)$; and that $\prof_{\et,x}(X_s)\ge n$. Then
\[
 \Phi^i_{\bar x}=0\qquad (i<n-1).
\]

\subsection*{5. Action of monodromy on fundamental groups}
In this section, which corresponds to Grothendieck's expose of 19 November 1968, we use the semistable reduction theorem for curves. Artin and Winters gave a direct proof; it can also be deduced from the analogous theorem for abelian varieties. We shall also use Schlessinger's theory as presented in Rim's Expose VI.

Let $k$ be algebraically closed, let $C$ be a projective curve over $k$, and let $x_0,\ldots,x_n$ be finitely many closed points of $C$. We say that $X=(C;x_0,\ldots,x_n)$ is stable if $C$ is reduced and has only ordinary double points with distinct tangents, the $x_i$ are distinct smooth points of $C$, and if $C_j$ is an irreducible component of $C$ and $E$ is the set of points of its normalization lying over a singular point of $C$ or one of the $x_i$, then $E$ has at least $3$ elements when $C_j$ has genus $0$ and at least $1$ element when $C_j$ has genus $1$.

The last condition is equivalent to the absence of non-trivial infinitesimal automorphisms, and also to the ampleness of the invertible sheaf
\[
 \omega\Bigl(\sum x_i\Bigr),
\]
where $\omega$ is the dualizing sheaf.

\textbf{Proposition 5.1.1.} Let $X=(C;x_0,\ldots,x_n)$ be stable over the generic point $\eta$ of a trait $S$, with $C_\eta$ smooth. There exists a finite separable extension $\eta'$ of $\eta$ and a stable curve $X'$ over the normalization $S'$ of $S$ in $\eta'$ such that
\[
 X'\otimes_{S'}\eta'=X\otimes_\eta\eta'.
\]

Let $X_\eta$ be a geometrically connected scheme of finite type over the field of fractions of a strictly local trait $S$. For a geometric point $\xi$ of $X_{\bar\eta}$ one has an exact sequence
\[
0\to \pi_1(X_{\bar\eta},\xi)\to \pi_1(X,\xi)\to\Gal(\bar\eta/\eta)\to0,
\tag{5.2.1}
\]
and hence a homomorphism from $\Gal(\bar\eta/\eta)$ to the group of outer automorphisms of $\pi_1(X_{\bar\eta})$. Passing to the maximal prime-to-$p$ quotient gives
\[
\Gal(\bar\eta/\eta)\to \Aut_{\mathrm{ext}}\bigl(\pi_1^{(p')}(X_{\bar\eta})\bigr).
\tag{5.2.2}
\]
If $X_\eta$ has a rational point $x$ and $\bar x$ is the geometric point over it, the sequence splits canonically and gives
\[
[x]:\Gal(\bar\eta/\eta)\to \Aut\bigl(\pi_1^{(p')}(X_{\bar\eta},\bar x)\bigr).
\tag{5.2.3}
\]

\textbf{Theorem 5.3.} The image of wild inertia $P$ under (5.2.2), or equivalently under (5.2.3), is finite.

The proof first reduces to the case where $X_\eta$ is geometrically normal and integral. Let $X'$ be the normalization and let $X''=X'\times_X X'$. After a finite extension of $k(\eta)$, the connected components of $X'$ may be assumed geometrically normal and integral with rational base points, and those of $X''$ geometrically connected with rational base points. Choosing prime-to-$p$ path classes invariant under $P$, the Van Kampen formulas show that if a finite-index subgroup $P'\subset P$ acts trivially on the fundamental groups of the components, then $P'$ acts trivially on $\pi_1^{(p')}(X_{\bar\eta})$.

For $X_\eta$ normal and $U\subset X_\eta$ open, $\pi_1(U)$ maps onto $\pi_1(X_\eta)$. Hence one may suppose $X_\eta$ smooth and affine. By cutting with generic hyperplane sections and using Bertini, one reduces to $X_\eta$ smooth, geometrically connected, and of dimension one. After finite extension of $k(\eta)$ one may suppose that $X_\eta$ is the complement, in a smooth projective curve $C$, of finitely many rational points; applying semistable reduction, one reduces to a stable curve over $S$. The continuation begins with the stable-curve calculation on the next source page.



% SGA7 pages 25--36 additions
\providecommand{\cA}{\mathcal A}
\providecommand{\cV}{\mathcal V}
\providecommand{\cF}{\mathcal F}
\providecommand{\Hom}{\operatorname{Hom}}
\providecommand{\coker}{\operatorname{coker}}
\providecommand{\codim}{\operatorname{codim}}
\providecommand{\Coker}{\operatorname{Coker}}
\providecommand{\Picard}{\operatorname{Pic}}
\providecommand{\End}{\operatorname{End}}
\providecommand{\Prof}{\operatorname{prof}}
\providecommand{\Ima}{\operatorname{Im}}
\providecommand{\Frob}{\operatorname{Frob}}
\providecommand{\cQ}{\mathcal Q}
\providecommand{\Et}{\mathrm{et}}
\providecommand{\profp}{\widehat{\mathbb Z}^{(p')}}


\subsubsection*{Theorem 5.4}
Let $S$ be a strictly local scheme, let $f:\bar X\to S$ be a proper and flat morphism whose special fibre is a curve whose only points of non-smoothness are ordinary double points, and let $Y$ be a subscheme of $\bar X$ which is the union of a finite number of disjoint sections contained in the smooth locus. Put
\[
 X=\bar X-Y.
\]
Let $x$ be a rational point of $X$, let $U$ be the open subset of $S$ over which $X$ is smooth, and let $\xi$ be a geometric point of $U$. Then the image of
\[
 \pi_1(U,\xi)
\]
in
\[
 \Aut\bigl(\pi_1^{(p')}(X_\xi,x_\xi)\bigr)
\]
is abelian and prime to $p$.

One must keep in mind that one cannot make $\pi_1(U,\xi)$ act on $\pi_1^{(p')}(X_\xi,x_\xi)$ except because the groups $\pi_1^{(p')}$ satisfy the specialization theorem.

The proof reduces first to the case where $S$ is complete, and then to the universal case in which $S$ is a formal moduli scheme for the special fibre equipped with its sections. Then $S$ is a formal power series ring over a Cohen ring,
\[
 S=\Spec W[[t_1,\ldots,t_n]],
\]
and $U$ is the complement of a normal-crossings divisor with equation
\[
 t_1\cdots t_m=0.
\]
In this universal case, by Abhyankar's lemma, $\pi_1(U,\xi)$ is abelian and prime to $p$, whence the theorem.

\textbf{Remark 5.5.} One can prove a result analogous to 5.4 in arbitrary relative dimension by considering a generic plane section of dimension one and applying Bertini.

\subsection*{6. Appendix: arithmetic proof of the stable reduction theorem}
\begin{center}
by P. Deligne
\end{center}
Let $S$ be a trait and let $A_\eta$ be an abelian variety over $k(\eta)$, of dimension $d\ne0$. For a local morphism of traits $S'\to S$, denote by $A_{\eta'}$ the abelian variety over $k(\eta')$ deduced from $A_\eta$ by extension of scalars. Let $\cA_{S'}$ be the Neron model of $A_{\eta'}$ over $S'$, let $\cA_{S',s'}$ be its special fibre, and let $\cA^0_{S',s'}$ be the neutral component of $\cA_{S',s'}$. The stable reduction theorem is the following.

\textbf{Theorem 6.1.} For a suitable $S'$, $A_{\eta'}$ has stable reduction, i.e. $\cA^0_{S',s'}$ is an extension of an abelian variety by a torus.

It follows rather easily from 6.1 that one may choose $S'$ so that $k(\eta')$ is a finite separable extension of $k(\eta)$. Suppose first that the residue field of $S$ is of finite type over the prime field; the same argument would work when $k(s)$ is radicial over such a field.

Let $\ell$ be a prime number prime to the residual characteristic exponent $p$, let $T_\ell(A)$ be the Tate module, let
\[
 V_\ell(A)=T_\ell(A)\otimes\Ql,
\]
and let
\[
 \rho:\Gal(\bar\eta/\eta)\longrightarrow \GL(T_\ell(A))
\]
be the natural representation. A preliminary extension of scalars reduces us to the case where the inertia group $I$ acts unipotently, by 1.2, hence through its largest pro-$\ell$ quotient $\Zl(1)$, cf. 0.3. Let $T$ be the image in $\GL(T_\ell(A))$ of a generator of $\Zl(1)$, and let $r$ be the largest integer $\ge0$ such that $(T-1)^r\ne0$.

\textbf{Lemma 6.2.} The map
\[
 \sigma\in I,
 \quad v\in V_\ell(A)
 \longmapsto (\rho(\sigma)-1)^r(v)
\]
induces morphisms and an isomorphism
\begin{equation}\tag{6.2.1}
\begin{tikzcd}[column sep=large,row sep=large]
 V_\ell(A)_I\otimes\Ql(r) \arrow[r,"M"] \arrow[d]
 & V_\ell(A)^I \\
 V_\ell(A)/\Ker(T-1)^r\otimes\Ql(r) \arrow[r,"\sim"']
 & \Ima(T-1)^r \arrow[u]
\end{tikzcd}
\end{equation}
The modules $V_\ell(A)_I$ and $V_\ell(A)^I$ are $\Gal(\bar s/s)$-modules, and $M$ is Galois-equivariant.

\textbf{6.3.} An $\ell$-adic representation
\[
 \tau:\Gal(\bar s/s)\longrightarrow \GL(V)
\]
is said to be of weight $n$ ($n\in\mathbb Z$) if the following condition is satisfied: for a suitable model $X$ of $k(s)$, with $X$ an irreducible variety of finite type over the prime field and $k(s)=k(X)$, the representation $\tau$ factors through $\pi_1(X,\bar s)$, and, for $x$ a closed point of $X$, the eigenvalues of $\tau(F_x)$ are algebraic numbers all of whose complex conjugates have absolute value
\[
 |k(x)|^{n/2},
\]
where $F_x$ denotes geometric Frobenius $\Phi_x^{-1}$.

If $V$ (respectively $W$) is of weight $n$ (respectively $m$) and $n\ne m$, then
\[
 \Hom_{\Gal}(V,W)=0.
\]

\textbf{6.4.} It follows from the universal property of the Neron model that
\[
 V_\ell(A)^I=V_\ell(\cA^0_{S,s}).
\]
If $a$ is the dimension of the largest abelian quotient of $\cA^0_{S,s}$, and $m$ the dimension of the multiplicative part of $\cA^0_{S,s}$, then, by Weil, $V_\ell(A)^I$ is an extension of a representation of dimension $2a$ and weight $-1$ by a representation of dimension $m$ and weight $-2$.

\textbf{6.5.} Let $A_\eta^*$ be the abelian variety dual to $A_\eta$. Recall that $V_\ell(A^*)$ and $V_\ell(A)$, hence also $V_\ell(A^*)^I$ and $V_\ell(A)_I$, are in duality with values in $\Ql(1)$. If $a^*$ and $m^*$ are defined for $A^*$ as $a$ and $m$ for $A$, it follows from 6.4 that $V_\ell(A)_I$ is an extension of a representation of dimension $m^*$ and weight $0$ by a representation of dimension $2a^*$ and weight $-1$.

\textbf{6.6.} The Galois module $\Ql(r)$ has weight $-2r$. Since $\Ima(T-1)^r\ne0$, one deduces from 6.4, 6.5, and the isomorphism 6.2 that $r\le1$. If $r=0$, then
\[
\begin{cases}
 a+m\le d,\\
 2a+m=2d,
\end{cases}
\]
hence $m=0$, $a=d$, and $A_\eta$ has good reduction. Suppose that $r=1$. Then 6.2 implies that $\Ima(T-1)$ has weight $-2$ and that $V_\ell(A)/V_\ell(A)^I$ has weight $0$; these spaces have the same dimension $m_1\le m$. Hence
\[
 \dim V_\ell(A)^I=2d-m_1=2a+m,
\]
and
\[
 2\dim \cA^0_{S,s}\ge 2(a+m)
 \ge 2a+m+m_1=2d=2\dim A_{S,s}.
\]
Thus $A_\eta$ has stable reduction, since $\dim \cA^0_{S,s}=a+m$. At the same time one obtains $m_1=m$, i.e.
\[
 \Ima(T-1)\subset V_\ell(A)^I\simeq V_\ell(\cA^0_{S,s})
\]
is the $V_\ell$ of the multiplicative part of $\cA^0_{S,s}$.

\textbf{6.7.} To treat the general case, we shall use 0.5, cf. 1.3. We reduce to the case where
\begin{enumerate}[label=\alph*)]
\item the conditions of 0.5.2 or 0.5.3 are satisfied;
\item $\Gal(\bar\eta/\eta)$ acts trivially on $A_\ell$;
\item the action of $I$ on $T_\ell(A)$ is unipotent.
\end{enumerate}
Under these hypotheses we shall prove that $A_\eta$ has stable reduction. Hypothesis a) permits one to apply 0.5.1 and to write
\[
 S=\limproj S_i,
\]
as in 1.3, whose notation we resume. Let $s_i$ be the closed point of $S_i$. For $i$ sufficiently large, the neutral component $\cA^0_S$ of the Neron model $\cA_S$ comes from a smooth group scheme with connected fibres $\cA_i$ over $S_i$, and $(\cA_i)_\ell$ is constant on $S_i-D_i$. The group $\pi_1(S_i-D_i,\bar\eta)$ acts on $T_\ell(A)$, and acts trivially on
\[
 T_\ell(A)/\ell T_\ell(A)\simeq A_\ell.
\]
As in 1.3, one sees that $P_i$ acts trivially. Then, by (1.3.1),
\[
 T_\ell(A)_I=T_\ell(A)_{I_i},
 \qquad
 T_\ell(\cA_{i,s_i})=T_\ell(\cA^0_S)=T_\ell(A)^I=T_\ell(A^{I_i}),
\]
and (6.2.1) is a diagram of $\Gal(\bar s_i/s_i)$-modules to which one can apply the arguments of 6.4--6.6.

\subsubsection*{Bibliography}
\begin{enumerate}[label={[\arabic*]}]
\item M. Artin, \emph{Algebraic approximation of structures over complete local rings}, Publ. Math. I.H.E.S. 36 (1969), 23--58.
\item M. Artin and G. Winters, \emph{Degenerate fibres and reduction of curves}.
\item P. Deligne and D. Mumford, \emph{The irreducibility of the space of curves of a given genus}, Publ. Math. I.H.E.S. 36 (1969), 75--110.
\item J.-P. Serre, \emph{Corps locaux}, Publ. Math. Univ. Nancago--Hermann, 1962.
\item J.-P. Serre and J. Tate, \emph{Good reduction of abelian varieties}, Ann. of Math. 88 (1968), 492--517.
\end{enumerate}

\section*{Expose II. Finiteness properties of the fundamental group}
\begin{center}
by Michele Raynaud
\end{center}
Let $k$ be a separably closed field of characteristic $p\ge0$, let $X$ be a connected scheme of finite type over $k$, and let
\[
 \pi_1^{(p')}(X)
\]
be the largest quotient ``prime to $p$'' of $\pi_1(X)$. We show in this expose that, if $X$ is a surface, $\pi_1^{(p')}(X)$ is topologically of finite presentation, and that the same is true for arbitrary $X$ if one admits resolution of singularities.

In the case of a surface, the method of proof, due to J.-P. Murre, consists in reducing to the case where one has a fibration over $\mathbb P^1$ having the properties (i), (ii), and (iii) of the following proposition.

\textbf{Proposition 2.1.} Let $k$ be an algebraically closed field and let $S$ be a projective, irreducible, smooth surface over $k$. Then one can find a blow-up
\[
 g:S'\longrightarrow S,
\]
where $S'$ is a regular surface equipped with a fibration
\[
 f:S'\longrightarrow \mathbb P^1,
\]
such that $f$ has a section $\sigma$ and the following conditions are satisfied:
\begin{enumerate}[label=\textup{(\roman*)}]
\item $f$ is smooth at the points of $\sigma(\mathbb P^1)$;
\item the fibres of $f$ are geometrically integral curves having only ordinary singularities;
\item there exists a non-empty open $U$ of $\mathbb P^1$ such that the fibres of $f$ at the points of $U$ are smooth.
\end{enumerate}

The proposition follows from the existence of a projective embedding $S\to \mathbb P^r$ such that one can find a Lefschetz pencil for that embedding. Let
\[
 D=\{H_t\}_{t\in\mathbb P^1}
\]
be such a pencil of hyperplanes. Recall that it has the following properties: its axis $\Delta$ cuts $S$ transversally, so that the closed subscheme $\Delta\cap S$ is reduced and consists of a finite set of closed points $x_1,\ldots,x_n$ of $S$; there exists a non-empty open $U\subset\mathbb P^1$ such that, for every $t\in U$, $H_t$ cuts $S$ transversally; and, for every $t\in\mathbb P^1-U$, $H_t$ cuts $S$ transversally except at one point, which is an ordinary quadratic singular point.

Through every point $x$ of $S$ not belonging to $\Delta$ there passes a unique hyperplane $H_t$; the application that, to $x$, associates $t\in\mathbb P^1$ is a rational map $a:S\dashrightarrow\mathbb P^1$ whose domain of definition is $S-(\Delta\cap S)$. Consider the blow-up
\[
 g:S'\longrightarrow S
\]
of the closed subscheme $S\cap\Delta$ of $S$. The scheme $S'$ is regular and the fibre $S'_{x_i}$ over a point $x_i$ is identified with the projective fibre defined by the tangent space to $S$ at $x_i$, hence is isomorphic to $\mathbb P^1$. Moreover the rational map $f=ag$ is everywhere defined. Finally, to each point $x_i$ one associates the section of $f$ which, to a point $t\in\mathbb P^1$, makes correspond the tangent vector at $x_i$ to $H_t\cap S$.

Let us verify (i), (ii), and (iii). For each $t\in\mathbb P^1$, the curves $S'_t$ and $S_t=H_t\cap S$ are isomorphic: the morphism $S'_t\to S_t$ is an isomorphism except perhaps at the points $x_i$, but $S_t$ is regular at those points; moreover $S'_t$ is a complete intersection in $S'$, hence integral, and therefore isomorphic to $S_t$. Conditions (i) and (iii) are just conditions a) and b) of the Lefschetz pencil. Finally $f$ is smooth at the points of $\sigma$, since $f$ is flat and the fibres $S'_t$ are geometrically regular at the points $\sigma(t)$.

\textbf{Corollary 2.2.} With the notation of 2.1, let $Y$ be a normal-crossings divisor in $S$, and let $Y'=Y\times_S S'$. Then one can choose $f,g,U$ so that $Y'\times_{\mathbb P^1}U$ is a normal-crossings divisor relative to $U$.

Indeed, if there exists a Lefschetz pencil for an embedding $S\to\mathbb P^r$, the Lefschetz pencils form an open subset of the variety of lines in $\mathbb P^r$. The same holds for the set of Lefschetz pencils whose axis does not meet $Y$ and for which there exists a non-empty open $U_1\subset\mathbb P^1$ such that $Y\cap H_t$ is smooth for $t\in U_1$. It then suffices to construct $f$ and $g$ from a pencil satisfying these two conditions, and to replace $U$ by $U\cap U_1$.

\textbf{2.3.0.} Let $G$ be a profinite group. Recall that $G$ is topologically of finite generation if there exists an integer $n$ such that $G$ is isomorphic to a quotient of the free pro-group on $n$ generators, $F(n)$. Let $K=\Ker(F(n)\to G)$. One says that $G$ is topologically of finite presentation if, in addition, one can find finitely many elements $k_1,\ldots,k_r$ of $K$ such that the smallest closed invariant subgroup of $K$ containing $k_1,\ldots,k_r$ is equal to $K$.

\textbf{Theorem 2.3.1.} Let $k$ be a separably closed field of characteristic $p\ge0$, and let $X$ be a connected scheme of finite type over $k$. Assume that schemes of finite type and of dimension $\le \dim(X)$ over an algebraic closure of $k$ are strongly desingularizable (SGA 5 I 3.1.5); this condition is satisfied when $\dim X=2$, by [1], or when $k=0$, by [2]. Then the fundamental group
\[
 \pi_1^{(p')}(X)
\]
is topologically of finite presentation.

By SGA 1 IX 4.10, one may suppose $k$ algebraically closed.

\emph{1) Reduction to the case where $X$ is a regular surface.} If $f:X'\to X$ is a morphism of finite type, set $X''=X'\times_X X'$, and denote by $(X'_a)_{a\in A}$ and $(X''_b)_{b\in B}$ the sets of connected components of $X'$ and $X''$, respectively. When $f$ is a morphism of effective descent for etale coverings, the fundamental group $\pi_1(X)$ can be computed explicitly in terms of the groups $\pi_1(X'_a)$ and $\pi_1(X''_b)$. It follows from SGA 1 IX 5.2, 5.3 that, if the $\pi_1(X'_a)$ are topologically of finite generation, so is $\pi_1(X)$; and, if the $\pi_1(X'_a)$ are topologically of finite presentation and the $\pi_1(X''_b)$ topologically of finite generation, then $\pi_1(X)$ is topologically of finite presentation.

Since $X$ is desingularizable, one can find a regular scheme $X'$ and a proper birational morphism $f:X'\to X$; since $f$ is a morphism of effective descent for etale coverings, it is enough to prove the theorem under the additional hypothesis that $X$ is regular. Indeed one first deduces that $\pi_1(X)$ is topologically of finite generation; knowing that $\pi_1(X')$ is topologically of finite presentation and the $\pi_1(X''_b)$ topologically of finite generation, one deduces that $\pi_1(X)$ is topologically of finite presentation.

Let now $(U_i)_{i\in I}$ be a finite covering of $X$ by affine open subsets, and let
\[
 f:\coprod_i U_i\longrightarrow X
\]
be the canonical morphism. Since $f$ is a morphism of effective descent for etale coverings, one sees that it is enough to prove the theorem for the $U_i$; one is therefore reduced to the case where $X$ is affine and smooth.

Assume from now on that $X$ is affine, smooth over $k$, and $\dim X>2$. One may then apply the Lefschetz theorem for quasi-projective schemes: if $Y$ is a sufficiently general hyperplane section of $X$, one has an isomorphism
\[
 \pi_1(Y)\simeq \pi_1(X).
\]
It is therefore enough to prove the theorem for $Y$, which reduces us to the case $\dim X=2$.

\emph{2) Case where $X$ is a smooth surface over $k$.} By the theorem of strong desingularization of surfaces, one can find a smooth, integral, projective surface $S$ and an open immersion $i:X\to S$ such that the divisor $Y=S-X$ has normal crossings. Apply Corollary 2.2 and keep its notation.

It is enough to prove that the fundamental group of $X'=X\times_S S'$ is topologically of finite presentation. Since $X$ and $X'$ are regular and since
\[
 \codim\left(X'-\bigcup_{1\le i\le n}\{x_i\},X\right)\ge2,
\]
every etale covering of $X'$ induces on $X-\bigcup_{1\le i\le n}\{x_i\}$ an etale covering which extends uniquely to an etale covering of $X$ (SGA 2 1.11). This shows that one has an isomorphism
\[
 \pi_1(X')\simeq \pi_1(X).
\]

Consider the Cartesian diagram
\[
\begin{tikzcd}
X' & X'_U \arrow[l]\\
\mathbb P^1 \arrow[u,"h"] & U \arrow[l] \arrow[u,"h_U"']
\end{tikzcd}
\]
where $h=f|_{X'}$. Let $L$ be the set of prime numbers distinct from $p$, and let $\bar\eta$ be a geometric point above the generic point $\eta$ of $\mathbb P^1$. Let us verify that the hypotheses of SGA 1 XIII 4.8 are satisfied. The morphism $h$ has a section $\sigma$, because $f$ does. Condition a) of SGA 1 XIII 4.7 is proved as follows: the fibres of $h$ being geometrically integral, the morphism $h$ is $0$-acyclic and locally $0$-acyclic; it remains to show that, for every locally constant sheaf of sets $F$ on $X'$, the pair $(F,h)$ is cohomologically proper in dimension $\le0$. If $\xi$ is a closed point of $\mathbb P^1$ and $Z$ denotes the integral closure of $\Spec \cO_{\mathbb P^1,\xi}$ in $\bar\eta$, the scheme $X'_Z=X'\times_{\mathbb P^1}Z$ satisfies condition $(S_2)$ at the closed points of $X'$ and is regular at every other point, hence is normal. Therefore an etale covering of $X'_Z$ whose restriction to $X'_{\bar\eta}$ is connected is necessarily connected. This proves condition a). Since $X'_U$ is the complement in $S'_U$ of a normal-crossings divisor relative to $U$, $h_U$ is locally $1$-aspherical for $L$ (SGA 4 XV 2.1), and, for every finite constant sheaf $F$ of $L$-groups on $X'_U$, the pair $(F,h_U)$ is cohomologically proper in dimension $\le1$ (SGA 1 XIII 2.4); this is condition b). Since $\mathbb P^1$ is normal and $T=\mathbb P^1-U$ has codimension $1$, one has
\[
 \prof^{\et}_T(S)\ge2,
\]
which gives condition c). The condition
\[
 \prof^{\mathrm{top}}_x(X')\ge3
\]
holds at every closed point of $X'$ by purity, and gives condition e). Finally, for every point $t\in\mathbb P^1$, $h$ is smooth at $\sigma(t)$, hence is locally $1$-aspherical for $L$ at $\sigma(t)$, giving condition d).

By SGA 1 XIII 4.7, if $a$ is a geometric point of $X'_{\bar\eta}$, one has an exact sequence
\[
1\longrightarrow \pi_1^{(p')}(X'_{\bar\eta},a)
\longrightarrow \pi_1'(X'_U,a)
\longrightarrow \pi_1(U,a)
\longrightarrow 1,
\]
and the choice of the section $\sigma$ defines a splitting of this exact sequence. Let $K$ denote the image of $\pi_1(U,a)$ in
\[
 \Aut\bigl(\pi_1^{(p')}(X'_{\bar\eta},a)\bigr).
\]
The group of coinvariants under $K$,
\[
 \pi_1^{(p')}(X'_{\bar\eta},a)_K,
\]
is the quotient of $\pi_1^{(p')}(X'_{\bar\eta},a)$ by the closed invariant subgroup generated by the elements of the form $x^{-1}q(x)$, where $x\in\pi_1^{(p')}(X'_{\bar\eta},a)$ and $q\in K$. By loc. cit. 4.7.4, one has an isomorphism
\[
 \pi_1^{(p')}(X',a)\simeq \pi_1^{(p')}(X'_{\bar\eta},a)_K.
\]

Let
\[
 \mathbb P^1-U=\coprod_{1\le i\le n}\{t_i\},
\]
and, for each $i$, let $\bar T_i$ be the strict localization of $\mathbb P^1$ at a geometric point $\bar t_i$ over $t_i$, and let $s_i$ be the generic point of $\bar T_i$. An etale covering of $U$ whose restriction to each $s_i$ is trivial extends to an etale covering of $\mathbb P^1$, hence is trivial. Thus the assertion above amounts to saying that $\pi_1(U,a)$ is the closed invariant subgroup generated by the images of the Galois groups $\Gal(\bar k(s_i),k(s_i))$ in $\pi_1(U,a)$. By I 5.3, the image $K_i$ of $\Gal(\bar k(s_i),k(s_i))$ in $\Aut(\pi_1^{(p')}(X'_{\bar\eta},a))$ is finite. Since $K$ is the smallest closed invariant subgroup containing all the $K_i$, and since $\pi_1^{(p')}(X'_{\bar\eta},a)$ is topologically of finite presentation (SGA 1 XIII 3.2), this proves that $\pi_1^{(p')}(X'_{\bar\eta},a)_K$, and therefore also $\pi_1^{(p')}(X)$, to which it is isomorphic, is topologically of finite presentation.

\textbf{Remark 2.3.2.} Resolution of singularities is not necessary for proving Theorem 2.3.1 in the case where $X$ is the open complement of a normal-crossings divisor in a projective scheme smooth over $k$.

\subsubsection*{Bibliography to Expose II}
\begin{enumerate}[label={[\arabic*]}]
\item S. Abhyankar, \emph{Local uniformization of algebraic surfaces over ground field of characteristic $p=0$}, Amer. J. Math. 63 (1956).
\item H. Hironaka, \emph{Resolution of singularities of an algebraic variety over a field of characteristic zero}, Ann. of Math. 79 (1964).
\item M. Raynaud, \emph{Theoreme de Lefschetz en cohomologie des faisceaux coherents et en cohomologie etale}, thesis, to appear.
\end{enumerate}



\begin{center}
{\Large SGA 7-I, Expos\'e VI}\\[0.4em]
{\large Formal Deformation Theory}\\[0.4em]
by D. S. Rim
\end{center}

\section*{1. Formal existence theorem}

We recall the notion of cofibered category (SGA 1, IV).

Let $C$ be a fixed category. A cofibered category $A$ over $C$ is, by definition, a category $A$ together with a functor
\[
 P:A\longrightarrow C
\]
such that:

\textbf{Ax. 1.} For any map $f$ in $C$ and any object $a$ in $A$ with $P(a)$ equal to the source of $f$, there exists an $f$-morphism
\[
 \alpha:a\longrightarrow b
\]
which is cocartesian, i.e. for any $f$-morphism $\alpha':a\to b'$ there exists a unique $T$-morphism $\tau:b\to b'$ in $A$ such that
\[
 \alpha'=\tau\alpha,
\]
where $T$ is the target of $f$.

\textbf{Ax. 2.} A composite of cocartesian morphisms in $A$ is again cocartesian.

The following properties of a cofibered category $P:A\to C$ follow immediately from the definitions.

\textbf{Cancellation.} Let $\alpha,\beta:a\to a'$ be two cocartesian maps such that $P(\alpha)=P(\beta)$. If there exists a cocartesian map $\gamma:b\to a$ such that $\alpha\gamma=\beta\gamma$, then $\alpha=\beta$.

\textbf{Factorization.} Given cocartesian maps $\alpha:a\to a'$ and $\beta:a\to a''$, there exists a cocartesian map $\gamma:a'\to a''$ such that $\gamma\alpha=\beta$ if and only if there exists $f:P(a')\to P(a'')$ in $C$ such that $fP(\alpha)=P(\beta)$. Furthermore $\gamma$ is unique.

Let $A$ be a cofibered category over $C$. For each object $S$ in $C$ we denote by $A(S)$ the subcategory of $A$ whose objects are the objects $a$ in $A$ such that $P(a)=S$, and
\[
 \Hom_{A(S)}(a,a')=\{\alpha\in\Hom_A(a,a')\mid P(\alpha)=\id_S\}.
\]
A cofibered groupoid over $C$ is a cofibered category $A$ over $C$ such that $A(S)$ is a groupoid for every object $S$ in $C$. We recall that a groupoid is a category in which every morphism is an isomorphism. Thus, a cofibered groupoid over $C$ is a cofibered category $P:A\to C$ such that every map in $A$ is cocartesian. Furthermore, for every map $\alpha\in\operatorname{Fl}(A)$, $\alpha$ is an isomorphism in $A$ if and only if $P(\alpha)$ is an isomorphism.

Throughout the rest of this section, we fix a complete noetherian local ring $\Lambda$ with residue field $k$, and we set $\CL$ to be the category of artinian local $\Lambda$-algebras with the same residue field $k$, where the maps are local homomorphisms over $\Lambda$. From 1.2 on, a cofibered groupoid $A$ over $\CL$ is always assumed to have the property that $A(k)=\{e\}$, where $\{e\}$ denotes the category in which $\Hom_{\{e\}}(a,b)$ consists of one and only one element for any two objects $a,b$ in $\{e\}$. A map $a'\xrightarrow{\alpha}a$ in $A$ will be called simply a $\CL$-map if $P(a')=P(a)$ and $P(\alpha)=1_{P(a)}$. A map $a'\xrightarrow{\alpha}a$ in $A$ will be called surjective, by abuse of language, if $P(\alpha):P(a')\to P(a)$ is surjective in $\CL$.

\textbf{Example 1.1(a).} Let $F:C\to(\mathrm{categories})$ be a functor. Define the category $\int F$ as follows: an object in $\int F$ is a pair $(S,a)$ where $S\in\ob(C)$ and $a\in\ob F(S)$, and a map
\[
 (S,a)\longrightarrow (S',a')
\]
is a pair $(f,u')$ where $f:S\to S'$ is a map in $C$ and $u':F(f)a\to a'$ is a map in $F(S')$. Composition of two maps
\[
 (S,a)\xrightarrow{(f,u')}(S',a')\xrightarrow{(g,u'')}(S'',a'')
\]
is given by
\[
 (gf,u''\cdot F(g)u'):(S,a)\longrightarrow (S'',a'').
\]
With respect to the functor $P:\int F\to C$ given by $(S,a)\mapsto S$, $\int F$ becomes a cofibered category over $C$. If $F(S)$ is a groupoid for every $S\in\ob(C)$, then $\int F$ is a cofibered groupoid over $C$. In particular, a functor $F:C\to(\mathrm{groups})$ yields a cofibered group, and a functor $F:C\to(\Ens)$ yields a cofibered discrete groupoid. We also note that a natural transformation $\varphi:F\to G$, where $F,G:C\to(\mathrm{categories})$ are functors, induces a cocartesian functor $\varphi:\int F\to\int G$ over $C$.

\textbf{Example 1.1(b).} Let $X,Y$ be schemes over $\Lambda$ together with a fixed isomorphism $X\otimes_\Lambda k\xrightarrow{\sim}Y\otimes_\Lambda k$. Then the functor
\[
 \Isom_{X,Y}:\CL\longrightarrow(\mathrm{groupoids})
\]
given by
\[
 \Isom_{X,Y}(S)=\{\text{isomorphisms }X\otimes_\Lambda S\xrightarrow{\sim}Y\otimes_\Lambda S\text{ inducing the fixed one over }k\}
\]
yields a cofibered groupoid $P:\int\Isom_{X,Y}\to\CL$.

\textbf{Example 1.1(c).} Let $X_0$ be a fixed scheme over $k$, and let $M_{X_0}$ be the category consisting of deformations of $X_0$ over $\CL$ (see \S4). Thus an object in $M_{X_0}$ is a pair $(R,X)$ where $R$ is in $\CL$ and $X$ is a deformation of $X_0$ to $R$. Then the functor
\[
 P:M_{X_0}\longrightarrow\CL,
 \qquad (R,X)\longmapsto R
\]
defines a cofibered groupoid over $\CL$.

\textbf{Example 1.1(d).} For any groupoid $X$ we denote by $[X]$ the discrete groupoid consisting of the isomorphism classes of objects in $X$. If $P:F\to C$ is a cofibered groupoid, we obtain a functor $[F]:C\to(\Ens)$ given by $[F](S)=[F(S)]$, and in turn a cofibered discrete groupoid, which is, by abuse of notation, denoted by $[F]$.

\textbf{Definition 1.2.} A cofibered groupoid $A$ over $\CL$ is called semi-homogeneous if the following properties hold.

\begin{enumerate}[label=(\arabic*)]
\item Every diagram of solid arrows in $A$ where $a'\to a$ is surjective can be completed into a commutative diagram including the dotted arrow:
\[
\begin{tikzcd}
 a'' \arrow[r,dashed] \arrow[d] & a' \arrow[d]\\
 a_1 \arrow[r] & a .
\end{tikzcd}
\]

\item Let
\[
 R\times_k k[\varepsilon]\mathrel{\substack{\longrightarrow\\[-.6ex]\longrightarrow}} R
\]
be the projection maps, where $k[\varepsilon]$ is the ring of dual numbers over $k$, and let
\[
 a'\mathrel{\substack{\longrightarrow\\[-.6ex]\longrightarrow}} a
\]
be a pair of $\pi_1$-, $\pi_2$-maps in $A$. Then there exists a $\CL$-map $\gamma:a'\to a''$ with $\alpha=\beta\gamma$ if and only if $(\pi_1)_*(a')=(\pi_2)_*(a'')$.
\end{enumerate}

\textbf{Remark 1.3(a).} A reformulation of the notion in terms of 2-categories will be given later in 2.6(b). We note that, if $A$ is semi-homogeneous, then so is $[A]$; in particular 1.2(2) entails that the canonical map
\[
 [A(R\times_k k[\varepsilon])]\longrightarrow [A(R)]\times [A(k[\varepsilon])]
\]
is bijective.

\textbf{Remark 1.3(b).} Consider a commutative diagram in a cofibered category $A$ over $\CL$:
\[
\begin{tikzcd}
& a' \arrow[dl] \arrow[dr] & \\
a_1 \arrow[dr] && a_2 \arrow[dl]\\
& a &
\end{tikzcd}
\]
Set $R=P(a)$, $R'=P(a')$, and $R_i=P(a_i)$ for $i=1,2$. If we set $a''$ to be the direct image under the map
\[
 (P(\nu_1),P(\nu_2)):R'\longrightarrow R_1\times_R R_2,
\]
we obtain a commutative diagram. Thus 1.2(1) can be replaced by:

\textbf{(1)'} Every diagram in $A$ above $R_1\leftarrow R'\to R_2$ in $\CL$, where
\[
 R'\longrightarrow R_1\times_R R_2
\]
is surjective, can be completed to a commutative diagram, where the maps $\mu_i$ ($i=1,2$) are projections.

\textbf{Remark 1.3(c).} If $A$ is a semi-homogeneous cofibered groupoid over $\CL$, then $[A(k[\varepsilon])]$ is canonically provided with a vector-space structure over $k$, where $k[\varepsilon]$ is the ring of dual numbers over $k$. Indeed, let $V$ be the category of finite-dimensional vector spaces over $k$. Then any functor
\[
 H:V\longrightarrow(\Ens)
\]
preserving the product is canonically factored through the forgetful functor $V\to(\Ens)$. On the other hand, the functor
\[
 G:V\longrightarrow\CL,
 \qquad W\longmapsto k[W]=k\oplus W,
 \quad w^2=0,
\]
transforms products in $V$ into fibre products over $k$ in $\CL$. If $A$ is a semi-homogeneous cofibered groupoid over $\CL$, then
\[
 [A(k[V]\times_k k[W])]\longrightarrow [A(k[V])]\times[A(k[W])]
\]
is bijective by 1.3(a), and the functor $V\mapsto [A(k[V])]$ commutes with products, hence the result.

\textbf{Definition 1.4.} (1) Let $R\in\ob(\CL)$. A small extension of $R$ in $\CL$ is a surjective map $R'\to R$ in $\CL$ such that $(\Ker f)^2=0$ and $\Ker f\simeq k$. In other words, $R'\xrightarrow{f}R$ in $\CL$ is a small extension of $R$ if
\[
 0\longrightarrow k\xrightarrow{t}R'\xrightarrow{f}R\longrightarrow0
\]
is exact for some $t\in R'$ with $t^2=0$.

(2) Let $A$ be a cofibered groupoid over $\CL$, and let $a$ be an object in $A$. A small prolongation of $a$ is an arrow $a'\to a$ in $A$ such that $P(a')\to P(a)$ is a small extension in $\CL$. An arrow $a'\to a$ in $A$ will be called versal for small prolongations of $a$ if, for every small prolongation $a_1\to a$ in $A$, there exists a commutative diagram
\[
\begin{tikzcd}
 a' \arrow[rr] \arrow[dr] && a_1 \arrow[dl]\\
& a &
\end{tikzcd}
\]

\textbf{Proposition 1.5.} Let $A$ be a semi-homogeneous cofibered groupoid over $\CL$.

\begin{enumerate}[label=(\arabic*)]
\item Let $R\times_k k[\varepsilon]\mathrel{\substack{\longrightarrow\\[-.6ex]\longrightarrow}} R$ be the projection maps in $\CL$, and let $a'\xrightarrow{\alpha}a$ be a map. Then there exists a map $a\xrightarrow{\beta}a'$ such that $\alpha\beta=1_a$ if and only if there exists a map $a\to(\pi_2)_*a'$.

\item Let $R'\xrightarrow{f}R$ be a small extension in $\CL$, and consider a commutative diagram
\[
\begin{tikzcd}
 R'\times_R R' \arrow[r] \arrow[d] & R' \arrow[d,"f"]\\
 R' \arrow[r,"f"'] & R
\end{tikzcd}
\qquad\text{above}\qquad
\begin{tikzcd}
 a'' \arrow[r] \arrow[d] & a' \arrow[d]\\
 a' \arrow[r] & a .
\end{tikzcd}
\]
Define the canonical map
\[
 \mu:R'\times_R R'\longrightarrow k[\Ker f],\qquad
 \mu(r_1,r_2)=r_1(0)+(r_2-r_1).
\]
If there exists a map $a'\to(\mu)_*(a'')$, then there exists a map $\gamma:a'\to a$ such that $\alpha=a\gamma$.
\end{enumerate}

\textbf{Proof.} (1) The ``only if'' part is clear. Now assume that there exists a map $\varphi:a\to(\pi_2)_*a'$. If we set
\[
 h:R\longrightarrow R\times_k k[\varepsilon],\qquad h=1\times P(\varphi),
\]
then the composite map $R\xrightarrow{h}R\times_k k[\varepsilon]\xrightarrow{\pi_1}R$ is the identity, and therefore we get maps $a\xrightarrow{\gamma}a'$ such that $\pi_1\gamma=1_a$, where $P(\gamma)=h$ and $P(a')=\pi_1$. On the other hand, $\pi_2h=P(\varphi)$ and hence there exists a $\pi_2$-map $h^*a\to(\pi_2)_*a'$. This means that $(\pi_1)^*h^*a=(\pi_2)^*\varphi$, and therefore by 1.2(2) we have a $\CL$-map $u:h^*a\to a'$ such that $\alpha u=a_1$. Define $\beta:a\to a'$ by $\beta=u\gamma$. Then $\alpha\beta=\alpha u\gamma=a_1\gamma=1_a$.

(2) The map $1\times f:R'\times_RR'\to R'\times_k k[\Ker f]$ is an isomorphism compatible with the projections on the first factor, and $\pi_2(1\times f)=\mu$, where $\pi_2:R'\times_k k[\Ker f]\to k[\Ker f]$ is the projection map. We note that $k[\Ker f]\simeq k[\varepsilon]$ since $R'\to R$ is a small extension. Therefore, if there exists a map $a'\to(\mu)_*a''=(\pi_2)_*(1\times f)_*a''$, then by (1) above there exists a map $\gamma':a'\to a''$ such that $\alpha'\gamma'=1_{a'}$. Define $\gamma:a'\to a$ by $\gamma=\beta\gamma'$. We then have $\alpha\gamma=a$.

\textbf{Corollary 1.6.} Let $A$ be a semi-homogeneous cofibered groupoid over $\CL$ and consider a diagram
\[
\begin{tikzcd}
 a' \arrow[dr,"\alpha"'] && a_1 \arrow[dl,"\alpha_1"]\\
& a &
\end{tikzcd}
\]
in $A$, where $a_1$ is a small prolongation. Assume that $a$ is versal for small prolongations of $e$. Then there exists a map $a'\xrightarrow{\rho}a_1$ such that $\alpha=\alpha_1\rho$ in $A$ if and only if there exists a map $f:P(a')\to P(a_1)$ such that $P(\alpha)=P(\alpha_1)\cdot f$ in $\CL$.

\textbf{Proof.} Assume that $P(\alpha)=P(\alpha_1)\cdot f$, where $f:P(a')\to P(a_1)$. Replacing $a'$ by its direct image under $f$, we may and shall assume that $P(\alpha)=P(\alpha_1)$. Set $R'=P(a_1)$ and $R=P(a)$. Since $A$ is semi-homogeneous, we get a commutative diagram
\[
\begin{tikzcd}
& a'' \arrow[dl,"\beta"'] \arrow[dr,"\beta_1"] &\\
a' \arrow[dr,"\alpha"'] && a_1 \arrow[dl,"\alpha_1"]\\
& a &
\end{tikzcd}
\qquad
\begin{tikzcd}
& R'\times_R R' \arrow[dl] \arrow[dr] &\\
R' \arrow[dr] && R' \arrow[dl]\\
& R &
\end{tikzcd}
\]
above the corresponding diagram of rings. Since $a$ is versal for small prolongations of $e$, so is $a'$, and therefore our statement follows from 1.5(2).

\textbf{Definition 1.7.} For each object $R$ in $\CL$ we put
\[
 \bar\Omega_R=M/(\ma R+M^2),
\]
where $\ma$ and $M$ are the maximal ideals of $\Lambda$ and $R$ respectively. It is a finite-dimensional vector space over $k$, and $R\mapsto\bar\Omega_R$ defines a functor
\[
 \bar\Omega:\CL\longrightarrow V,
\]
where $V$ is the category of finite-dimensional vector spaces over $k$. If $A$ is a cofibered category over $\CL$, then the composite functor
\[
 A\xrightarrow{P}\CL\xrightarrow{\bar\Omega}V
\]
is, by abuse of notation, denoted by $\bar\Omega$ again.

\textbf{Proposition 1.8.} Let $A$ be a semi-homogeneous cofibered groupoid over $\CL$. Then, given an object $a$ in $A$ versal for small prolongations of $e$, there exists $a'\to a$ such that:
\begin{enumerate}[label=(\roman*)]
\item $a'$ is versal for small prolongations of $a$;
\item $P(a')\to P(a)$ is surjective, and the kernel is annihilated by the maximal ideal of $P(a)$;
\item $\bar\Omega_{a'}\to\bar\Omega_a$ is an isomorphism.
\end{enumerate}

\textbf{Proof.} Set $R=P(a)$ and consider an exact sequence
\[
 0\longrightarrow J\longrightarrow S\longrightarrow R\longrightarrow0
\]
where $S$ is a formal power-series ring over $\Lambda$ in a finite number of variables, and $J\subset M^2$, where $M$ is the maximal ideal of $S$. Let $\Sigma$ be the set of ideals $J_1$ in $S$ such that (i) $J\supset J_1\supset MJ$ and (ii) there exists a $\pi_1$-morphism $a_1\to a$ in $A$, where $\pi_1:S/J_1\to R$ is the canonical map. The property 1.2(1) of $A$, together with the fact that $J/MJ$ is a finite-dimensional vector space over $k$, entails that there exists a smallest ideal in the family $\Sigma$. Let it be $J'$ and choose a $\pi'$-morphism $a'\to a$, where $\pi':S/J'\to R$ is the canonical map. Note that $\bar\Omega_{a'}\to\bar\Omega_a$ is an isomorphism since $J'\subset J\subset M^2$. We claim that $a'\to a$ is versal for small prolongations of $a$. Indeed, let $a_1\to a$ be a small prolongation. If we put $R_1=P(a_1)$, we obtain an exact commutative diagram
\[
\begin{tikzcd}
0 \arrow[r] & J/MJ \arrow[r] \arrow[d,"h"'] & S/MJ \arrow[r] \arrow[d,"h_1"] & R \arrow[r] \arrow[d,equal] & 0\\
0 \arrow[r] & k \arrow[r] & R_1 \arrow[r,"\pi_1"'] & R \arrow[r] & 0.
\end{tikzcd}
\]
If $h=0$, then $\pi_1:R_1\to R$ admits a cross-section so that $h_1$ induces $f:S/J'\to R_1$. If $h\neq0$, then $h_1:S/MJ\to R_1$ is surjective and hence $h_1$ induces $f:S/J'\to R_1$ by the definition of the ideal $J'$. Therefore in either case we obtain a commutative diagram
\[
\begin{tikzcd}
S/J' \arrow[r,"f"] \arrow[dr,"\pi'"'] & R_1 \arrow[d,"\pi_1"]\\
& R .
\end{tikzcd}
\]
Since $a$ is versal for small prolongations of $e$, it follows from 1.6 that we obtain a commutative diagram
\[
\begin{tikzcd}
 a' \arrow[r,"\rho"] \arrow[dr] & a_1 \arrow[d]\\
& a .
\end{tikzcd}
\]

Let $P:A\to\CL$ be a cofibered category (in groupoids). It is then clear that we obtain a cofibered category
\[
 \pro(P):\pro(A)\longrightarrow\pro(\CL)
\]
(in groupoids). For the definition of the category of pro-objects $\pro(A)$, see A. Grothendieck, S\'eminaire Bourbaki 195.
\[
\begin{tikzcd}
 A \arrow[r] \arrow[d,"P"'] & \pro(A) \arrow[d,"\pro(P)"]\\
 \CL \arrow[r] & \pro(\CL).
\end{tikzcd}
\]
Now let $\widehat{\CL}$ be the full subcategory of $\pro(\CL)$ in which
\[
 \ob(\widehat{\CL})=\{(R_i)_{i\in\mathbb N},\ R_i\in\ob(\CL),\ \varphi_i:R_i\to R_{i-1}\text{ is surjective for all }i,
\]
and induces isomorphisms
\[
 \mathfrak m_i/(\ma+\mathfrak m_i^2)\xrightarrow{\sim}
 \mathfrak m_{i-1}/(\ma+\mathfrak m_{i-1}^2)
\]
for all sufficiently large $i\}$. Here $\mathbb N$ stands for the ordered category of natural numbers, and $\mathfrak m_i$ is the maximal ideal of $R_i$. We set $\widehat A$ to be the full subcategory of $\pro(A)$ in which
\[
 \ob(\widehat A)=\{a\in\ob(\pro(A))\mid \pro(P)(a)\in\ob(\widehat{\CL})\}.
\]
We then obtain a cofibered category
\[
 \widehat P:\widehat A\longrightarrow\widehat{\CL}
\]
(in groupoids) with a commutative diagram
\[
\begin{tikzcd}
 A \arrow[r,hook] \arrow[d,"P"'] & \widehat A \arrow[d,"\widehat P"]\\
 \CL \arrow[r,hook] & \widehat{\CL}.
\end{tikzcd}
\]
We note that the category $\widehat{\CL}$ is canonically equivalent to the category of complete local $\Lambda$-algebras of formally finite type over $\Lambda$ with fixed residue field $k$. Henceforth we shall make such identification.

\textbf{Definition 1.9.} Let $A$ be a cofibered groupoid over $\CL$. An object $\xi$ in $\widehat A$ is called a quasi-initial object of $\widehat A$ if for every object $a$ in $\widehat A$ there exists a map $\xi\to a$ in $\widehat A$. A quasi-initial object $\xi$ of $\widehat A$ is called minimal if the morphism $\xi:h_R\to[\widehat A]$ induces a bijection
\[
 \xi(k[\varepsilon]):\Hom_{\Lambda\text{-}\alg}(R,k[\varepsilon])\longrightarrow [A(k[\varepsilon])],
\]
where $R=P(\xi)$. An object $\xi$ in $\widehat A$ is called a versal formal object of $A$ if every surjective map $\alpha:a'\to a$ in $A$ induces a surjective map
\[
 \Hom_{\widehat A}(\xi,a')\longrightarrow\Hom_{\widehat A}(\xi,a).
\]
A versal formal object $\xi$ of $A$ is called minimal if the morphism $\xi:h_R\to[\widehat A]$ induces a bijection
\[
 \xi(k[\varepsilon]):\Hom_{\Lambda\text{-}\alg}(R,k[\varepsilon])\longrightarrow [A(k[\varepsilon])],
\]
where $R=P(\xi)$.

\textbf{Proposition 1.10.} Let $A$ be a cofibered groupoid over $\CL$.
\begin{enumerate}[label=(\roman*)]
\item Every (minimally) versal formal object of $A$ is a quasi-initial (minimal) object of $\widehat A$.
\item If $\xi$ is an object in $\widehat A$ such that
\[
 \xi(k[\varepsilon]):\Hom_{\Lambda\text{-}\alg}(R,k[\varepsilon])\longrightarrow [A(k[\varepsilon])]
\]
is injective, where $P(\xi)=R$, then every endomorphism of $\xi$ in $\widehat A$ is an automorphism. In particular a quasi-initial minimal object of $\widehat A$, if it exists, is unique up to non-canonical isomorphism.
\end{enumerate}

\textbf{Proof.} (i) Let $\xi$ be a versal formal object of $A$. We must show that, for any object $a$ in $\widehat A$, we have a map $\xi\to a$. By the definition of $\widehat A$, there exists a sequence of surjective maps in $A$
\[
 \cdots\longrightarrow a_n\xrightarrow{\alpha_n}a_{n-1}\longrightarrow\cdots\longrightarrow a_1\xrightarrow{\alpha_1}a_0
\]
such that $a=\lim a_n$. Assume that there exist maps $\varphi_i:\xi\to a_i$ such that $\alpha_i\varphi_i=\varphi_{i-1}$ for all $i<n$. Since $\xi$ is a versal formal object of $A$ and $\alpha_n:a_n\to a_{n-1}$ is surjective in $A$, we obtain a map $\xi\xrightarrow{\varphi_n}a_n$ such that $\alpha_n\varphi_n=\varphi_{n-1}$. If we set $\varphi=\lim \varphi_n$, we have $\varphi:\xi\to a$.

(ii) Let $\xi\xrightarrow{\alpha}\xi$ be a map in $\widehat A$. Then, for every $\lambda\in\Hom_{\Lambda\text{-}\alg}(R,k[\varepsilon])$, we have
\[
 \lambda\xi=\lambda\alpha\xi=(\lambda\cdot P(\alpha))\xi,
\]
and hence by the hypothesis on $\xi$ we have $\lambda=\lambda\cdot P(\alpha)$ for all $\lambda\in\Hom_{\Lambda\text{-}\alg}(R,k[\varepsilon])$. Therefore $P(\alpha):R\to R$ is surjective. However, every surjective endomorphism of a noetherian ring is an automorphism; hence $P(\alpha)$ is an automorphism of $R$, and therefore $\alpha$ is an automorphism of $\xi$.

\textbf{Theorem 1.11 (Schlessinger).} Let $A$ be a cofibered groupoid over $\CL$. Then $A$ admits a minimally versal formal object if and only if:
\begin{enumerate}[label=(\roman*)]
\item $A$ is semi-homogeneous;
\item $[A(k[\varepsilon])]$ is a finite-dimensional vector space over $k$ (cf. 1.3(c)).
\end{enumerate}

\textbf{Proof.} Necessity. Let $\xi$ in $\widehat A$ be a minimally versal formal object of $A$, and consider a diagram
\[
\begin{tikzcd}
 a'' \arrow[r,dashed] \arrow[d] & a' \arrow[d]\\
 a_1 \arrow[r] & a
\end{tikzcd}
\]
in $A$ where $q$ is surjective. Since $\xi$ is a quasi-initial object of $\widehat A$, there exists a map $\xi\to a_1$. Since $q$ is surjective, the composite map $\beta=\beta_1\rho$ can be lifted to $a'$, i.e. we obtain a commutative diagram
\[
\begin{tikzcd}
& \xi \arrow[dl] \arrow[d] \\
a_1 \arrow[r] & a' \arrow[d]\\
& a .
\end{tikzcd}
\]
This proves condition 1.2(1).

We now prove condition 1.2(2). Let $S\times_k k[\varepsilon]\mathrel{\substack{\longrightarrow\\[-.6ex]\longrightarrow}} S$ be the projections, and let
\[
\begin{tikzcd}
 a' \arrow[r,shift left,"u"] \arrow[r,shift right,"v"'] & b'
\end{tikzcd}
\]
be the two maps in $A$. By the definition of $\xi$, we obtain a commutative diagram
\[
\begin{tikzcd}
& \xi \arrow[dl] \arrow[dr] &\\
a' \arrow[rr,shift left,"u"] \arrow[rr,shift right,"v"'] && b' .
\end{tikzcd}
\]
In particular $\pi_1P(u)=\pi_1P(v)$. Assume now that $(\pi_2)^*a'=(\pi_2)^*b'$. Then $\pi_2P(u)=\pi_2P(v)$ because of the minimality of $\xi$, and hence
\[
 P(u)=\pi_1P(u)\times\pi_2P(u)=\pi_1P(v)\times\pi_2P(v)=P(v).
\]
Since $A$ is a cofibered groupoid, we obtain a unique isomorphism $a'\xrightarrow{\sim}b'$ in $A(S\times_k k[\varepsilon])$ such that $\beta u=v$. Then $\alpha u=\alpha v$ entails that $\beta\alpha=\alpha$, since $A$ is a cofibered category.

Sufficiency. Let $e_1,\ldots,e_n$ be objects in $A(k[\varepsilon])$ which form a $k$-basis of the vector space $[A(k[\varepsilon])]$, and choose $a_1$ in $A(k[\varepsilon]\times_k\cdots\times_k k[\varepsilon])$ such that $(\pi_i)_*a_1=e_i$. Then
\[
 a_1(k[\varepsilon]):\Hom_{\Lambda\text{-}\alg}(R_1,k[\varepsilon])\longrightarrow [A(k[\varepsilon])]
\]
is bijective, where $R_1=P(a_1)$; in particular $a_1\to e$ is versal for small prolongations of $e$. Choose $a_2\to a_1$ which is versal for small prolongations of $a_1$, and such that $\bar\Omega_{a_2}\to\bar\Omega_{a_1}$ is bijective. We then choose $a_3\to a_2$, versal for small prolongations of $a_2$, such that $\bar\Omega_{a_3}\to\bar\Omega_{a_2}$ is bijective, and so on, and set
\[
 \xi=\lim_n a_n.
\]
We claim that $\xi$ is a minimally versal formal object of $A$. Indeed, since $\bar\Omega_{a_n}\to\bar\Omega_{a_{n-1}}$ is bijective for all $n$, $\xi(k[\varepsilon])$ is bijective. Now consider a diagram
\[
\begin{tikzcd}
& a' \arrow[d,"q"]\\
\xi \arrow[r,"h"'] \arrow[ur,dashed] & a
\end{tikzcd}
\]
where $q$ is surjective in $A$. We must show that $h$ can be lifted to $a'$. For this, one may assume, by induction, that $a'$ is a small prolongation of $a$. Since $P(a)$ is an artinian ring, $h:\xi\to a$ is factored as a composite map $\xi\to a_n\to a$ for some $a_n$. Since $A$ is semi-homogeneous, we get a diagram
\[
\begin{tikzcd}
 a'' \arrow[r] \arrow[d,"q'"'] & a' \arrow[d,"q"]\\
 a_n \arrow[r] & a
\end{tikzcd}
\]
where $q'$ is also a small prolongation. We then obtain a commutative diagram
\[
\begin{tikzcd}
& a'' \arrow[d]\\
\xi \arrow[r] \arrow[ur,dashed] & a_n,
\end{tikzcd}
\]
and therefore we are done.

\textbf{Remark 1.12.} Let
\[
 \cdots\longrightarrow a_n\longrightarrow\cdots\longrightarrow a_2\longrightarrow a_1\longrightarrow e
\]
be an infinite sequence of arrows in $A$ such that $a_i$ is versal for small prolongations of $a_{i-1}$, $P(a_i)\to P(a_{i-1})$ is surjective, and $\bar\Omega_{a_i}\to\bar\Omega_{a_{i-1}}$ is bijective for $i\gg0$. Then the above proof shows that $\lim a_i$ is a versal formal object of $A$.

\textbf{Remark 1.13.} Let $A$ be a cofibered groupoid over $\CL$. If $A$ has the properties 1.11(i) and (ii), then so does $[A]$, and $\xi\in\ob(\widehat A)$ is a minimally versal formal object of $A$ if and only if it is so for $[A]$, by 1.10(ii).

\textbf{Remark 1.14.} Let $A$ be a cofibered groupoid over $\CL$, admitting a minimally versal formal object $\xi$. Then an invariant of $R=P(\xi)$, which is a complete local $\Lambda$-algebra of formally finite type over $\Lambda$, is automatically an invariant of $A$ over $\CL$. Therefore, for example, we shall say that $A$ is formally smooth over $\CL$ if $R$ is formally smooth over $\Lambda$, i.e. a formal power-series algebra over $\Lambda$, or we shall say that $A$ is irreducible over $\CL$ if $\Spec(R)$ is irreducible, etc. In the next section we introduce the notion of smoothness of a functor $A\to B$ over $\CL$, where $A,B$ are cofibered categories in groupoids over $\CL$. It turns out (cf. 2.4) that ``$A$ is formally smooth over $\CL$'' is equivalent to ``$P:A\to\CL$ is a smooth functor'', and therefore no confusion arises. We also define the dimension of $A$ over $\CL$ by
\[
 \dim_\Lambda A=\dim_k R.
\]
For this we note the following fact. Let $\mathcal C_k$ be the category of artinian local $k$-algebras with residue field $k$. Then $\mathcal C_k\to\CL$ is a fully faithful embedding whose left adjoint functor $\CL\to\mathcal C_k$ is given by
\[
 R\longmapsto k\otimes_\Lambda R.
\]
Let $A|\mathcal C_k$ be the pull-back of $P:A\to\CL$ under $\mathcal C_k\to\CL$. Then it is trivial to see that if $\xi$ is a minimally versal formal object for $A$ over $\CL$, then $k\otimes_\Lambda\xi$ is a minimally versal formal object for $A|\mathcal C_k$ over $\mathcal C_k$, where $\xi\to k\otimes_\Lambda\xi$ is a cocartesian arrow above $R\to k\otimes_\Lambda R$. It follows that
\[
 \dim_\Lambda A=\dim_k(A|\mathcal C_k).
\]

\textbf{1.15.} In the remainder of this paragraph, we briefly explain how to modify the preceding theory in order to allow residue-field extensions.

Let $\Lambda$ and $K$ be complete noetherian local rings, with $K$ a $\Lambda$-algebra. We assume that:
\begin{enumerate}[label=(\alph*)]
\item the natural map from $\Lambda$ to $K$ is a local homomorphism;
\item the residue field $K'$ of $K$ is an extension of finite type of the residue field $k$ of $\Lambda$.
\end{enumerate}
The most interesting case is when $K=K'$ and $K$ is a finite inseparable extension of $k$.

We denote by $\widehat{\mathcal C}_{\Lambda,K}$ the category whose objects are the complete noetherian local rings $R$ provided with a structure of $\Lambda$-algebra and with a $\Lambda$-epimorphism $s:R\to K$.

We denote by $\CLK$ the subcategory of $\widehat{\mathcal C}_{\Lambda,K}$ consisting of those $(R,s)$ such that $\Ker(s)$ is an $R$-module of finite length. The category $\CLK$ can be identified with a full subcategory of $\pro(\CL)$. If $K$ is artinian, so is any object in $\CLK$.

For any finite-dimensional vector space $V$ over $K'$, we denote by $D_V(K)$ the algebra of dual numbers over $K$ defined by $V$; its elements are of the form $\kappa+v\varepsilon$ with $\kappa\in K$, $v\in V$ and $\varepsilon^2=0$, and
\[
 (\kappa+v\varepsilon)(\kappa'+v'\varepsilon)=\kappa\kappa'+(\kappa v'+\kappa'v)\varepsilon.
\]
The map $e:D_V(K)\to K$, $e(\kappa+v\varepsilon)=\kappa$, turns $D_V(K)$ into an object of $\CLK$. The role played before by $k[\varepsilon]$ will be played here by $D_{K'}(K)$.

As explained in 1.8, any cofibered groupoid $A$ over $\CLK$ extends in a natural way to a cofibered category $\widehat A_K$ over $\widehat{\mathcal C}_{\Lambda,K}$. A map in $A$ or in $\widehat A_K$ will be called surjective if the corresponding homomorphism of algebras in $\CLK$ or in $\widehat{\mathcal C}_{\Lambda,K}$ is surjective.




% SGA7 batch 100 macro additions
\providecommand{\Fl}{\operatorname{Fl}}
\providecommand{\coker}{\operatorname{coker}}
\providecommand{\Exal}{\operatorname{Exal}}


\subsection*{Continuation of 1.15: residue-field extensions}

We assume that \(A(K)=\{e\}\).

\textbf{Definition 1.16.} A cofibered groupoid \(A\) over \(\mathcal C_{\Lambda,K}\) is called
\emph{semi-homogeneous} if the following properties hold.

\begin{enumerate}[label=(\Roman*)]
\item Every diagram of solid arrows in \(A\)
\[
\begin{tikzcd}
a'' \arrow[r,dashed] \arrow[d,dashed] & a' \arrow[d]\\
a_1 \arrow[r] & a
\end{tikzcd}
\]
where \(a'\to a\) is surjective, can be completed to a commutative diagram including
the dotted arrows.

\item For every \(R\), the natural map
\[
 [A(R\times_K D_{K'}(K))]
 \longrightarrow
 [A(R)]\times [A(D_{K'}(K))]
\]
is bijective, where the right side is the corresponding fibre product of isomorphism
classes.
\end{enumerate}
If \(A\) is semi-homogeneous, so is \([A]\).

\textbf{1.17.} Let us first restrict \(A\) to the full subcategory of
\(\mathcal C_{\Lambda,K}\) whose objects are the \(D_V(K)\). Let
\[
 \Omega=\Omega^1_{K/\Lambda}\otimes_K K',
\]
and let \(d:K\to \Omega\) be the derivation. A morphism
\(f:D_V(K)\to D_W(K)\) can be written
\[
 f(k+v\varepsilon)=k+(u(v)+D\,dk)\varepsilon,
\]
with \(u\in\Hom_{K'}(V,W)\) and \(D\in\Hom_{K'}(\Omega,W)\). If \(f\) is identified
with the pair \((u,D)\), the composition law is
\[
 (u,D)\circ(u',D')=(uu',\,D+uD').
\]
Condition (II) entails the following condition:
\[
\tag{II'}
 V\longmapsto [A(D_V(K))]
\]
commutes with products as a functor from \(K'\)-vector spaces to sets.

\textbf{Lemma 1.18.} Assume that condition \((\mathrm{II}')\) is satisfied by \(A\). Then:

\begin{enumerate}[label=(\arabic*)]
\item \([A(D_{K'}(K))]\) carries a unique \(K'\)-vector-space structure such that, functorially in
\(V\),
\[
 [A(D_V(K))]\simeq V\otimes_{K'}[A(D_{K'}(K))].
\]
\item There is a linear map
\[
 \gamma:\Hom_{K'}(\Omega,K')\longrightarrow [A(D_{K'}(K))]
\]
such that, for any \(D\in\Hom_{K'}(\Omega,K')\), the endomorphism induced by
\((1,D):D_{K'}(K)\to D_{K'}(K)\) is
\[
 x\longmapsto x+\gamma(D).
\]
\item More generally, for any \((u,D):D_V(K)\to D_W(K)\), the corresponding map
\[
 [A(D_V(K))]\simeq V\otimes[A(D_{K'}(K))]
 \longrightarrow
 [A(D_W(K))]\simeq W\otimes[A(D_{K'}(K))]
\]
is
\[
 x\longmapsto u(x)+\gamma(D).
\]
\end{enumerate}

\emph{Proof.} Put
\[
 F(V)=[A(D_V(K))],\qquad E=F(K').
\]
Condition \((\mathrm{II}')\), together with the fact that \(K'\) is a vector-space
object in the category of finite-dimensional \(K'\)-vector spaces, gives the vector-space
structure on \(F(K')\) and the functorial identification \(F(V)\simeq V\otimes E\).

We write \(F(u,D)\) for the map induced by \((u,D)\). Since
\[
 (u,D)=(1,D)\circ(u,0),
\]
one has
\[
 F(u,D)=F(D)\circ(u\otimes 1_E).
\]
The functoriality of \(F\) amounts to
\[
 F(0)=\operatorname{id},\qquad
 F(D)\circ F(D')=F(D+D'),\qquad
 (u\otimes 1_E)\circ F(D)=F(uD)\circ(u\otimes 1_E).
\]
For \(V=K'^n\), using the maps \(i:V\to V\oplus K'^n\) and
\(p:V\oplus K'^n\to K'^n\), 
The commutative diagram used in this reduction is
\[
\begin{tikzcd}[column sep=large,row sep=large]
V\otimes E \arrow[r] \arrow[d,"F(D)"'] &
(V\oplus K'^n)\otimes E \arrow[r] \arrow[d,"{F(i,D)}"] &
K'^n\otimes E \arrow[d,equal]\\
V\otimes E \arrow[r] &
(V\oplus K'^n)\otimes E \arrow[r] &
K'^n\otimes E .
\end{tikzcd}
\]

 one sees that \(F(D)\) is a translation:
\[
 F(D)(y)=y+\alpha(D).
\]
The identities above imply
\[
 \alpha(D+D')=\alpha(D)+\alpha(D'),\qquad
 \alpha(uD)=u\,\alpha(D),
\]
and hence \(\alpha(D)\) is defined by a unique linear map
\(\gamma:\Hom(\Omega,K')\to E\). This proves the assertions.

\textbf{Definition 1.19.} Let \(\xi\) be a formal object of \(A\).

\begin{enumerate}[label=(\arabic*)]
\item The object \(\xi\) is called \emph{versal} if every surjective map \(a'\to a\) in \(A\)
induces a surjective map
\[
 \Hom_A(\xi,a')\longrightarrow \Hom_A(\xi,a).
\]
\item A versal formal object \(\xi\), with image \((R,\xi)\), is called \emph{minimal} if,
for any two maps \(\varphi,\psi:R\to D_V(K)\), whenever \(\varphi(\xi)\) and
\(\psi(\xi)\) are isomorphic there exists a derivation \(D:K\to V\) such that
\[
 (\psi,D)\circ \xi=\xi,
\]
or equivalently \(\psi=\varphi-D\circ d\).
\end{enumerate}

\textbf{Theorem 1.20.} If \(A\) is semi-homogeneous and
\[
 \dim_{K'}[A(D_{K'}(K))]<\infty,
\]
then \(A\) admits a minimally versal formal object, and any two such objects are
isomorphic.

\emph{Proof.} With the notation of Lemma 1.18, let \(H^\ast\) be a subspace of
\([A(D_{K'}(K))]\) supplementary to the image of \(\gamma\), and let \(H\) be the dual
of \(H^\ast\). Put
\[
 R_1=D_H(K).
\]
Choose \(\xi_1\in A(R_1)\) whose image in
\[
 [A(R_1)]=H^\ast\otimes [A(D_{K'}(K))]
       =\Hom\bigl(H^\ast,[A(D_{K'}(K))]\bigr)
\]
is the inclusion of \(H^\ast\) in \([A(D_{K'}(K))]\). Then
\[
 \Hom_{\mathcal C_{\Lambda,K}}(D_{K'}(K),D_V(K))
   \longrightarrow [A(D_V(K))]
\]
is surjective, and the minimality condition of Definition 1.19 is valid for
\((R_1,\xi_1)\).

Choose \(S\in\widehat{\mathcal C}_{\Lambda,K}\), formally smooth over \(\Lambda\), and an
epimorphism
\[
 \pi:S\longrightarrow R_1
\]
such that \(\Ker(\pi)\cap\mathfrak m_S^2=0\). If \(\mathfrak m\) is the maximal ideal
of \(S\), define inductively, for \(n\geq 1\),
\[
 R_n=S/\mathfrak m^n,\qquad \xi_n\in A(R_n),
\]
by requiring that:
\begin{enumerate}[label=(\alph*)]
\item \(\mathfrak m^n\) be the intersection of all ideals \(I\supset\mathfrak m^n\) of \(S\)
such that \(\xi_n\) extends to an object of \(A(S/I)\);
\item \(\xi_n\) extend \(\xi_{n-1}\).
\end{enumerate}
For
\[
 R=\varprojlim R_n,\qquad \xi=\varprojlim \xi_n,
\]
one verifies that \((R,\xi)\) is a minimally versal formal object. The construction also gives
uniqueness up to isomorphism.

\section*{2. Pro-representable cofibered groupoids}

Let \(A\) be a cofibered groupoid over \(\mathcal C_\Lambda\), admitting a minimally
versal formal object \(\xi\). In general, \(\xi\) alone does not allow one to recover
\(A\), up to \(\mathcal C_\Lambda\)-equivalence.

Many cofibered groupoids which occur in practice satisfy a stronger property than
semi-homogeneity. To study them, it is convenient first to rephrase the preceding
construction in terms of smooth functors. The notation and terminology of §1 remain
in force.

\textbf{Definition 2.1.} Let
\[
 \varphi:A\longrightarrow B
\]
be a functor of cofibered groupoids over \(\mathcal C_\Lambda\). We say that
\(\varphi\) is \emph{smooth} if every surjective map
\(\beta:b'\to\varphi(a)\) in \(B\) can be completed to a commutative diagram
\[
\begin{tikzcd}
\varphi(a') \arrow[r] \arrow[d] & b' \arrow[d,"\beta"]\\
\varphi(a) \arrow[r,equal] & \varphi(a)
\end{tikzcd}
\tag{\(*\)}
\]
for some \(a'\to a\) in \(A\). We say that \(\varphi\) is \emph{minimally smooth}
if it is smooth and the induced map
\[
 [A(k[\varepsilon])]\longrightarrow [B(k[\varepsilon])]
\]
is bijective.

\textbf{Remark 2.2.}
\begin{enumerate}[label=(\alph*)]
\item The condition that \(\varphi\) be smooth is equivalent to asking that, in the preceding
diagram, \(\varphi(a')\to b'\) may be chosen to be a \(\mathcal C_\Lambda\)-isomorphism; it is
also equivalent to imposing the lifting condition only for small prolongations
\(b'\to\varphi(a)\).

\item For every cofibered groupoid \(A\), the natural functor
\[
 A\longrightarrow [A]
\]
is minimally smooth.

\item If \(R\in\widehat{\mathcal C}_\Lambda\), put
\[
 h_R(S)=\Hom_{\Lambda\text{-}\mathrm{alg}}(R,S)\qquad(S\in\mathcal C_\Lambda).
\]
This is a discrete cofibered groupoid over \(\mathcal C_\Lambda\). A map \(R\to S\) in
\(\widehat{\mathcal C}_\Lambda\) induces \(h_S\to h_R\); the latter is smooth if and
only if \(S\) is a formal power-series ring over \(R\).

\item Given \(\xi\in A\) with \(P(\xi)=R\), choose, for each \(f:R\to S\), a cocartesian
image \(\xi(f)\). This defines a functor
\[
 \xi^\ast:h_R\longrightarrow A.
\]
It is smooth precisely when \(\xi\) is a versal formal object; it is minimally smooth
precisely when \(\xi\) is minimally versal.
\end{enumerate}

\textbf{Proposition 2.3.} For functors of cofibered groupoids
\[
 A\xrightarrow{\varphi}B\xrightarrow{\psi}C
\]
over \(\mathcal C_\Lambda\):
\begin{enumerate}[label=(\alph*)]
\item if \(\varphi\) and \(\psi\) are smooth, then \(\psi\varphi\) is smooth;
\item if \(\psi\) and \(\psi\varphi\) are smooth, then \(\varphi\) is smooth;
\item if \(\psi\varphi\) is smooth and \(\varphi\) is surjective on isomorphism classes
of objects, then \(\psi\) is smooth.
\end{enumerate}

\textbf{Remark 2.4.} If \(A\) has a versal formal object \(\xi\), then
\(\xi^\ast:h_R\to A\) is smooth, where \(R=P(\xi)\). Hence
\[
 A\text{ is formally smooth over }\mathcal C_\Lambda
\]
is equivalent to saying that \(R\) is a formal power-series algebra over \(\Lambda\),
and also to saying that the structural functor \(A\to\mathcal C_\Lambda\) is smooth.

\textbf{Definition 2.5.} A cofibered groupoid \(A\) over \(\mathcal C_\Lambda\) is called
\emph{homogeneous} if every diagram
\[
\begin{tikzcd}
a' \arrow[r] \arrow[d] & a \arrow[d]\\
b' \arrow[r] & b
\end{tikzcd}
\]
where \(a'\to a\) is surjective admits a fibre product \(a'\times_a b'\) in \(A\) such
that
\[
 P(a'\times_a b')=P(a')\times_{P(a)}P(b').
\]

\textbf{Remark 2.6.}
\begin{enumerate}[label=(\alph*)]
\item Homogeneity implies semi-homogeneity. Indeed, for any solid diagram in \(A\) above
a diagram in \(\mathcal C_\Lambda\), the fibre-product condition supplies the missing
object and the missing arrows whenever the relevant map in \(\mathcal C_\Lambda\) is
surjective.

\item The usual fibre product of categories \(X\times_ZY\) has objects \((x,y)\) with
\(\Phi(x)=\Psi(y)\). The fibre product in the sense of \(2\)-categories has objects
\((x,y;u)\), where \(u:\Phi(x)\simeq\Psi(y)\) is an isomorphism, and arrows are pairs
which make the evident square commute:
\[
\begin{tikzcd}
\Phi(x) \arrow[r,"u"] \arrow[d] & \Psi(y) \arrow[d]\\
\Phi(x') \arrow[r,"u'"'] & \Psi(y').
\end{tikzcd}
\]
Thus the notions of homogeneous and semi-homogeneous groupoids, and of smooth functors,
may be expressed by requiring certain functors to such \(2\)-categorical fibre products
to be equivalences, or to be essentially surjective on isomorphism classes.

\item If \(A\) is homogeneous, then for finite-dimensional \(k\)-vector spaces \(V,W\), the
functor
\[
 A(k[V\oplus W])\longrightarrow A(k[V])\times A(k[W])
\]
is an equivalence. In particular \(\Aut(1_V)\) and \([A(k[\varepsilon])]\) have canonical
\(k\)-vector-space structures.

\item For every \(R\in\mathcal C_\Lambda\), the discrete groupoid \(h_R\) is homogeneous.
If \(A\) is semi-homogeneous, then \([A]\) is semi-homogeneous; homogeneity of \(A\) does
not, in general, imply homogeneity of \([A]\).
\end{enumerate}

\textbf{Proposition 2.7.} Let \(A\) be a homogeneous cofibered groupoid over
\(\mathcal C_\Lambda\). The following conditions are equivalent:
\begin{enumerate}[label=(\roman*)]
\item \([A]\) is homogeneous, that is,
\[
 [A(R'\times_R S)]\longrightarrow [A(R')]\times_{[A(R)]}[A(S)]
\]
is bijective for every small extension \(R'\to R\).
\item For every small prolongation \(a'\to a\) in \(A\), the map
\[
 \Aut_{A(R')}(a')\longrightarrow \Aut_{A(R)}(a)
\]
is surjective, where \(R'=P(a')\) and \(R=P(a)\).
\end{enumerate}

\emph{Proof.} The implication (i)\(\Rightarrow\)(ii) is obtained by applying homogeneity of
\([A]\) to the two objects \(a'\times_a a'\) and \(a'\times_a a'_\tau\), where
\(\tau\in\Aut_A(a)\). The reverse implication follows by comparing two objects over
\(R'\times_R S\) whose images over \(R'\) and \(S\) are isomorphic; the assumed lifting
of automorphisms permits the comparison isomorphisms to be modified so that the two
squares commute simultaneously.

The comparison is represented schematically by the solid-arrow diagram
\[
\begin{array}{ccccc}
 & a''_2 & & R'\times_R S &\\
 \swarrow & \downarrow & \searrow & \swarrow & \searrow\\
a'_2 & a''_1 & b_1 \xrightarrow{\tau} b_2 & R' & S\\
 & \swarrow & \searrow & \searrow & \swarrow\\
 & a'_1 & a_1 & & R ,
\end{array}
\]
with the left part in \(A\) and the right diamond in \(\mathcal C_\Lambda\).


Let \(R'\xrightarrow f R\) be a small extension. There is a canonical map
\[
 \Ker(f)\longrightarrow R'\times_R R',
\qquad
(r_1,r_2)\longmapsto r_1(0)+(r_2-r_1),
\]
and the corresponding map
\[
D_f:\Ker(f)\longrightarrow R'\times_R R'[\Ker(f)]
\]
is an isomorphism. If \(a_i\to a\) (\(i=1,2\)) are two \(f\)-maps in \(A\), we write
\[
 \tau(a_1,a_2)\in A(k[\Ker(f)])
\]
for the object obtained by transporting \(a_1\times_a a_2\) through this identification.
Then
\[
 a_1\times_a a_2\simeq a_1\times \tau(a_1,a_2),
\]
compatibly with the projection onto the first factor.

\textbf{Lemma 2.8.} Let \(A\) be homogeneous over \(\mathcal C_\Lambda\), and let
\(a_i\to a\) (\(i=1,2\)) be \(f\)-maps in \(A\).
\begin{enumerate}[label=(\arabic*)]
\item There exists \(p:a_1\to a_2\) in \(A\), with \(a_2p=a_1\), if and only if there exists
an arrow \(a_1\to\tau(a_1,a_2)\) in \(A\).
\item There exists such a \(p\) with \(P(p)=\operatorname{id}\) if and only if
\[
 \tau(a_1,a_2)=0\quad\text{in }[A(k[\Ker(f)])].
\]
\end{enumerate}

\emph{Proof.} This is the precise form of the preceding fibre-product description. Such a
\(p\) is equivalent to a section of
\[
 a_1\times\tau(a_1,a_2)\longrightarrow a_1,
\]
and the condition \(P(p)=\operatorname{id}\) is exactly the condition that the corresponding
arrow to \(k[\Ker(f)]\) factor through the zero section.

\textbf{Proposition 2.9.} Let
\[
 A\xrightarrow{\varphi}B\xrightarrow{\psi}C
\]
be functors of cofibered groupoids over \(\mathcal C_\Lambda\).
\begin{enumerate}[label=(\alph*)]
\item Assume that \(B\) is semi-homogeneous, that \(\psi\) is homogeneous, and that
\[
 [A(k[\varepsilon])]\longrightarrow [C(k[\varepsilon])]
\]
is injective. Then \(\psi\varphi\) is smooth if and only if \(\varphi\) and \(\psi\) are both
smooth.
\item If \(\varphi:A\to B\) is smooth and \(A\) is semi-homogeneous, then \(B\) is
semi-homogeneous.
\end{enumerate}

\emph{Proof.} For (a), assume \(\psi\varphi\) is smooth. Since smoothness descends as in
Proposition 2.3, it remains to lift a small prolongation \(b'\to\varphi(a)\). Smoothness of
\(\psi\varphi\) gives a lift in \(A\) after applying \(\psi\), and the semi-homogeneity of \(B\),
together with the injectivity on tangent classes, removes the obstruction.

In the proof of smoothness one completes, by dotted arrows, the diagram
\[
\begin{array}{ccccc}
\varphi(a'_{n+1}) & \dashrightarrow & \varphi(a'_n)\\
\downarrow & \searrow & \downarrow\\
\varphi(a_{n+1}) & \longrightarrow & \varphi(a_n)
\end{array}
\qquad
\begin{array}{ccc}
b'_{n+1} & \longrightarrow & b'_n\\
\downarrow{\scriptstyle \beta_{n+1}} && \downarrow{\scriptstyle \beta_n}\\
\varphi(a_{n+1}) & \longrightarrow & \varphi(a_n).
\end{array}
\]
After using the semi-homogeneity of \(A\), the corresponding diagram in \(B\) is a fibre
product; this gives the required lift.

Conversely, if
\(\varphi\) and \(\psi\) are smooth, their composite is smooth. Part (b) is obtained by applying
the lifting condition defining smoothness to the semi-homogeneous diagrams in \(A\).

\textbf{Corollary 2.10.}
\begin{enumerate}[label=(\alph*)]
\item Let \(\varphi:A\to B\) be a smooth functor of cofibered groupoids over
\(\mathcal C_\Lambda\). Let \(\xi\) be a versal formal object of \(A\), and let \(\eta\) be a
minimally versal formal object of \(B\). If \(B\) is homogeneous, then for any map
\[
 \beta:\eta\longrightarrow \varphi(\xi)
\]
in \(\widehat B\), the ring \(P(\xi)\) is a formal power-series ring over \(P(\eta)\) via
\[
 P(\beta):P(\eta)\longrightarrow P(\varphi(\xi))=P(\xi).
\]

\item Let \(A\) be a homogeneous cofibered groupoid over \(\mathcal C_\Lambda\), and let
\(\xi,\eta\) be versal formal objects of \(A\). If \(\eta\) is minimal, then for any map
\[
 \alpha:\eta\longrightarrow \xi
\]
in \(\widehat A\), the ring \(P(\xi)\) is a formal power-series ring over \(P(\eta)\) through
\(P(\alpha)\).
\end{enumerate}

\emph{Proof.} A map \(\beta:\eta\to\varphi(\xi)\) induces the commutative square
\[
\begin{tikzcd}
h_{P(\xi)} \arrow[r,"\widetilde\beta"] \arrow[d,"\widetilde\xi"'] &
h_{P(\eta)} \arrow[d,"\widetilde\eta"]\\
A \arrow[r,"\varphi"'] & B .
\end{tikzcd}
\]
Since \(\widetilde\xi\) and \(\varphi\) are smooth, \(\widetilde\eta\,\widetilde\beta
=\varphi\widetilde\xi\) is smooth. Applying Proposition 2.9 and the minimality of \(\eta\)
shows that \(\widetilde\beta\) is smooth; equivalently \(P(\xi)\) is a formal power-series
ring over \(P(\eta)\). Part (b) follows from (a) by taking \(A=B\).

We now consider representability of a cofibered groupoid. Let \(C\) be a category. Each
object \(R\in C\) defines the functor
\[
 h_R:C\longrightarrow(\mathrm{Ens}),\qquad h_R(S)=\Hom_C(R,S),
\]
and the Yoneda functor \(h:C^\circ\to\Hom(C,\mathrm{Ens})\) is fully faithful.

A \emph{cocategory} in \(C\) is a diagram
\[
\begin{tikzcd}[column sep=small]
h_{R_1}\times_{h_{R_0}}h_{R_1} \arrow[r,"c"] &
h_{R_1}
\arrow[r,shift left=0.55ex,"s_1"]
\arrow[r,shift right=0.55ex,"s_2"'] &
h_{R_0}
\arrow[l,bend left=50,"\varepsilon"]
\end{tikzcd}
\]
with \(s_i\varepsilon=\operatorname{id}\) for \(i=1,2\), such that \(c\) is associative and
compatible with the \(s_i\). Equivalently, the two composites
\[
 h_{R_1}\xrightarrow{\varepsilon}
 h_{R_0}\times_{h_{R_0}}h_{R_1}
 \xrightarrow{\varepsilon\times 1}
 h_{R_1}\times_{h_{R_0}}h_{R_1}
 \xrightarrow{c}h_{R_1}
\]
and
\[
 h_{R_1}\xrightarrow{\varepsilon}
 h_{R_1}\times_{h_{R_0}}h_{R_0}
 \xrightarrow{1\times\varepsilon}
 h_{R_1}\times_{h_{R_0}}h_{R_1}
 \xrightarrow{c}h_{R_1}
\]
are both identities.

If \((R_0,R_1,s_1,s_2,\varepsilon,c)\) is a cocategory in \(C\), then for every
\(S\in C\) one obtains a category \(h_{(R_0,R_1,\ldots)}(S)\) by putting
\[
 \operatorname{Ob}h_{(R_0,R_1,\ldots)}(S)=\Hom_C(R_0,S),\qquad
 \operatorname{Fl}h_{(R_0,R_1,\ldots)}(S)=\Hom_C(R_1,S),
\]
with composition induced by
\[
 c:h_{R_1}\times_{h_{R_0}}h_{R_1}\longrightarrow h_{R_1}.
\]
This defines a functor
\[
 h_{(R_0,R_1,\ldots)}:C\longrightarrow(\mathrm{categories}),
\]
hence a cofibered category over \(C\). A cogroupoid is a cocategory for which each
\(h_{(R_0,R_1,\ldots)}(S)\) is a groupoid.

In practice, cocategories and cogroupoids arise as follows. Let \(P:F\to C\) be a cofibered
category and let \(\xi\) be a fixed object of \(F\). Put \(R_0=P(\xi)\), let \((R_0,C)\)
be the category of arrows in \(C\) with source \(R_0\), and choose a cocartesian section
\(\xi^\#\) of \(F'\to(R_0,C)\) with \(\xi^\#(\operatorname{id}_{R_0})=\xi\). The functor
\[
 \Hom_\xi:(R_0\amalg R_0,C)\longrightarrow(\mathrm{Ens})
\]
is given by
\[
 X=(f_1,f_2)\longmapsto
 \Hom_{F(S)}\bigl(\xi^\#(f_2),\xi^\#(f_1)\bigr),
\]
where \(S\) is the target of \(f_1\) and \(f_2\). If this functor is represented by
\[
 X_0=(R_0 \mathrel{\substack{\xrightarrow{s_1}\\[-.6ex]\xrightarrow[s_2]{}}} R_1),
\]
then \(R_1\), together with the evident source, target, identity, and composition maps, defines
a cocategory \((R_0,R_1,\ldots)\). The resulting functor
\[
 h_{(R_0,R_1,\ldots)}\longrightarrow F
\]
is a \(C\)-equivalence if and only if \(\xi\) is quasi-initial in \(F\) over \(C\).

\textbf{Definition and Proposition 2.11.} A cogroupoid
\((R_0,R_1,\ldots)\) in \(\widehat{\mathcal C}_\Lambda\) is called \emph{smooth} if either of
the following equivalent conditions holds:
\begin{enumerate}[label=(\roman*)]
\item \(R_1\) is a formal power-series ring over \(R_0\) via \(s_1\), and also via \(s_2\);
\item the functor
\[
 h_{R_0}\longrightarrow h_{(R_0,R_1,\ldots)}
\]
is smooth over \(\mathcal C_\Lambda\).
\end{enumerate}
It is called \emph{minimally smooth}, or \emph{normalized smooth}, if it is smooth and
\(h_{(R_0,R_1,\ldots)}(k[\varepsilon])\) is totally disconnected, i.e.
\[
 s_1(k[\varepsilon])=s_2(k[\varepsilon]).
\]

\emph{Proof.} The smoothness of
\(\varepsilon:h_{R_0}\to h_{(R_0,R_1,\ldots)}\) means precisely that, for every surjective map
\((S',f')\to(S,f)\), a lifting of \(f\) to \(S'\), together with an isomorphism over \(S\),
can be lifted after replacing the isomorphism by one over \(S'\). This is exactly the smoothness
of \(s_1:h_{R_1}\to h_{R_0}\), and similarly for \(s_2\).

\textbf{Proposition 2.12.} Let \((R_0,R_1,\ldots)\) be a cogroupoid in
\(\widehat{\mathcal C}_\Lambda\).
\begin{enumerate}[label=(\alph*)]
\item If \((R_0,R_1,\ldots)\) is smooth, then \(h_{(R_0,R_1,\ldots)}\) is homogeneous.
\item If \(s_1(k[\varepsilon])=s_2(k[\varepsilon])\), then the following statements are equivalent:
\begin{enumerate}[label=(\roman*)]
\item \((R_0,R_1,\ldots)\) is smooth;
\item \(h_{(R_0,R_1,\ldots)}\) is homogeneous;
\item \(h_{(R_0,R_1,\ldots)}\) is semi-homogeneous.
\end{enumerate}
\item If \(h_{(R_0,R_1,\ldots)}\) is semi-homogeneous, then it is equivalent to
\(h_{(S_0,S_1,\ldots)}\) for some minimally smooth cogroupoid \(S_\ast\).
\end{enumerate}

\emph{Proof.} For (a), consider a diagram
\[
\begin{tikzcd}
(S_1,f_1) \arrow[dr,"{(h_1,u_1)}"'] && (S_2,f_2) \arrow[dl,"{(h_2,u_2)}"]\\
& (S,f) &
\end{tikzcd}
\]
in \(h_{(R_0,R_1,\ldots)}\), with \(h_2:S_2\to S\) surjective. Since
\(s_1:h_{R_1}\to h_{R_0}\) is smooth, \(u_1^{-1}u_2\) lifts to some
\(v\in h_{R_1}(S_2)\) such that \(s_1(v)=f_2\). This defines the object
\[
 (S_1\times_S S_2,\; f_1\times_S s_2(v))
\]
and the commutative diagram
\[
\begin{tikzcd}
& (S_1\times_S S_2,\; f_1\times s_2(v))
   \arrow[dl,"\pi_1"'] \arrow[dr,"{(\pi_2,v^{-1})}"] &\\
(S_1,f_1)\arrow[dr,"{(h_1,u_1)}"'] &&
(S_2,f_2)\arrow[dl,"{(h_2,u_2)}"]\\
& (S,f) &
\end{tikzcd}
\]
which is a fibre product. This proves homogeneity.

For (b), it remains to show (iii)\(\Rightarrow\)(i). If \(h_{(R_0,R_1,\ldots)}\) is
semi-homogeneous, then \(h_{(R_0,R_1,\ldots)}\) admits a minimally versal formal object.
On the other hand \((R_0,1_{R_0})\) is a quasi-initial minimal object, and the equality on
\(k[\varepsilon]\)-points makes it minimally versal. Hence
\[
 h_{R_0}\longrightarrow h_{(R_0,R_1,\ldots)}
\]
is minimally smooth, which is the required smoothness.

For (c), let \((S_0,\xi_0)\) be a minimally versal formal object of
\(h_{(R_0,R_1,\ldots)}\). For any object \(\xi:R_0\to T\), the functor
\(\Isom(\xi,\xi)\) on \(\mathcal C_\Lambda/(T\times T)\) is represented by
\[
 R_1\otimes_{R_0\otimes_\Lambda R_0}(T\otimes_\Lambda T).
\]
Passing to the inverse limit, \(\Isom(\xi_0,\xi_0)\) is pro-represented by some \(S_1\), and
\(h_{(R_0,R_1,\ldots)}\) is equivalent to \(h_{(S_0,S_1,\ldots)}\). By (b), \(S_\ast\) is
minimally smooth.

\textbf{Definition 2.13.} A cofibered groupoid \(A\) over \(\mathcal C_\Lambda\) is called
\emph{pro-representable} (respectively \emph{minimally pro-representable}) if it is
\(\mathcal C_\Lambda\)-equivalent to
\[
 h_{(R_0,R_1,\ldots)}
\]
for some cogroupoid (respectively normalized cogroupoid)
\((R_0,R_1,\ldots)\) in \(\widehat{\mathcal C}_\Lambda\).

\textbf{Proposition 2.14.} If \((R_0,R_1,\ldots)\) is a normalized cogroupoid in
\(\widehat{\mathcal C}_\Lambda\), then every homomorphism
\[
 (\sigma_0,\sigma_1):(R_0,R_1,\ldots)\longrightarrow(R_0,R_1,\ldots)
\]
which induces a \(\mathcal C_\Lambda\)-equivalence
\[
 h_{(\sigma_0,\sigma_1)}:
 h_{(R_0,R_1,\ldots)}\longrightarrow h_{(R_0,R_1,\ldots)}
\]
is an isomorphism.



% SGA7 batch 101 additions
\providecommand{\Der}{\operatorname{Der}}
\providecommand{\Coker}{\operatorname{Coker}}
\providecommand{\Cov}{\operatorname{Cov}}
\providecommand{\Exalcomm}{\operatorname{Exalcomm}}
\providecommand{\id}{\operatorname{id}}
\providecommand{\cC}{\mathcal C}
\providecommand{\cH}{\mathcal H}
\providecommand{\cP}{\mathcal P}
\providecommand{\cE}{\mathcal E}
\providecommand{\what}{\widehat}
\providecommand{\wt}{\widetilde}
\providecommand{\mhat}{\widehat{\mathfrak m}}
\providecommand{\Rhat}{\widehat R}
\providecommand{\Ahat}{\widehat A}
\providecommand{\Bhat}{\widehat B}
\providecommand{\Ehat}{\widehat E}


\begin{center}
{\Large SGA 7-I — Exposé VI}\\[0.3em]
{\large Source scan pages 71--96: pro-representability and the sheaves \(D^i\)}
\end{center}

\subsection*{2. Pro-representable groupoids, continued}

\textbf{Proof of Proposition 2.14.}  We keep the notation of the preceding statement.  Suppose that a homomorphism \((\sigma_0,\sigma_1)\) of a normalized cogroupoid induces a \(\cC_\Lambda\)-equivalence on the associated cofibered groupoids.  Applying the equivalence to the normalized base object gives a diagram of the form
\[
\begin{tikzcd}[column sep=large,row sep=large]
 B_{R_0}(x,e) \arrow[r] \arrow[d] & h_{R_0}(x,e) \arrow[r] \arrow[d] & h_{(R_0,R_1,\ldots)}(x,e) \arrow[d,"h_{(\sigma_0,\sigma_1)}"] \\
 B_{R_0}(x,e) \arrow[r] & h_{R_0}(x,e) \arrow[r] & h_{(R_0,R_1,\ldots)}(x,e).
\end{tikzcd}
\]
Since, after passage to \(k[\varepsilon]\), the relevant groupoid is totally disconnected, the induced tangent map is bijective.  Hence the endomorphism \(R_0\to R_0\) is a surjective endomorphism of a complete local noetherian ring, and therefore an automorphism.  The defining pull-back formula for \(R_1\) then implies that \(R_1\to R_1\) is an automorphism as well.  Thus \((\sigma_0,\sigma_1)\) is an isomorphism.

\textbf{Lemma 2.15.}  Let \(A,B\) be discrete groupoids over \(\cC_\Lambda\), and let \(\varphi:A\to B\) be a minimally smooth functor.  If \(A\) and \(B\) are homogeneous, then \(\varphi\) is a \(\cC_\Lambda\)-equivalence.

\textit{Proof.}  We must prove that \(\varphi(S):A(S)\to B(S)\) is bijective for every \(S\).  It is surjective by minimal smoothness, and bijective for \(S=k[\varepsilon]\).  We argue by induction on the length of \(S\).  If \(S'\twoheadrightarrow S\) is a small extension, homogeneity gives
\[
 S'\times_S S' \simeq S'\times_k k[\Ker f],
\]
and therefore the commutative diagram
\[
\begin{tikzcd}[column sep=large]
 A(S')\times_{A(S)}A(S') \arrow[r] \arrow[d,"\varphi\times\varphi"'] & A(S')\times A(k[\Ker f]) \arrow[d,"\varphi\times\varphi"] \\
 B(S')\times_{B(S)}B(S') \arrow[r] & B(S')\times B(k[\Ker f]).
\end{tikzcd}
\]
Thus \(A(S')\) and \(B(S')\) are principal homogeneous sets under tangent groups identified by \(\varphi\).  If \(\varphi(S)\) is bijective, then \(\varphi(S')\) is injective and hence bijective. \(\square\)

\textbf{Corollary 2.16.}  Let \(A\) be homogeneous over \(\cC_\Lambda\), with finite-dimensional tangent space.  For every object \(\xi\) of \(A\), the functor
\[
 \Isom_\xi:\cC_\Lambda/R\longrightarrow \mathrm{Ens},\qquad R=P(\xi),
\]
is pro-representable.

\textit{Proof.}  The functor \(\Isom_\xi\) is homogeneous: for a diagram \(T_1\to T\leftarrow T_2\) with one arrow surjective,
\[
 \Isom_\xi(T_1\times_TT_2) \simeq \Isom_\xi(T_1)\times_{\Isom_\xi(T)}\Isom_\xi(T_2).
\]
The formal existence theorem gives a minimally smooth functor from a pro-representable functor to \(\Isom_\xi\); Lemma 2.15 makes it an equivalence. \(\square\)

\textbf{Theorem 2.17.}  A cofibered groupoid \(A\) over \(\cC_\Lambda\) is pro-representable by a normalized smooth cogroupoid if and only if
\begin{enumerate}[label=\textup{(\arabic*)}]
\item \(A\) is homogeneous;
\item \([A(k[\varepsilon])]\) and \(\Aut(1_{k[\varepsilon]})\) are finite-dimensional vector spaces over \(k\).
\end{enumerate}

\textit{Proof.}  Smooth cogroupoids are homogeneous, and the tangent and automorphism spaces are kernels and cokernels of maps between spaces \(\Hom(R_i,k[\varepsilon])\), hence finite-dimensional.  Conversely, choose a minimally versal formal object \(\xi\) represented by \(R_0\).  Corollary 2.16 pro-represents the functor of isomorphisms between the two pull-backs of \(\xi\) over \(R_0\widehat\otimes_\Lambda R_0\) by an object \(R_1\), carrying a universal isomorphism.  Composition of isomorphisms gives a cogroupoid \((R_0,R_1,\ldots)\).  The induced functor \(h_{(R_0,R_1,\ldots)}\to A\) is minimally smooth, hence, by Lemma 2.15, an equivalence. \(\square\)

\textbf{Corollary 2.18.}  A functor \(F:\cC_\Lambda\to\mathrm{Ens}\) is pro-representable if and only if it is homogeneous and \(\dim_k F(k[\varepsilon])<\infty\).

\textbf{Corollary 2.19.}  Every smooth cogroupoid in \(\widehat{\cC}_\Lambda\) can be normalized, up to isomorphism.

For a cogroupoid \((R_0,R_1,\ldots)\), let \(\overline R_1=R_1/J\), where \(J\) is generated by \(s_1(r)-s_2(r)\), \(r\in R_0\).  Then \((R_0,\overline R_1,\ldots)\) defines a formal group scheme over \(R_0\).

\textbf{Corollary 2.20.}  Let \(A\) be homogeneous and pro-represented by a normalized smooth cogroupoid \((R_0,R_1,\ldots)\).  Then \((R_0,\overline R_1,\ldots)\) and \((k,k\widehat\otimes_{R_0}\overline R_1,\ldots)\) pro-represent the automorphism functors \(\Aut_\xi\) and \(\Aut^0\).  Moreover the following are equivalent:
\begin{enumerate}[label=\textup{(\roman*)}]
\item the formal group scheme is formally smooth over \(R_0\);
\item for every surjection \(a'\to a\) in \(A\), the map \(\Aut(a')\to\Aut(a)\) is surjective;
\item \(\overline R_1\) pro-represents the functor of isomorphism classes over \(\cC_\Lambda\).
\end{enumerate}
This is the combination of Proposition 2.7 with Lemma 2.15.

\textbf{Remarks 2.21--2.22.}  Pulling a pro-representation back to \(\cC_\Lambda/k\) replaces the cogroupoid by its base change to \(k\), preserving normalization and relative dimension.  If \(A\) has two smooth pro-representations and one is normalized, then there is a comparison morphism of cogroupoids; the non-normalized base is a formal power-series extension of the normalized one, and the comparison becomes an isomorphism after this smooth base change.

\subsection*{3. The coherent sheaves \(D^i\)}

We now recall the elementary obstruction theory of commutative algebra extensions in the language of Grothendieck cohomology on a site.  Rings are commutative with unit.  For an \(R\)-algebra \(A\), let \(\cC_{A|R}\) be the category of arrows \(B\to A\) in \(R\)-algebras, endowed with the topology whose coverings are surjective arrows.  It has final object \(A\to A\), initial object \(R\to A\), fibre products and coproducts.

A sheaf \(F\) of \(A\)-modules on \(\cC_{A|R}\) is a functor such that, for every cover \(B'\twoheadrightarrow B\),
\[
0\to F(B)\to F(B')\mathrel{\substack{\longrightarrow\\[-.6ex]\longrightarrow}} F(B'\times_B B')
\]
is exact.  Additive sheaves preserve coproducts.  An \(A\)-module \(M\) gives the additive sheaf
\[
 \Der_R(-,M):B\longmapsto \Der_R(B,M),
\]
where \(M\) is viewed as a \(B\)-module through \(B\to A\).

Let \(\widetilde{\cC}_{A|R}\) and \(\widehat{\cC}_{A|R}\) denote sheaves and presheaves.  The inclusion of sheaves into presheaves has right derived functors \(\cH^q_{A|R}\), and
\[
 H^q(A|R,F)=\bigl[\cH^q_{A|R}(F)\bigr](A).
\]
An object \(B\in\cC_{A|R}\) and a ring homomorphism \(S\to R\) give the functorial diagrams
\[
\begin{tikzcd}[column sep=large]
\widehat{\cC}_{B|R} \arrow[r] & \widehat{\cC}_{A|R} \\
\widetilde{\cC}_{B|R} \arrow[u] \arrow[r] & \widetilde{\cC}_{A|R} \arrow[u]
\end{tikzcd}
\qquad
\begin{tikzcd}[column sep=large]
\widehat{\cC}_{A|S} \arrow[r] & \widehat{\cC}_{A|R} \\
\widetilde{\cC}_{A|S} \arrow[u] \arrow[r] & \widetilde{\cC}_{A|R} \arrow[u].
\end{tikzcd}
\]
For a covering \(B'\to B\), the groups \(H^q(\{B'\to B\},F)\) are the cohomology groups of the Čech semi-simplicial module
\[
F(B')\mathrel{\substack{\longrightarrow\\[-.6ex]\longrightarrow}} F(B'\times_BB')\rightrightarrows F(B'\times_BB'\times_BB')\rightrightarrows\cdots .
\]
For a presheaf \(F\), set
\[
\cH^q_{A|R}(F)(B)=\varinjlim_{\{B'\to B\}}H^q(\{B'\to B\},F).
\]
There is the usual spectral sequence
\[
H^p(A|R,\cH^q(F))\Rightarrow H^{p+q}(A|R,F),
\]
and hence, in low degree,
\[
\begin{aligned}
H^1(A|R,F)&=\cH^1(A|R,F),\\
0&\to\cH^2(A|R,F)\to H^2(A|R,F)\to H^1(A|R,\cH^1(F))\\
 &\to\cH^3(A|R,F).
\end{aligned}
\]

\textbf{Lemma 3.2.}  If \(P\twoheadrightarrow B\) is a cover by a polynomial \(R\)-algebra, then
\[
H^n(\{P\to B\},F)\xrightarrow{\sim}\cH^n_{A|R}(F)(B).
\]
Polynomial covers dominate all coverings, which proves the lemma.

\textbf{Corollary 3.3.}  (a) If \(A\) is noetherian and \(F\) is of finite type, then \(\cH^n(F)\) is of finite type.  (b) If \(P\) is polynomial over \(R\), then \(H^n(P|R,F)=0\) for all \(n>0\).  (c) Exactness of sheaves may be checked on polynomial \(R\)-algebras.  Indeed, the polynomial subsite is cofinal and every surjection in it admits a section.

\textbf{Lemma 3.5.}  For a diagram of \(R\)-algebras
\[
\begin{tikzcd}
B' \arrow[r] \arrow[d,two heads] & C \arrow[d] \\
B \arrow[r] & A
\end{tikzcd}
\]
and an \(R\)-algebra \(S\):
\begin{enumerate}[label=\textup{(\roman*)}]
\item if \(\Tor^R_1(B,S)=0\), then
\[
(B'\times_BC)\otimes_RS \simeq (B'\otimes_RS)\times_{B\otimes_RS}(C\otimes_RS);
\]
\item if \(B'\) is \(R\)-flat and \(\Tor_i^R(B,S)=\Tor_i^R(C,S)=0\) for \(0<i\le n\), then \(\Tor_i^R(B'\times_BC,S)=0\) for \(0<i\le n\).
\end{enumerate}
The proof compares the exact rows obtained from \(0\to J\to B'\to B\to0\):
\[
\begin{tikzcd}
0\arrow[r]&J\otimes_RS\arrow[r]\arrow[d,equal]&(B'\times_BC)\otimes_RS\arrow[r]\arrow[d]&C\otimes_RS\arrow[r]\arrow[d,equal]&0\\
0\arrow[r]&J\otimes_RS\arrow[r]&(B'\otimes_RS)\times_{B\otimes_RS}(C\otimes_RS)\arrow[r]&C\otimes_RS\arrow[r]&0.
\end{tikzcd}
\]

\textbf{Corollary 3.6.}  If \(\Tor^R_1(A,S)=0\), then
\[
H^n(A\otimes_RS|S,G)\simeq H^n(A|R,SG).
\]
If \(S\) is \(R\)-flat, this gives \(S\cH^n(G)\simeq\cH^n(SG)\).  Polynomial covers remain Tor-acyclic under the hypotheses of Lemma 3.5.

\textbf{Proposition 3.7.}  (a) If \(\Tor_i^R(A,S)=0\) for \(0<i\le n\), then
\[
\cH^i(A\otimes_RS|S,F)\xrightarrow{\sim}\cH^i(A|R,SF),\quad i\le n.
\]
(b) If \(P\) is polynomial over \(R\), then for an additive sheaf \(F\) on \(\cC_{A\otimes_RP|R}\),
\[
H^n(A\otimes_RP|R,F)\simeq H^n(A|R,F),\quad n>0.
\]
(c) For \(S\to R\to A\) and additive \(F\) on \(\cC_{A|S}\), there is an exact sequence
\[
0\to H^0(A|R,F_R)\to H^0(A|S,F)\to H^0(R|S,F)\to H^1(A|R,F_R)\to\cdots,
\]
where \(F_R(X)=\Ker(F(X)\to F(R))\).  The proof compares the spectral sequences
\[
\begin{tikzcd}[column sep=large]
H^p(A\otimes_RS|S,\cH^q(F)) \arrow[r,Rightarrow] \arrow[d] & H^{p+q}(A\otimes_RS|S,F) \arrow[d] \\
H^p(A|R,\cH^q(SF)) \arrow[r,Rightarrow] & H^{p+q}(A|R,SF)
\end{tikzcd}
\]
and then uses the Leray sequence for the forgetful functor.

\textbf{Corollary 3.8.}  If \(\Tor_i^R(A,B)=0\) for \(0<i\le n\), then for additive \(F\)
\[
H^i(A\otimes_RB|R,F)\simeq H^i(A|R,F)\oplus H^i(B|R,F),\quad i\le n.
\]
This follows from the splitting diagram
\[
\begin{tikzcd}[column sep=large]
H^i(A\otimes_RB|A,F_A) \arrow[r] \arrow[d,"\sim"'] & H^i(A\otimes_RB|R,F) \\
H^i(B|R,F_A) \arrow[r] & H^i(B|R,F).
\end{tikzcd}
\]
If moreover \(A\otimes_RA\simeq A\) and \(\Tor_i^R(A,A)=0\), then \(H^i(A|R,F)=0\) for \(i>0\).

\textbf{Corollary 3.9.}  For a multiplicative subset \(S\subset A\), if \(F\) and \(F_S\) are additive, then
\[
H^i(S^{-1}A|R,F)\simeq H^i(A|R,F).
\]

For an \(A\)-module \(M\), define
\[
D^i(A|R,M)=H^i(A|R,\Der_R(-,M)).
\]
Choose \(P\twoheadrightarrow A\) polynomial over \(R\), with kernel \(I\).

\textbf{Lemma 3.10.}
\[
D^1(A|R,M)=\Coker\bigl(\Der_R(P,M)\to\Hom_A(I/I^2,M)\bigr),
\]
\[
D^2(A|R,M)=H^1(A|R,\cH^1(\Der_R(-,M))).
\]
The proof uses the simplicial algebra \(P^{(n)}=P\times_A\cdots\times_AP\) and the exact sequence
\[
0\to\Der_R(A,M)\to\Der_R(P,M)\to\Hom_A(I/I^2,M)\to H^1(\{P\to A\},\Der_R(-,M))\to0.
\]

Let \(F\twoheadrightarrow I\) be a free \(P\)-module mapping onto \(I\), and let \(K_\bullet\) be the associated Koszul complex.

\textbf{Lemma 3.11.}
\[
D^2(A|R,M)=\Coker\bigl(H^1(K_\bullet,M)\to\Hom_A(H_1(K_\bullet),M)\bigr).
\]
This follows by replacing the cover \(P\to A\) by the polynomial algebra \(P[F]\); the obstruction term is
\[
(F)\cap(F')/((F)(F'))\simeq \Tor^P_1(A,A)=H_1(K_\bullet).
\]

\textbf{Theorem 3.12.}  The groups \(D^i\) satisfy:
\begin{enumerate}[label=\textup{(\arabic*)}]
\item a long exact sequence in \(D^i(A|R,-)\) for every exact sequence of \(A\)-modules;
\item for \(S\to R\to A\), an exact sequence
\[
0\to D^0(A|R,M)\to D^0(A|S,M)\to D^0(R|S,M)\to D^1(A|R,M)\to\cdots;
\]
\item Tor-independent base change \(D^i(A\otimes_RS|S,M)\simeq D^i(A|R,M)\);
\item localization \(D^i(S^{-1}A|R,S^{-1}M)\simeq S^{-1}D^i(A|R,M)\);
\item the product formula \(D^i(A\otimes_RB|R,M)\simeq D^i(A|R,M)\oplus D^i(B|R,M)\) under Tor-independence;
\item compatibility with tensoring by a flat \(A\)-module, when \(A\) is essentially of finite type over a noetherian \(R\);
\item finiteness of \(D^i(A|R,M)\) for \(M\) finite;
\item the explicit formulas
\[
D^0(A|R,M)=\Der_R(A,M),
\]
\[
D^1(A|R,M)=\Coker(\Der_R(P,M)\to\Hom_A(I/I^2,M))=\Exalcomm(A|R,M),
\]
\[
D^2(A|R,M)=\Coker(H^1(K_\bullet,M)\to\Hom_A(H_1(K_\bullet),M)).
\]
\end{enumerate}
The extension interpretation of \(D^1\) is given by the push-out diagram
\[
\begin{tikzcd}
0\arrow[r]&I\arrow[r]\arrow[d]&P\arrow[r]\arrow[d]&A\arrow[r]\arrow[d,equal]&0\\
0\arrow[r]&M\arrow[r]&A'\arrow[r]&A\arrow[r]&0.
\end{tikzcd}
\]
Two maps \(I\to M\) differing by a derivation of \(P\) give isomorphic extensions, and every extension arises this way.

\textbf{Corollary 3.13.}
\begin{enumerate}[label=\textup{(\alph*)}]
\item \(A\) is formally smooth over \(R\) iff \(D^1(A|R,M)=0\) for every \(M\).
\item If \(R\) is noetherian and \(A\) finite type over \(R\), then \(A\) is a relative complete intersection iff \(D^2(A|R,M)=0\) for every \(M\).
\item Localization commutes with \(D^i\).
\item If \(A\) is local, essentially of finite type over noetherian \(R\), and \(\Ahat\) is its completion, then
\[
D^i(\Ahat|R,\Ahat\otimes_AE)\simeq \Ahat\otimes_AD^i(A|R,E)
\]
for every finite \(A\)-module \(E\).
\end{enumerate}
For (d), the final step compares algebra extensions and uses Artin--Rees.  The completion comparison is expressed by the exact commutative diagram
\[
\begin{tikzcd}[column sep=large]
0\arrow[r]&\Ehat\arrow[r]\arrow[d,equal]&B\arrow[r]\arrow[d]&A\arrow[r]\arrow[d]&0\\
0\arrow[r]&\Ehat\arrow[r]&\Bhat\arrow[r]&\Ahat\arrow[r]&0,
\end{tikzcd}
\]
which proves the needed injectivity and surjectivity for \(D^1\), and hence the assertion.



% SGA7 batch 102 additions
\providecommand{\Sing}{\operatorname{Sing}}
\providecommand{\Supp}{\operatorname{Supp}}
\providecommand{\depth}{\operatorname{depth}}
\providecommand{\Fitt}{\operatorname{Fitt}}
\providecommand{\Jac}{\operatorname{Jac}}
\providecommand{\sHom}{\mathcal Hom}
\providecommand{\RHom}{\operatorname{RHom}}
\providecommand{\Aut}{\operatorname{Aut}}
\providecommand{\Der}{\operatorname{Der}}
\providecommand{\Spf}{\operatorname{Spf}}
\providecommand{\cM}{\mathcal M}
\providecommand{\cB}{\mathcal B}
\providecommand{\cJ}{\mathcal J}
\providecommand{\cI}{\mathcal I}
\providecommand{\cL}{\mathcal L}
\providecommand{\cR}{\mathcal R}
\providecommand{\cT}{\mathcal T}
\providecommand{\cU}{\mathcal U}
\providecommand{\T}{\mathcal T}
\providecommand{\cV}{\mathcal V}
\providecommand{\kdual}{k[\varepsilon]}


\begin{center}
{\Large SGA 7-I — Exposé VI}\\[0.3em]
{\large Source scan pages 97--120: coherent deformation sheaves, formal moduli, and the beginning of formal Jacobian subschemes}
\end{center}

\subsection*{3. The coherent sheaves \(D^i\), conclusion}

\textbf{Corollary 3.14.} Let \(R\) be noetherian and let \(A\) be an \(R\)-algebra of finite type, smooth everywhere except at one closed point \(x\in\Spec(A)\). Then there are natural isomorphisms
\[
 D^1(A\mid R,A)\simeq D^1(A_x\mid R,A_x)
 \simeq D^1(\widehat A_x\mid R,\widehat A_x).
\]
Indeed \(D^1(A\mid R,A)\) is an \(A\)-module of finite type whose support is \(\{x\}\). Hence localization and completion at \(x\) do not change it, and the assertion follows from 3.13.

\textbf{Remark 3.15.} Let \(R\) be a topological ring and \(A\) a topological \(R\)-algebra. For an \(A\)-module \(E\) set
\[
 D^i_{\mathrm{top}}(A\mid R,E)=
 \varinjlim_{\alpha} D^i(A_\alpha\mid R_\alpha, A_\alpha\otimes_A E),
\]
where \(A_\alpha=A/J_\alpha\), \(R_\alpha=R/I_\alpha\), and \(J_\alpha\), \(I_\alpha\) range through open ideals of \(A\), respectively \(R\). With this convention, saying that \(A\) is a formally smooth topological \(R\)-algebra means that
\[
 D^1_{\mathrm{top}}(A\mid R,E)=0
\]
for every \(A\)-module \(E\).

Let now \(X\) be a scheme over a ring \(R\), and let \(E\) be a quasi-coherent \(\mathcal O_X\)-module. The sheaves
\[
 D^i_{X\mid R}(E)=D^i(X\mid R,E)
\]
are defined on affine opens by
\[
 \bigl[D^i_{X\mid R}(E)\bigr](U)=
 D^i\bigl(\Gamma(U,\mathcal O_X)\mid R,\Gamma(U,E)\bigr).
\]
If \(R\) is noetherian and \(X\) is of finite type over \(R\), 3.13(a) gives this construction canonically as a quasi-coherent \(\mathcal O_X\)-module. When no ambiguity is possible we write
\[
 D^i_{X\mid R}=D^i_{X\mid R}(\mathcal O_X).
\]
The preceding affine statements globalize to these sheaves. The following are the properties needed below.

\textbf{Corollary 3.15.} Let \(R\) be noetherian and \(X\) an \(R\)-scheme of finite type.
\begin{enumerate}[label=(\alph*)]
\item For every coherent \(\mathcal O_X\)-module \(E\), the sheaf \(D^i_{X\mid R}(E)\) is coherent for all \(i\), and
\[
 \Supp D^i_{X\mid R}(E)\subset X^{\mathrm{Sing}}
 \quad (i>0),
\]
where \(X^{\mathrm{Sing}}\) is the set of points at which \(X\) is not smooth over \(R\).
\item If \(X\) is flat over \(R\), then for every \(R\)-algebra \(S\) there is the canonical base-change isomorphism
\[
 D^i_{X_S\mid S}(E_S)\simeq D^i_{X\mid R}(E)_S
\]
for all \(i\).
\item If \(X\) is a relative complete intersection over \(R\), then, for every surjection of rings \(S\twoheadrightarrow R\), there is an exact sequence
\[
 0\to D^1_{X\mid R}(E)\to D^1_{X\mid S}(E)
 \to \mathcal Hom_{\mathcal O_X}(I/I^2\otimes_R\mathcal O_X,E)\to 0,
\]
where \(I=\Ker(S\to R)\). If, moreover, \(R\) is local with residue field \(k\) and \(X\) is \(R\)-flat, then
\[
 0\to D^1_{X_0\mid k}(\mathcal O_{X_0})\to
 D^1_{X\mid S}(\mathcal O_X)\otimes_R k
 \to D^1(R\mid S,k)\otimes_k\mathcal O_{X_0}\to 0,
\]
where \(X_0=X\otimes_R k\).
\end{enumerate}
The proof is exactly the passage from the affine statements 3.12 and 3.13 to sheaves. In (c), relative complete-intersection hypotheses give \(D^i_{X\mid R}(E)=0\) for \(i\geq 2\), and the exact sequence of 3.12(2), applied to \(S\to R\), becomes the displayed exact sequence after identifying
\[
 D^1_{R\mid S}(E)=\mathcal Hom_{\mathcal O_X}(I/I^2\otimes_R\mathcal O_X,E).
\]
The final sequence follows by reducing modulo the maximal ideal of \(R\) and using the base-change isomorphism.

\textbf{Remark 3.16.} The quasi-coherent sheaves \(D^i_{X\mid R}(E)\) should occur as the \(E^{0,i}_2\)-terms of a spectral sequence relating the local theory to the global theory. L. Illusie has obtained such a theory by homotopical algebra, generalizing methods of M. Andre and D. Quillen rather than only cohomological algebra. The present approach should also globalize by using Grothendieck cohomology on the category of schemes with a suitable topology extending the affine one.

\section*{4. Formal moduli of deformations}

The cofibered categories of deformation problems are among the useful examples in which isomorphism classes of objects are often not pro-representable, but nevertheless admit formal versal objects. We return to the notation of paragraph 1. Thus \(\Lambda\) is a fixed complete local noetherian ring with residue field \(k\), and \(\mathcal C_\Lambda\) is the category of artinian local \(\Lambda\)-algebras with residue field \(k\).

Let \(X_0\) be a fixed scheme over \(k\). An infinitesimal deformation of \(X_0\) is a pair \((R,X)\), where \(R\in\operatorname{ob}(\mathcal C_\Lambda)\), \(X\) is a flat \(R\)-scheme, and there is a cartesian diagram
\[
\begin{tikzcd}[column sep=large,row sep=large]
 X_0 \arrow[r,"i_X"] \arrow[d] & X \arrow[d,"\mathrm{flat}"]\\
 \Spec k \arrow[r] & \Spec R .
\end{tikzcd}
\]
The induced morphism \(X_0\to X\times_{\Spec R}\Spec k\) is required to be an isomorphism. Since \(R\) is artinian local, \(i_X\) is a closed immersion and a homeomorphism on underlying spaces.

For two deformations \((R,X)\) and \((R',X')\), an arrow \((R',X')\to(R,X)\) consists of a morphism \(\varphi:R'\to R\) in \(\mathcal C_\Lambda\) and a morphism \(\xi:X\to X'\) fitting into the diagram
\[
\begin{tikzcd}[column sep=large,row sep=large]
 X \arrow[rr,"\xi"] \arrow[d] & X_0 \arrow[l,"i_X"'] \arrow[r,"i_{X'}"] \arrow[d] & X' \arrow[d]\\
 \Spec R & \Spec k \arrow[l] \arrow[r] & \Spec R'.
\end{tikzcd}
\]
This gives a category \(\mathcal M_{X_0}\) and a functor
\[
 p:\mathcal M_{X_0}\to\mathcal C_\Lambda,
 \qquad (R,X)\mapsto R.
\]
If \((R,X)\) is an object and \(R\to S\) is a morphism in \(\mathcal C_\Lambda\), then
\[
 X_S=X\times_{\Spec R}\Spec S
\]
is a deformation of \(X_0\) over \(S\), and the cartesian square
\[
\begin{tikzcd}[column sep=large,row sep=large]
 X_S \arrow[r,"p_1"] \arrow[d,"\mathrm{flat}"'] & X \arrow[d,"\mathrm{flat}"]\\
 \Spec S \arrow[r,"\Spec(\varphi)"'] & \Spec R
\end{tikzcd}
\]
shows that \(\mathcal M_{X_0}\) is cofibered over \(\mathcal C_\Lambda\).

\textbf{4.1.} The following elementary properties will be used repeatedly.
\begin{enumerate}[label=(\alph*)]
\item Every arrow over the identity of \(R\) is an isomorphism. Thus \(\mathcal M_{X_0}\) is cofibered in groupoids.
\item If \(X_0=\Spec A_0\) is affine, then every deformation \((R,X)\) is affine. The fiber \(\mathcal M_{X_0}(R)\) is equivalent to the category whose objects are flat \(R\)-algebras \(A\), together with an \(R\)-algebra map \(j:A\to A_0\) inducing \(A\otimes_R k\simeq A_0\).
\item The corresponding formal category \(\widehat{\mathcal M}_{X_0}\) has, over \(R\), formal schemes \(\mathfrak X\) adic over \(\Spf R\), flat modulo every power of the maximal ideal. If \(X_0\) is noetherian, this is the expected category of formal deformations of \(X_0\).
\item If \(U_0\subset X_0\) is open, restriction defines \(\mathcal M_{X_0}\to\mathcal M_{U_0}\). Likewise a point \(x_0\in X_0\) gives functors to the deformation categories of \(\Spec \mathcal O_{X_0,x_0}\) and \(\Spec \widehat{\mathcal O}_{X_0,x_0}\).
\end{enumerate}

\textbf{Lemma 4.2.} Consider a diagram of rings
\[
\begin{tikzcd}
 R' \arrow[r] & R\\
 R'' \arrow[u] \arrow[r,two heads] & R_0 \arrow[u]
\end{tikzcd}
\]
in which the indicated map is surjective. If \(E\) and \(E'\) are flat modules over the two upper rings and are identified after base change to \(R_0\), then the fiber product module is flat over the fiber product ring. Equivalently, if
\[
 0\to I\to R\to R'\to 0
\]
is exact, the sequence
\[
 0\to I\otimes E\to E\times E'\to E'\to 0
\]
is exact and gives the Tor-vanishing needed for flatness.

\textbf{Corollary 4.3.} For a diagram
\[
 (R,X)\longleftarrow(R_0,X_0)\longrightarrow(R',X')
\]
in \(\mathcal M_{X_0}\), with the ring map \(R'\to R_0\) surjective, the amalgamated sum \(X\amalg_{X_0}X'\) is flat over \(R\times_{R_0}R'\). Hence \(\mathcal M_{X_0}\) is a homogeneous cofibered category in groupoids over \(\mathcal C_\Lambda\).

\textbf{Lemma 4.4.} There is a canonical exact sequence
\[
 0\to H^1(X_0,D^0_{X_0})\to [\mathcal M_{X_0}(k[\varepsilon])]
 \to H^0(X_0,D^1_{X_0}).
\]
If \(X_0=\Spec A_0\), the middle set is identified with isomorphism classes of \(k\)-algebra extensions of \(A_0\) by \(A_0\), hence with \(D^1(A_0\mid k,A_0)\). In general the map is obtained by restricting to affine opens. Its kernel is the set of locally trivial deformations, and the sheaf of germs of automorphisms of the trivial deformation \(X_0\times\Spec k[\varepsilon]\) is \(D^0_{X_0}\); the usual Cech argument gives the displayed exact sequence.

\textbf{Theorem 4.5.} Let \(X_0\) be a scheme of finite type over \(k\). Suppose either that \(X_0\) is proper over \(k\), or that \(X_0\) is affine and
\[
 X_0^{\mathrm{Sing}}=\{x\in X_0\mid X_0\text{ is not smooth over }k\text{ at }x\}
\]
is finite. Then:
\begin{enumerate}[label=(\alph*)]
\item \(\mathcal M_{X_0}\) admits a formal versal object over \(\mathcal C_\Lambda\). If \((R,\mathfrak X)\) is such a formal versal deformation and \((S,Y)\) is an infinitesimal deformation, there is a morphism \(R\to S\) and a cartesian diagram
\[
\begin{tikzcd}
 Y \arrow[r] \arrow[d] & \mathfrak X \arrow[d]\\
 \Spf S \arrow[r] & \Spf R.
\end{tikzcd}
\]
Moreover \((R,\mathfrak X)\) pro-represents the functor of isomorphism classes if and only if, for every surjection \(R'\to R\) in \(\mathcal C_\Lambda\) and every deformation \((R',X')\), the homomorphism
\[
 \Aut_{R'}(X'\mid X_0)\to \Aut_R(X'\otimes_{R'}R\mid X_0)
\]
is surjective.
\item If \(X_0\) is proper over \(k\), then \(\mathcal M_{X_0}\) is smoothly pseudo-representable over \(\mathcal C_\Lambda\).
\end{enumerate}
The proof uses homogeneity from 4.3 and the tangent exact sequence 4.4. In the proper case, the coherent sheaves \(D^i\) have finite-dimensional cohomology. In the affine isolated-singularity case, \(H^1(X_0,D^0_{X_0})=0\) and \(H^0(X_0,D^1_{X_0})\) is finite-dimensional because the support is finite. The criterion for pro-representability is the automorphism criterion from 2.7.

\textbf{Remarks 4.6.} (a) If \(X_0\) is affine, a formal versal deformation exists exactly when \(D^1_{X_0}\) has finite support; this need not mean that the non-smooth locus itself is finite. (b) From now on, ``deformation of \(X_0\)'' means an object of \(\mathcal M_{X_0}\); when ambiguity is possible, an object in the fiber over an artinian ring is called an infinitesimal deformation. (c) Although \(\Lambda\) is fixed, the number \(\dim_k \mathfrak m_R/\mathfrak m_R^2\) for a formal versal deformation depends only on \(X_0\), not on the choice of \(\Lambda\).

The same method applies to a pair \(Y_0\subset X_0\), where \(Y_0\) is a closed subscheme. A deformation of the pair is a diagram
\[
\begin{tikzcd}[column sep=large,row sep=large]
 Y_0 \arrow[r] \arrow[d] & Y \arrow[d,"\mathrm{flat}"]\\
 X_0 \arrow[r] \arrow[d] & X \arrow[d,"\mathrm{flat}"]\\
 \Spec k \arrow[r] & \Spec R,
\end{tikzcd}
\]
which becomes the given pair after base change to \(k\). These deformations form a homogeneous cofibered category \(\mathcal M_{X_0,Y_0}\) over \(\mathcal C_\Lambda\).

\textbf{Lemma 4.7.} Let \(I\) be the ideal sheaf of \(Y_0\) in \(X_0\). There is an exact sequence
\[
0\to H^0\bigl(Y_0,\mathcal Hom_{\mathcal O_{Y_0}}(I/I^2,\mathcal O_{Y_0})\bigr)
\to [\mathcal M_{X_0,Y_0}(k[\varepsilon])]
\to [\mathcal M_{X_0}(k[\varepsilon])].
\]
The proof identifies the kernel with deformations of \(Y_0\) inside the trivial first-order deformation \(X_0\times\Spec k[\varepsilon]\). Such an object is equivalently a diagram
\[
\begin{tikzcd}[column sep=large,row sep=large]
 0 \arrow[r] & I \arrow[r] \arrow[d] & \mathcal O_{X_0} \arrow[r] \arrow[d] & \mathcal O_{Y_0} \arrow[r] \arrow[d,equal] & 0\\
 0 \arrow[r] & \mathcal O_{Y_0} \arrow[r,"\varepsilon"] & \mathcal O_Y \arrow[r] & \mathcal O_{Y_0} \arrow[r] & 0,
\end{tikzcd}
\]
and hence is determined by a morphism \(I/I^2\to\mathcal O_{Y_0}\).

\textbf{Remark 4.8.} A more general formula, implying Lemma 4.7, is given in SGA 5, III, 4.4. One can also define, for any \(S\)-scheme \(X\), quasi-coherent sheaves \(D^i_{X\mid S}\), extending the affine case, and one has a spectral sequence relating the global Ext-groups for the cotangent complex to the cohomology of these sheaves. In particular the low-degree terms recover the tangent, indeterminacy, and obstruction groups used in the preceding deformation arguments.

\textbf{Theorem 4.9.} Let \(X_0\) be a \(k\)-scheme of finite type and let \(Y_0\subset X_0\) be a closed subscheme. Assume either that \(X_0\) is proper over \(k\), or that \(X_0\) is affine, smooth outside a finite set of points, and \(Y_0\) is proper over \(k\). Then \(\mathcal M_{X_0,Y_0}\) admits a formal versal deformation over \(\mathcal C_\Lambda\), with the same pro-representability criterion by surjectivity on automorphism groups. If, moreover,
\[
 H^0\bigl(Y_0,\mathcal Hom(I/I^2,\mathcal O_{Y_0})\bigr)
\]
is finite-dimensional over \(k\), then \(\mathcal M_{X_0,Y_0}\) is smoothly pseudo-representable. This follows at once from 4.4, 4.6, 4.7 and the criteria 1.11, 2.7, and 2.13.


We now consider passage to localization and completion in the deformation theory of schemes.  If \(x_0\in X_0\) is fixed, one has natural functors
\[
 \mathcal M_{X_0}\longrightarrow
 \mathcal M_{\Spec(\mathcal O_{X_0,x_0})}
 \longrightarrow
 \mathcal M_{\Spec(\widehat{\mathcal O}_{X_0,x_0})},
\]
and after completing formal deformations one obtains
\[
 \widehat{\mathcal M}_{X_0}\longrightarrow
 \widehat{\mathcal M}_{\Spec(\mathcal O_{X_0,x_0})}
 \longrightarrow
 \widehat{\mathcal M}_{\Spec(\widehat{\mathcal O}_{X_0,x_0})}.
\]
If \((R,\mathfrak X)\) is a formal deformation and
\[
 \mathfrak X=\varprojlim_\nu X_\nu,\qquad
 X_\nu=\mathfrak X\times_{\Spf R}\Spec(R/\mathfrak m^{\nu+1}),
\]
then the two images at \(x_0\) are
\[
 \mathfrak X_{(x_0)}=\varprojlim_\nu\Spec(\mathcal O_{X_\nu,x_0}),\qquad
 \widehat{\mathfrak X}_{(x_0)}=\varprojlim_\nu\Spec(\widehat{\mathcal O}_{X_\nu,x_0}).
\]

\textbf{Theorem 4.10.} Let \(X_0=\Spec(A_0)\) be an affine scheme of finite type over \(k\).
\begin{enumerate}[label=(\alph*)]
\item If \(X_0\) is locally a complete intersection, then \(\mathcal M_{X_0}\) is smooth over \(\mathcal C_\Lambda\).
\item Assume that, for a fixed closed point \(x\in X_0\), the complement \(X_0-\{x\}\) is smooth over \(k\). Then the functors
\[
 \mathcal M_{X_0}\longrightarrow \mathcal M_{\Spec(\mathcal O_{X_0,x})}
 \longrightarrow \mathcal M_{\Spec(\widehat{\mathcal O}_{X_0,x})}
\]
are both minimally smooth. In particular a formal versal deformation of \(X_0\) induces formal versal deformations of \(\Spec(\mathcal O_{X_0,x})\) and \(\Spec(\widehat{\mathcal O}_{X_0,x})\).
\end{enumerate}
\textit{Proof.}  For (a), it is enough to prove smoothness of the functor \([\mathcal M_{X_0}]\to\mathcal C_\Lambda\). Let \((R,\Spec A)\) be a deformation of \(\Spec A_0\), and let \(R'\to R\) be a small extension. The exact sequence for the composite \(R'\to R\to A\), with coefficients in \(A_0\), gives
\[
 \cdots\to D^1(A\mid R,A_0)\to D^1(A\mid R',A_0)\to
 D^1(R\mid R',k)\otimes_k A_0\to D^2(A\mid R,A_0)\to\cdots .
\]
Because \(A\) is \(R\)-flat,
\[
 D^i(A\mid R,A_0)=D^i(A_0\mid k,A_0)\qquad (i>0),
\]
and the last group vanishes for \(i=2\) when \(A_0\) is locally a complete intersection. Thus one finds an \(R'\)-algebra extension \(A'\) of \(A\) by \(A_0\) mapping to the class of \(R'\), i.e. a lifting \((R',\Spec A')\) of \((R,\Spec A)\).

For (b), Corollary 3.14 gives canonical identifications of tangent spaces
\[
 D^1(A_0\mid k,A_0)
i \simeq D^1((A_0)_x\mid k,(A_0)_x)
 \simeq D^1(\widehat{(A_0)_x}\mid k,\widehat{(A_0)_x}).
\]
Hence
\[
 [\mathcal M_{X_0}(k[\varepsilon])]
i\simeq
 [\mathcal M_{\Spec(\mathcal O_{X_0,x})}(k[\varepsilon])]
i\simeq
 [\mathcal M_{\Spec(\widehat{\mathcal O}_{X_0,x})}(k[\varepsilon])].
\]
To prove minimal smoothness it remains, by 2.7(a), to show smoothness of the functors. Given \((R,\\Spec A)\) and a small extension \(R'\to R\), the same exact commutative diagram for \(R'\to R\to A\), after localization and completion at \(x\), shows that any local or completed local lifting may be lifted globally, because the groups \(D^i(A\mid R,A_0)\) agree with their local and completed counterparts in positive degree.

\textbf{Theorem 4.11.} Let \(X_0\) be a separated \(k\)-scheme of finite type, smooth outside the finite closed set \(\{x_1,\ldots,x_r\}\). Let \(U_0^1,\ldots,U_0^r\) be affine open neighbourhoods of these points such that \(U_0^i\cap U_0^j\) is smooth over \(k\) for \(i\neq j\). If
\[
 H^2(X_0,D^0_{X_0})=0,
\]
then the natural functor
\[
 \mathcal M_{X_0}\to \mathcal M_{U_0^1}\times\cdots\times\mathcal M_{U_0^r}
\]
is smooth.
\textit{Proof.}  Complete \(U_0^1,\ldots,U_0^r\) to an affine open covering by choosing smooth affine opens \(U_0^j\), \(j>r\), such that \(U_0^i\cap U_0^j\) is affine for all \(i,j\). Let \(R'\to R\) be a small extension. Suppose that a deformation \((R,X)\) of \(X_0\) and liftings \((R',U'^i)\) of the prescribed deformations over \(U_0^i\), \(1\leq i\leq r\), are given. For \(j>r\), since \(U_0^j\) is smooth, 4.9(a) lifts \(X|_{U_0^j}\) to a deformation \((R',U'^j)\).

For each intersection \(U_0^i\cap U_0^j\), smoothness gives isomorphisms
\[
 \tau_{ij}: U'^i|_{U_0^i\cap U_0^j}\xrightarrow{\sim}
 U'^j|_{U_0^i\cap U_0^j}
\]
which reduce to the identity over \(R\). The defect \(\tau_{ij}\tau_{jk}\tau_{ik}^{-1}\) is a Cech 2-cocycle with values in \(D^0_{X_0}\). Its cohomology class lies in \(H^2(X_0,D^0_{X_0})\); this class vanishes by hypothesis. Therefore the local liftings can be modified so that the cocycle condition holds, and then they glue to a deformation \((R',X')\) with \(X'\otimes_{R'}R=X\).

\textbf{Remark 4.12.} The preceding proof unnecessarily assumes that \(X_0\) is separated. In fact, for a fixed deformation \((R,X)\) and compatible deformations \((R',U'^j)\) on an open covering, the deformations \((R',V')\) of any open subset \(V\subset X\), compatible with the already given data, form a gerbe in the sense of Giraud. Its band is \(D^0_{X_0}\). Hence the obstruction to a section of this gerbe, i.e. to a global deformation compatible with the local data, is again in \(H^2(X_0,D^0_{X_0})\).

\textbf{Corollary 4.13.} Let \(X_0\) be a \(k\)-scheme of finite type, smooth outside a finite set
\[
 X_0^{\mathrm{Sing}}=\{x_1,\ldots,x_r\},
\]
and assume \(H^2(X_0,D^0_{X_0})=0\). Then
\[
 \mathcal M_{X_0}\to
 \mathcal M_{\Spec(\widehat{\mathcal O}_{X_0,x_1})}\times\cdots\times
 \mathcal M_{\Spec(\widehat{\mathcal O}_{X_0,x_r})}
\]
is smooth. In particular, if \((R,\mathfrak X)\) is a formal versal deformation of \(X_0\), and if \((R_i,\mathfrak X_i)\) are formal versal deformations of the completed local schemes, then
\[
 \dim \mathcal M_{X_0}
 =\dim H^1(X_0,D^0_{X_0})+
 \sum_{i=1}^r \dim \mathcal M_{\Spec(\widehat{\mathcal O}_{X_0,x_i})}.
\]
If moreover \(X_0\) is locally a complete intersection, then \(\mathcal M_{X_0}\) is smooth.
\textit{Proof.} The first assertion is obtained from 4.10 and 4.9(b). Choose formal versal deformation rings \(R\) and \(R_i\). The diagram
\[
\begin{tikzcd}[column sep=large,row sep=large]
 h_R \arrow[r] \arrow[d] & h_{R_1\widehat\otimes\cdots\widehat\otimes R_r} \arrow[d]\\
 \mathcal M_{X_0} \arrow[r] & \prod_i\mathcal M_{\Spec(\widehat{\mathcal O}_{X_0,x_i})}
\end{tikzcd}
\]
shows, by 2.7(a), that \(h_R\to h_{R_1\widehat\otimes\cdots\widehat\otimes R_r}\) is smooth after allowing \(d=\dim_kH^1(X_0,D^0_{X_0})\) formal parameters. Taking tangent spaces gives the displayed dimension formula.

\textbf{4.14. Explicit complete-intersection construction.} Suppose
\[
 A_0=k[x_1,\ldots,x_n]/(f_1,\ldots,f_r),
\]
where \(f_1,\ldots,f_r\) is a regular sequence. The exact sequence
\[
 \cdots\to D^1(k[x_1,\ldots,x_n]\mid k,A_0)\to
 \Hom_{A_0}(J/J^2,A_0)\to D^1(A_0\mid k,A_0)\to0
\]
allows one to choose maps \(P_{
u j}\) representing a \(k\)-basis of \(D^1(A_0\mid k,A_0)\). Put
\[
 R=\Lambda[[t_1,\ldots,t_d]],\qquad
 B=R[[x_1,\ldots,x_n]]/(G_1,\ldots,G_r),
\]
where \(G_j=F_j-\sum_\nu t_\nu P_{
u j}\), and where \(F_j\) is obtained from \(f_j\) by lifting its coefficients from \(k\) to \(\Lambda\). Since the sequence remains regular modulo the maximal ideal, \(B\) is flat over \(R\), and \(k\otimes_RB=A_0\). The induced map
\[
 \Hom_{\Lambda\text{-alg}}(R,k[\varepsilon])\to D^1(A_0\mid k,A_0)
\]
is bijective by construction. Thus \((R,\Spf B)\) is a formal versal object. The same formula applies to the completed local complete-intersection algebra by 4.10(b).

\textbf{Remarks 4.15--4.16.} If \(X_0\) is affine and locally a complete intersection with only isolated singularities, the construction shows that a formal versal deformation is the formalization of an \(R\)-scheme of finite type. If \(X_0\) is not necessarily affine but \(H^2(X_0,D^0_{X_0})=0\), the same statement globalizes.

The global theory of the relative cotangent complex gives a uniform formulation. Let \(S_0\subset S'\subset S\) be defined by ideals \(J\supset I\) with \(JI=0\), let \(X'\) be flat over \(S'\), and put \(L=L_{X_0/S_0}\). For the category \(F\) of flat schemes over \(S\) extending \(X'\), one has:
\begin{enumerate}[label=(\alph*)]
\item the automorphism group of an object is \(\Ext^0(X_0,L,I_{X_0})\);
\item the set of isomorphism classes of objects, if non-empty, is a torsor under \(\Ext^1(X_0,L,I_{X_0})\);
\item there is a canonical obstruction class in \(\Ext^2(X_0,L,I_{X_0})\), whose vanishing is necessary and sufficient for non-emptiness.
\end{enumerate}
Using
\[
 \Ext^i(X_0,L,I_{X_0})=H^i(X_0,\RHom(L,I_{X_0})),
\]
and the fact that the cohomology sheaves of \(\RHom(L,I_{X_0})\) are the \(D^q_{X_0}(I_{X_0})\), one obtains the spectral sequence
\[
 E^{p,q}_2=H^p(X_0,D^q_{X_0}(I_{X_0}))
 \Longrightarrow \Ext^{p+q}(X_0,L,I_{X_0}).
\]
The first obstruction lies in \(H^0(X_0,D^2)\); when it vanishes, the local solutions form a torsor under \(D^1\), with obstruction in \(H^1(X_0,D^1)\); after choosing a coherent system of local isomorphism classes, the remaining obstruction lies in a quotient of \(H^2(X_0,D^0)\). This is the intrinsic form of the gluing argument used above.


\section*{5. Formal Jacobian subschemes}

Let \(X_0\) be a scheme of finite type over a field \(k\), and let \((R,\mathfrak X)\) be a formal versal deformation of \(X_0\). Since \(\mathfrak X\) is an \(R\)-adic formal scheme rather than an ordinary scheme over \(R\), we must say what is meant by its non-smooth locus.

We first recall Fitting ideals. Let \(A\) be a commutative ring and \(E\) an \(A\)-module of finite type. Choose a presentation
\[
 K\to F\to E\to 0,
\]
where \(F\) is locally free of constant rank \(n\). Define
\[
 I^p(E)=\operatorname{Im}\bigl(\bigwedge^{n-p}K\otimes\bigwedge^{p}F^\vee\to A\bigr)
 \qquad (p=0,1,2,\ldots).
\]
These ideals are independent of the chosen presentation and are the Fitting ideals of \(E\). They satisfy:
\begin{enumerate}[label=(\alph*)]
\item For every ring map \(A\to B\), base change gives
\[
 B\otimes_A A/I^p(E)\simeq B/I^p(B\otimes_AE).
\]
In particular localization commutes with \(I^p\).
\item The ideals are increasing, and \(I^p(E)_x=A_x\) whenever \(p>\dim_{k(x)} E(x)\). The support of \(A/I^0(E)\) is \(\Supp E\).
\item Direct sums satisfy the usual convolution formula for Fitting ideals.
\item If \(A\) is noetherian adic, \(A=\varprojlim A_n\), and \(E=\varprojlim E_n\), then
\[
 I^p(E)=\varprojlim I^p(E_n).
\]
\end{enumerate}
Thus for a scheme \(X\) and a quasi-coherent module \(E\) of finite presentation one obtains quasi-coherent ideals \(I^p_X(E)\) by the affine formula
\[
 I^p_X(E)(U)=I^p\bigl(\Gamma(U,E)\bigr)
\]
on affine opens \(U\). If \(X\to Y\) is essentially of finite presentation, the relative Kahler differentials \(\Omega^1_{X/Y}\) are of finite presentation, so the ideals \(I^p_X(\Omega^1_{X/Y})\) define closed subschemes. This is the starting point for the formal Jacobian subschemes.



% SGA7 pages 121--137 macro additions
\providecommand{\Der}{\operatorname{Der}}
\providecommand{\Supp}{\operatorname{Supp}}
\providecommand{\Ann}{\operatorname{Ann}}
\providecommand{\Fitt}{\operatorname{Fitt}}
\providecommand{\Jac}{\operatorname{Jac}}
\providecommand{\Frac}{\operatorname{Frac}}
\providecommand{\Sym}{\operatorname{Sym}}
\providecommand{\codim}{\operatorname{codim}}
\providecommand{\cJ}{\mathcal J}
\providecommand{\cI}{\mathcal I}
\providecommand{\cK}{\mathcal K}
\providecommand{\cM}{\mathcal M}
\providecommand{\cN}{\mathcal N}
\providecommand{\cX}{\mathfrak X}
\providecommand{\cY}{\mathfrak Y}
\providecommand{\cS}{\mathfrak S}
\providecommand{\cZ}{\mathfrak Z}
\providecommand{\Oo}{\mathcal O}
\providecommand{\widebar}{\overline}


\section*{5. Formal Jacobian subschemes, continued}

We use the notation introduced above for the Fitting ideals \(I^p(E)\).  For a morphism
\[
        f:X\longrightarrow Y
\]
which is essentially of finite presentation, we denote by
\[
        I^p(X\mid Y)=I^p_X(\Omega^1_{X/Y})
\]
the corresponding ideal sheaf on \(X\), and by \(J^p(X\mid Y)\) the closed subscheme which it defines.  Thus \(J^0(X\mid Y)\supset J^1(X\mid Y)\supset\cdots\).

\textbf{Theorem 5.3.}  Let \(f:X\to Y\) be essentially of finite presentation.  Then:
\begin{enumerate}[label=(\arabic*)]
\item For every base change \(Y'\to Y\), with \(X'=X\times_Y Y'\), the canonical morphism
\[
        J^p(X'\mid Y')\longrightarrow J^p(X\mid Y)\times_Y Y'
\]
is an isomorphism, for every \(p\).
\item For every \(x\in X\), \(y=f(x)\), the canonical local comparison is an isomorphism:
\[
        I^p_X(\Omega^1_{X/Y})_x
        \simeq
        I^p\bigl(\Omega^1_{\Oo_{X,x}/\Oo_{Y,y}}\bigr).
\]
\item The \(J^p(X\mid Y)\) form a decreasing sequence of closed subschemes.  A point \(x\) does not belong to \(J^p(X\mid Y)\) once \(p\) is larger than
\[
        \dim_{\kappa(x)}\bigl(\Omega^1_{X/Y}\otimes\kappa(x)\bigr).
\]
\item One has \(x\notin J^p(X\mid Y)\) if and only if, after replacing \(X\) by an open neighbourhood \(U\) of \(x\), there exists a closed immersion
\[
        i:U\hookrightarrow U'
\]
into a \(Y\)-scheme \(U'\) which is smooth over \(Y\) at \(x'=i(x)\), and whose relative embedding dimension at \(x'\) is \(p\).
\item If \(X_1\) and \(X_2\) are two \(Y\)-schemes essentially of finite presentation and \(X=X_1\times_YX_2\), then
\[
        I^p(X\mid Y)
        =
        \sum_{p_1+p_2=p}
        \pr_1^*I^{p_1}(X_1\mid Y)\,
        \pr_2^*I^{p_2}(X_2\mid Y).
\]
\end{enumerate}

\emph{Proof.}  Assertion (3) is the elementary monotonicity of Fitting ideals.  Assertion (1) follows from the base-change formula for Fitting ideals and from
\[
        \Omega^1_{X'/Y'}\simeq \pr^*\Omega^1_{X/Y}.
\]
Assertion (2) is the same statement after localization at \(x\).  Assertion (5) follows from the direct-sum formula for Fitting ideals applied to
\[
        \Omega^1_{X_1\times_YX_2/Y}
        \simeq
        \pr_1^*\Omega^1_{X_1/Y}\oplus
        \pr_2^*\Omega^1_{X_2/Y}.
\]

For (4), suppose first that such an immersion \(U\hookrightarrow U'\) exists.  Since \(U'\) is smooth over \(Y\) at \(x'\), \(\Omega^1_{U'/Y,x'}\) is free of rank \(p\), and the exact conormal sequence
\[
        I/I^2\longrightarrow
        \Omega^1_{U'/Y}\otimes\Oo_U\longrightarrow
        \Omega^1_{U/Y}\longrightarrow 0
\]
shows that the \(p\)-th Fitting ideal is the unit ideal at \(x\).  Hence \(x\notin J^p(U\mid Y)\), and therefore \(x\notin J^p(X\mid Y)\).

Conversely, the question is local.  Write \(X=\Spec B\), \(Y=\Spec A\), and choose a presentation
\[
        0\longrightarrow J\longrightarrow P\longrightarrow B\longrightarrow 0,
        \qquad
        P=A[x_1,\ldots,x_n].
\]
The exact sequence
\[
        J/J^2\longrightarrow B\otimes_P\Omega^1_{P/A}
        \longrightarrow \Omega^1_{B/A}\longrightarrow 0
\]
has middle term free of rank \(n\).  The condition \(x\notin J^p(X\mid Y)\) means that a suitable \((n-p)\)-minor of the Jacobian matrix of chosen generators of \(J\) is invertible at \(x\).  Thus there are elements
\[
        f_1,\ldots,f_{n-p}\in J
\]
such that
\[
        U'=\Spec\bigl(P/(f_1,\ldots,f_{n-p})\bigr)
\]
is smooth over \(Y\) at the image \(x'\) of \(x\), and \(U=\Spec B\) embeds as a closed subscheme of \(U'\).  This gives the required local presentation.

The following result is the corresponding Jacobian criterion.

\textbf{Theorem 5.4.}  Let \(f:X\to Y\) be essentially of finite presentation.  Suppose that there is an integer \(n\geq 0\) such that either:
\begin{enumerate}[label=(\arabic*)]
\item \(f\) is flat and all fibers have dimension \(n\);
\item \(f\) is a relative complete intersection of virtual dimension \(n\) in the sense of SGA 6, VIII 1.9.
\end{enumerate}
Then
\[
        x\notin J^n(X\mid Y)
        \quad\Longleftrightarrow\quad
        f\text{ is smooth at }x.
\]

\emph{Proof.}  By Theorem 5.3(4), \(x\notin J^n(X\mid Y)\) if and only if locally \(X\) embeds into a \(Y\)-scheme \(U'\), smooth over \(Y\) at \(x'\), with relative dimension \(n\).  It remains to see that such an embedding is an isomorphism at \(x\).

Let
\[
        0\longrightarrow I\longrightarrow B'\longrightarrow B\longrightarrow 0
\]
be the corresponding exact sequence of local rings, with \(B'\) the local ring of \(U'\) at \(x'\) and \(B\) the local ring of \(X\) at \(x\).  Put \(A=\Oo_{Y,y}\).  In case (1), flatness gives the exact sequence on the closed fiber
\[
        0\longrightarrow I/\mathfrak m I
        \longrightarrow B'/\mathfrak mB'
        \longrightarrow B/\mathfrak mB
        \longrightarrow 0,
\]
where \(\mathfrak m\) is the maximal ideal of \(A\).  The two fiber local rings have the same dimension \(n\), and \(B'/\mathfrak mB'\) is regular; hence \(I/\mathfrak m I=0\), whence \(I=0\).

In case (2), the conormal sequence
\[
        I/I^2\longrightarrow
        B\otimes_{B'}\Omega^1_{B'/A}
        \longrightarrow
        \Omega^1_{B/A}\longrightarrow 0
\]
is exact.  Since \(X\) is a relative complete intersection and \(U'\) is smooth, the first map is injective and the two modules which occur have the ranks prescribed by the virtual dimension.  These ranks force \(I/I^2=0\), hence \(I=0\).

\textbf{Definition 5.5.}  Under the hypotheses of Theorem 5.4, one writes simply
\[
        I(X\mid Y)=I^n(X\mid Y),\qquad
        J(X\mid Y)=J^n(X\mid Y).
\]
Thus \(J(X\mid Y)\) is the Jacobian subscheme measuring failure of smoothness.

We now extend the construction to noetherian formal schemes.  Let \(\mathfrak X\) be a noetherian formal scheme and \(E\) a coherent \(\Oo_{\mathfrak X}\)-module.  Define
\[
        I^p_{\mathfrak X}(E)(U)=I^p\bigl(\Gamma(U,E)\bigr)
\]
for every affine open \(U\subset\mathfrak X\).  If \(\mathfrak X=\varprojlim X_v\) and \(E=\varprojlim E_v\), then
\[
        I^p_{\mathfrak X}(E)=\varprojlim_v I^p_{X_v}(E_v).
\]

Let \(S\) be a noetherian formal scheme and let \(\mathfrak X\) be an \(S\)-adic formal scheme, essentially of finite type.  Let
\[
        \mathfrak X=\varprojlim X_v,\qquad S=\varprojlim S_v,\qquad
        X_v=\mathfrak X\times_S S_v.
\]
Then
\[
        \Omega^1_{\mathfrak X/S}\simeq \varprojlim_v\Omega^1_{X_v/S_v}
\]
is a coherent \(\Oo_{\mathfrak X}\)-module and is compatible with finite-type formal base change.

\textbf{Definition 5.6.}  Define
\[
        I^p(\mathfrak X\mid S)=I^p_{\mathfrak X}(\Omega^1_{\mathfrak X/S})
\]
and let \(J^p(\mathfrak X\mid S)\) be the closed formal subscheme of \(\mathfrak X\) defined by this ideal.  Equivalently
\[
        I^p(\mathfrak X\mid S)=\varprojlim_v I^p(X_v\mid S_v),
        \qquad
        J^p(\mathfrak X\mid S)=\varprojlim_v J^p(X_v\mid S_v).
\]
There is the natural cartesian comparison with the closed fiber:
\[
\begin{tikzcd}
J^p(X_0\mid S_0) \arrow[r] \arrow[d] &
J^p(\mathfrak X\mid S) \arrow[d]\\
X_0 \arrow[r] & \mathfrak X .
\end{tikzcd}
\]
If \(x\) is an isolated point of \(J^p(\mathfrak X\mid S)\), then the formal completion of \(J^p(\mathfrak X\mid S)\) at \(x\) is canonically the corresponding connected component.  In particular, if
\[
        |J^p(X_0\mid S_0)|=\{x_1,\ldots,x_m\}
\]
is finite and discrete, then
\[
        J^p(\mathfrak X\mid S)
        \simeq
        \coprod_{1\leq i\leq m} J^p(\mathfrak X\mid S)_{(x_i)}.
\]

\textbf{Corollary 5.7.}  Let \(S\) be noetherian and formal, and let \(\mathfrak X\) be an \(S\)-adic formal scheme essentially of finite type.
\begin{enumerate}[label=(\arabic*)]
\item The construction \(J^p(\mathfrak X\mid S)\) commutes with formal base change:
\[
        J^p(\mathfrak X'\mid S')=
        J^p(\mathfrak X\mid S)\times_S S',
        \qquad \mathfrak X'=\mathfrak X\times_S S'.
\]
\item The connected components of \(J^p(\mathfrak X\mid S)\) are in bijection with those of \(J^p(X_0\mid S_0)\), and for each isolated point \(x\) of \(J^p(X_0\mid S_0)\) there is a canonical isomorphism of completed local pieces
\[
        J^p(\mathfrak X\mid S)_{(x)}
        \simeq
        J^p(X_0\mid S_0)_{(x)} .
\]
In particular, in the finite discrete case one has the above disjoint-union decomposition.
\item If \(S\) is an ordinary noetherian scheme, \(S_0\subset S\) a closed subscheme, \(\widehat S\) the formal completion of \(S\) along \(S_0\), and
\[
        \widehat X=X\times_S\widehat S,
\]
then
\[
        J^p(\widehat X\mid \widehat S)
        =
        J^p(X\mid S)\times_S \widehat S.
\]
Thus the formal Jacobian subscheme is the completion of \(J^p(X\mid S)\) along \(J^p(X_0\mid S_0)\).
\item If a fixed integer \(n\) satisfies the hypotheses of Theorem 5.4 for all \(X_v\to S_v\), then \(x\notin J^n(\mathfrak X\mid S)\) if and only if \(\mathfrak X\) is formally smooth over \(S\) at \(x\).
\end{enumerate}

\emph{Proof.}  The assertions follow directly from the corresponding statements for ordinary schemes, the compatibility of \(\Omega^1_{\mathfrak X/S}\) with passage to \(X_v/S_v\), and the base-change formula for Fitting ideals.

\textbf{Remark 5.8.}  Let \(\mathfrak X\) be \(S\)-adic, essentially of finite type, and let \(x\in\mathfrak X\) with image \(y\in S\).  One can also consider the local formal \(S_{(y)}\)-scheme
\[
        \mathfrak X_{(x)}=\varprojlim_v \Spec(\Oo_{X_v,x_v}),
\]
with the natural morphism \(\mathfrak X_{(x)}\to\mathfrak X\).  Although \(\mathfrak X_{(x)}\) need not itself be essentially of finite type over \(S_{(y)}\), the completed module of differentials is still coherent in the sense needed here:
\[
        \Omega^1_{\mathfrak X_{(x)}/S_{(y)}}
        \simeq
        \varprojlim_v \widehat{\Omega^1_{\Oo_{X_v,x_v}/\Oo_{S_v,y_v}}}.
\]
This gives local formal subschemes \(J^p(\mathfrak X_{(x)}\mid S)\), and if \(x\) is an isolated point of \(J^p(\mathfrak X\mid S)\), the natural map
\[
        J^p(\mathfrak X_{(x)}\mid S)\longrightarrow J^p(\mathfrak X\mid S)
\]
identifies the source with the completed local component at \(x\).

We shall also need the image of the Jacobian subscheme in the base.  The set-theoretic image may be described by kernels
\[
        \Ker\bigl(\Oo_S\to f_*\Oo_{J^p(\mathfrak X\mid S)}\bigr),
\]
but the Fitting-ideal formulation is more useful.

\textbf{Definition 5.9.}  Let \(S\) be noetherian and formal, and let
\[
        f:\mathfrak X\longrightarrow S
\]
be \(S\)-adic, essentially of finite type.  Suppose \(J^p(\mathfrak X\mid S)\) is proper over \(S\).  Then \(f_*\Oo_{J^p(\mathfrak X\mid S)}\) is coherent, and we define an ideal sheaf on \(S\) by
\[
        I^p_S(\mathfrak X\mid S)
        =
        I^0_S\bigl(f_*\Oo_{J^p(\mathfrak X\mid S)}\bigr).
\]
Let \(J^p_S(\mathfrak X\mid S)\) be the closed formal subscheme of \(S\) defined by this ideal.  If the fixed-dimension hypotheses of Theorem 5.4 are in force, one writes
\[
        I_S(\mathfrak X\mid S)=I^n_S(\mathfrak X\mid S),
        \qquad
        J_S(\mathfrak X\mid S)=J^n_S(\mathfrak X\mid S).
\]

\textbf{Corollary 5.10.}  Under the hypotheses of Definition 5.9:
\begin{enumerate}[label=(\arabic*)]
\item The underlying set of \(J^p_S(\mathfrak X\mid S)\) is the image of \(J^p(\mathfrak X\mid S)\) under \(f:\mathfrak X\to S\).
\item The formation of \(J^p_S(\mathfrak X\mid S)\) commutes with base change.
\item If
\[
        J^p(\mathfrak X\mid S)=\coprod_i J^p(\mathfrak X\mid S)_{(x_i)}
\]
is the disjoint union of its isolated formal local components, then
\[
        I^p_S(\mathfrak X\mid S)
        =
        \prod_i I^p_{S,(x_i)}(\mathfrak X\mid S).
\]
\item In the case of completion of an ordinary noetherian \(S\)-scheme \(X\), the formal closed subscheme \(J^p_S(\widehat X\mid \widehat S)\) is the completion of \(J^p_S(X\mid S)\).
\end{enumerate}

\emph{Proof.}  Assertion (1) is the usual support property of the zeroth Fitting ideal.  Assertions (2) and (4) follow from base change for Fitting ideals.  Assertion (3) follows from the product formula for \(I^0\) applied to the decomposition of the coherent algebra \(f_*\Oo_{J^p(\mathfrak X\mid S)}\).

\textbf{Remark 5.11.}  The Jacobian ideal above is classical: it commutes with base change and, under the mild dimension hypotheses of Theorem 5.4, gives the smoothness criterion.  There are also more intrinsic ways of defining a canonical non-smooth locus in complete generality.  For an \(A\)-algebra \(B\), set
\[
        I_{B/A}
        =
        \bigcap_E
        \Ann_B D^1(B\mid A,E),
\]
where \(E\) runs through the finite-type \(B\)-modules.  This ideal commutes with localization by Corollary 3.13(c).  Hence for a morphism \(X\to Y\) one obtains an ideal sheaf \(I_{X/Y}\).  If \(X\to Y\) is essentially of finite type, then \(f\) is smooth at \(x\) if and only if
\[
        x\notin \Supp(\Oo_X/I_{X/Y})
\]
by 3.13(a).  The same construction carries over to adic formal schemes essentially of finite type and is therefore also suited to deformation theory.

\section*{6. Non-degenerate quadratic singularities}

The purpose of this section is to study deformations of a \(k\)-scheme whose singularities are isolated non-degenerate quadratic singularities.  In view of 4.13, under the hypotheses considered here the question becomes a local one.

We return to the notation of Section 4.  Thus \(A\) is a fixed complete noetherian local ring with residue field \(k\), and \(\mathcal C_A\) is the category of local artinian \(A\)-algebras with residue field \(k\).  Let \(X\) be a \(k\)-scheme of finite type and let \(x\) be a closed point of \(X\).  Recall that \(x\) is a non-degenerate quadratic singularity, rational over \(k\), if
\[
        \widehat{\Oo}_{X,x}
        \simeq
        k[[x_1,\ldots,x_n]]/(f)
\]
and
\[
        \left(\frac{\partial f}{\partial x_1},\ldots,
        \frac{\partial f}{\partial x_n}\right)=(x_1,\ldots,x_n).
\]
One may then arrange that \(f\) itself is a quadratic form.

\textbf{Lemma 6.1.}  Let
\[
        f\in k[[x_1,\ldots,x_n]]
\]
and suppose
\[
        \left(\frac{\partial f}{\partial x_1},\ldots,
        \frac{\partial f}{\partial x_n}\right)=\mathfrak m=(x_1,\ldots,x_n).
\]
Then there are parameters \(x'_1,\ldots,x'_n\), with
\[
        x_i\equiv x'_i\pmod{\mathfrak m^2},
\]
such that
\[
        f=q(x'_1,\ldots,x'_n)
\]
for a quadratic form \(q\).

\emph{Proof.}  Write \(f=f_2+h_1\), where \(f_2\) is quadratic and \(h_1\in\mathfrak m^3\).  Since the partial derivatives generate \(\mathfrak m\), there are elements \(e_i^{(1)}\in\mathfrak m^2\) such that
\[
        \sum_i e_i^{(1)}\frac{\partial f}{\partial x_i}
        \equiv h_1\pmod{\mathfrak m^4}.
\]
After replacing \(x_i\) by \(x_i-e_i^{(1)}\), the remainder is pushed into \(\mathfrak m^4\).  Repeating this construction, one obtains \(e_i^{(v)}\in\mathfrak m^{v+1}\) and parameters
\[
        x_i' = x_i-\sum_{v\geq 1}e_i^{(v)}
\]
for which all terms of degree at least \(3\) disappear.  Thus \(f\) becomes its quadratic part in the new variables.

\textbf{Lemma 6.2.}  Let
\[
        A_0\simeq k[[x_1,\ldots,x_n]]/(f),
\]
where \(f\) is a non-degenerate quadratic form, and let \((R,\mathfrak X)\) be a formal versal deformation of \(X_0=\Spec A_0\).  Then:
\begin{enumerate}[label=(\arabic*)]
\item The natural maps
\[
        J(\mathfrak X\mid R)\longrightarrow J_R(\mathfrak X\mid R)
        \longrightarrow \Spec A
\]
are isomorphisms.
\item The base Jacobian ideal is principal:
\[
        I_R(\mathfrak X\mid R)=(t),
        \qquad
        R=A[[t]],
\]
where \(t\) is a free formal generator of \(R\) over \(A\).
\item The morphism \(\mathfrak X\to\Spec A\) is formally smooth.
\end{enumerate}

\emph{Proof.}  By 4.10(b), the formal versal deformation has the explicit form
\[
        R=A[[t]],\qquad
        \mathfrak X=
        \Spf\bigl(R[[x_1,\ldots,x_n]]/(F-t)\bigr),
\]
where \(F\) is obtained from \(f\) by lifting the coefficients to \(A\).  The exact sequence
\[
        A_0\xrightarrow{dF}
        \bigoplus_i A_0\,dx_i
        \longrightarrow
        \Omega^1_{A_0/R}
        \longrightarrow 0
\]
shows that
\[
        I(\mathfrak X\mid R)=(\partial F/\partial x_1,\ldots,\partial F/\partial x_n)
        =(x_1,\ldots,x_n)
\]
on the special fiber.  Hence the Jacobian subscheme is obtained by setting \(x_1=\cdots=x_n=0\), and the remaining equation is \(t=0\).  Thus the maps in (1) are isomorphisms and \(I_R(\mathfrak X\mid R)=(t)\).  Finally
\[
        R[[x_1,\ldots,x_n]]/(F-t)
        \simeq
        A[[x_1,\ldots,x_n]]
\]
as an \(A\)-algebra, proving formal smoothness over \(A\).

\textbf{Corollary 6.3.}  Let \(x\in X_0\) be a non-degenerate quadratic singularity rational over \(k\).
\begin{enumerate}[label=(\arabic*)]
\item Let \((R,\mathfrak X_1)\) and \((R,\mathfrak X_2)\) be two deformations of \(X_0\) over an object \(R\) of \(\mathcal C_A\), with the same point above \(x\).  Then the completed local \(R\)-adic formal schemes \((\mathfrak X_1)_{(x)}\) and \((\mathfrak X_2)_{(x)}\) are isomorphic over \(R\) if and only if
\[
        I_R((\mathfrak X_1)_{(x)}\mid R)
        =
        I_R((\mathfrak X_2)_{(x)}\mid R).
\]
\item Let \(R\) be a regular local ring and \((R,\mathfrak X)\) a deformation of \(X_0\).  Then \(\Oo_{\mathfrak X,x}\) is regular if and only if
\[
        I_R(\mathfrak X_{(x)}\mid R)\not\subset \mathfrak m_R^2,
\]
that is, if and only if this Jacobian ideal defines a regular divisor in \(\Spf R\).
\end{enumerate}

\emph{Proof.}  Let the formal versal deformation be
\[
        A[[t]][[x_1,\ldots,x_n]]/(F-t).
\]
Two induced deformations over \(R\) are given by two elements \(a,b\in R\):
\[
        R[[x_1,\ldots,x_n]]/(F-a),
        \qquad
        R[[x_1,\ldots,x_n]]/(F-b).
\]
Their Jacobian ideals are \((a)\) and \((b)\).  Equality of these ideals gives \(a=vb\) for a unit \(v\).  Since \(R\) is complete and \(v\) has a square root after the standard lifting argument used here, the change of variables \(x_i\mapsto v^{1/2}x_i\) identifies the two equations, because \(F\) is quadratic.  The converse is immediate.  The second assertion follows from the fact that the quotient of the regular local ring \(R[[x_1,\ldots,x_n]]\) by \(F-a\) is regular precisely when \(a\notin\mathfrak m_R^2\).

\textbf{Theorem 6.4.}  Let \(X_0\) be a \(k\)-scheme of finite type with only isolated non-smooth points, and let \((R,\mathfrak X)\) be a formal versal deformation of \(X_0\).  Let
\[
        \{z_1,\ldots,z_m\}
\]
be the set of non-smooth points of \(X_0\).  Suppose
\[
        H^2(X_0,D^0)=0
\]
and suppose that all \(z_i\) are non-degenerate quadratic singularities, rational over \(k\).  Then:
\begin{enumerate}[label=(\arabic*)]
\item \(R\) is formally smooth over \(A\), i.e. \(R\) is a formal power-series ring over \(A\).
\item The Jacobian subscheme decomposes as a disjoint union
\[
        J(\mathfrak X\mid R)
        =
        \coprod_{i=1}^m J(\mathfrak X_{(z_i)}\mid R),
\]
and for each \(i\) the natural maps from the local Jacobian component to its image in \(\Spf R\) are isomorphisms.
\item The ideal \(I_R(\mathfrak X_{(z_i)}\mid R)\) is principal.  If
\[
        I_R(\mathfrak X_{(z_i)}\mid R)=(t_i),
\]
then \(t_1,\ldots,t_m\) can be extended to a system of free formal generators of \(R\) over \(A\), and
\[
        I_R(\mathfrak X\mid R)=\prod_{i=1}^m(t_i).
\]
\item The morphism \(\mathfrak X\to\Spec A\) is formally smooth.
\item If \(X_0\) is a complete irreducible curve, then
\[
        \dim_k\mathfrak m_R/\mathfrak m_R^2
        =
        3g-3+a,
\]
where
\[
        g=\dim_k H^1(X_0,\Oo_{X_0}),
        \qquad
        a=\dim_k H^0(X_0,D^1).
\]
\end{enumerate}

\emph{Proof.}  For each singular point \(z_i\), let \((R_i,\mathfrak U_i)\) be the formal versal deformation of \(\Spec(\Oo_{X_0,z_i})\).  Since \((R,\mathfrak X)\) induces a deformation of this local scheme, there is a continuous local map \(R_i\to R\) and a cartesian diagram
\[
\begin{tikzcd}
\mathfrak X_{(z_i)} \arrow[r] \arrow[d] &
\mathfrak U_i \arrow[d]\\
\Spf R \arrow[r] & \Spf R_i .
\end{tikzcd}
\]
Lemma 6.2 gives
\[
        J(\mathfrak X_{(z_i)}\mid R)
        \simeq
        J(\mathfrak U_i\mid R_i)\times_{\Spf R_i}\Spf R
\]
and
\[
        I_R(\mathfrak X_{(z_i)}\mid R)=(t_i),
\]
where \(t_i\) is the image in \(R\) of the corresponding local parameter.  Since \(H^2(X_0,D^0)=0\), the local-to-global criterion of 4.13 implies that
\[
        \Spf R\longrightarrow \prod_i\Spf R_i
\]
is formally smooth in the complementary directions; equivalently \(t_1,\ldots,t_m\) extend to a system of free formal generators of \(R\).  This proves (1), (2), and (3).  Assertion (4) follows because each local singular piece is formally smooth over \(\Spec A\) by Lemma 6.2(3), while away from the \(z_i\) the original scheme is already smooth.

For (5), formal smoothness of \(R\) over \(A\), together with the exact sequence obtained in 4.4 and 4.11,
\[
        0\to H^1(X_0,D^0)\to \mathfrak m_R/\mathfrak m_R^2
        \to H^0(X_0,D^1)\to 0,
\]
gives
\[
        \dim_k\mathfrak m_R/\mathfrak m_R^2
        =
        \dim_k H^1(X_0,D^0)+\dim_kH^0(X_0,D^1).
\]
Thus it remains to compute the Euler characteristic of \(D^0\) on the curve.  After extension of the base field to an algebraic closure, each non-degenerate quadratic singularity of a curve has completed local ring
\[
        k[[x,y]]/(xy).
\]
If \(A=\Oo_{X_0,z}\) and \(A'\) is its normalization, then
\[
        \ell(A'/A)=1,\qquad
        \Der_k(A,\mathfrak m_A)=\Der_k(A,A),
\]
and the conductor \(C\) of \(A\) in \(A'\) is \(\mathfrak m_A\).  Consequently
\[
        \Der_k(A,A)=\Der_k(A',A')
\]
as \(A\)-submodules of \(\Der_k(\Frac A,\Frac A)\), and one obtains an exact sequence
\[
        0\to D^0_{X_0}
        \to D^0_{\widetilde X_0}
        \to \Oo_{\widetilde X_0}/C
        \to 0,
\]
where \(\widetilde X_0\) denotes the normalization.  Since \(X_0\) is locally a complete intersection,
\[
        \ell(\Oo_{\widetilde X_0}/C)=2m.
\]
Riemann--Roch on the normalization gives
\[
        \chi(D^0_{\widetilde X_0})=3-3(g-m).
\]
Therefore
\[
        \chi(D^0_{X_0})
        =
        \chi(D^0_{\widetilde X_0})-\ell(\Oo_{\widetilde X_0}/C)
        =
        3-3g+m,
\]
which is equivalent to the stated formula.

\textbf{Remark 6.5.}  In the formula of Theorem 6.4(5), the integer
\[
        a=\dim_k H^0(X_0,D^1)
\]
is easily computed by Riemann--Roch.  Namely
\[
        a=
        \begin{cases}
        0,& g\geq 2,\\
        1,& g=1,\\
        3,& g=0.
        \end{cases}
\]
Finally, in characteristic \(2\) there is no non-degenerate quadratic form in an odd number of variables.  Equivalently, there is no non-degenerate quadratic singularity of even dimension when \(\operatorname{char}k=2\).  In that case one uses quadratic forms of maximum possible non-degeneracy.

\textbf{Definition 6.6.}  Let \(X\) be a \(k\)-scheme of finite type.  A closed point \(z\in X\) is called an ordinary quadratic singularity, rational over \(k\), if:
\begin{enumerate}[label=(\roman*)]
\item either \(\operatorname{char}k\neq 2\), or \(\operatorname{char}k=2\) and \(\dim\Oo_{X,z}\) is odd; in this case \(z\) is required to be a non-degenerate quadratic singularity rational over \(k\);
\item or \(\operatorname{char}k=2\) and \(\dim\Oo_{X,z}\) is even; after extension to an algebraic closure \(\bar k\), the completed local ring has the form
\[
        \bar k[[x_0,x_1,\ldots,x_{2n}]]
        /\bigl(x_0^2+x_1x_2+\cdots+x_{2n-1}x_{2n}\bigr).
\]
\end{enumerate}

\textbf{Lemma 6.7.}  Assume \(\operatorname{char}k=2\), and put
\[
        A_0=
        k[[x_0,x_1,\ldots,x_{2n}]]/(f),
        \qquad
        f=x_0^2+x_1x_2+\cdots+x_{2n-1}x_{2n}.
\]
Let \((R,\mathfrak X)\) be a formal versal deformation of \(X_0=\Spec A_0\).  Then
\[
        R=A[[t_1,t_2]],
\]
where \(t_1,t_2\) are free formal generators over \(A\),
\[
        I_R(\mathfrak X\mid R)=(t_1^2),
\]
and \(\mathfrak X\to\Spec A\) is formally smooth.

\emph{Proof.}  Since
\[
        \left(\frac{\partial f}{\partial x_0},
        \frac{\partial f}{\partial x_1},\ldots,
        \frac{\partial f}{\partial x_{2n}}\right)
        =
        (x_2,x_1,\ldots,x_{2n},x_{2n-1}),
\]
the deformation space has two formal parameters.  By 4.14 the versal deformation may be written
\[
        R=A[[t_1,t_2]],\qquad
        \mathfrak X=
        \Spf\bigl(R[[x_0,\ldots,x_{2n}]]/(F)\bigr),
\]
with
\[
        F=f+t_1+t_2x_0.
\]
The exact sequence for differentials gives
\[
        I(\mathfrak X\mid R)=(t_2,x_1,\ldots,x_{2n}),
\]
and the residual equation on the Jacobian locus is \(t_1\).  Thus the induced base ideal is \((t_1^2)\).  Moreover the equation can be rewritten so that the total formal scheme is formally smooth over \(A\), giving the last assertion.

\textbf{Corollary 6.8.}  Let \(X_0\) be a \(k\)-scheme of finite type with isolated non-smooth points only, and let \((R,\mathfrak X)\) be a formal versal deformation of \(X_0\).  Suppose
\[
        H^2(X_0,D^0)=0,
\]
\(\operatorname{char}k=2\), \(\dim X_0\) is even, and all non-smooth points \(z_1,\ldots,z_m\) are ordinary quadratic singularities.  Then:
\begin{enumerate}[label=(\arabic*)]
\item \(R\) is formally smooth over \(A\).
\item Each local base Jacobian ideal is principal of the form
\[
        I_R(\mathfrak X_{(z_i)}\mid R)=(t_i^2),
\]
and
\[
        I_R(\mathfrak X\mid R)=\prod_i (t_i^2).
\]
Moreover \(t_1,\ldots,t_m\) can be completed to a system of free formal generators of \(R\) over \(A\).
\item The morphism \(\mathfrak X\to\Spec A\) is formally smooth.
\item If \(k\) is algebraically closed and \(X_0\) is a complete irreducible curve, then
\[
        \dim_k\mathfrak m_R/\mathfrak m_R^2=3g-3+a,
\]
where
\[
        g=\dim_kH^1(X_0,\Oo_{X_0}),
        \qquad
        a=\dim_kH^0(X_0,D^1).
\]
\end{enumerate}

The proof is the same as that of Theorem 6.4, using Lemma 6.7 in place of Lemma 6.2 for the local calculation in characteristic \(2\).



% SGA 7-I Expose VII macro additions
\providecommand{\cU}{\mathcal U}
\providecommand{\T}{\mathcal T}
\providecommand{\e}{\mathrm e}
\providecommand{\Ens}{\operatorname{Ens}}
\providecommand{\Alg}{\operatorname{Alg}}
\providecommand{\Homens}{\operatorname{Hom}_{\mathrm{ens}}}
\providecommand{\Homalg}{\operatorname{Hom}_{\mathrm{alg}}}
\providecommand{\Hompt}{\operatorname{Hom}_{\mathrm{pt}}}
\providecommand{\EXT}{\operatorname{EXT}}
\providecommand{\TORS}{\operatorname{TORS}}
\providecommand{\BITORSRIG}{\operatorname{BITORSRIG}}
\providecommand{\BITORSBIRIG}{\operatorname{BITORSBIRIG}}
\providecommand{\BITORS}{\operatorname{BITORS}}
\providecommand{\pr}{\operatorname{pr}}
\providecommand{\id}{\mathrm{id}}
\providecommand{\sym}{\mathrm{sym}}
\providecommand{\h}{\mathrm h}
\providecommand{\cA}{\mathcal A}
\providecommand{\cB}{\mathcal B}


\section*{Exposé VII. Biextensions of sheaves of groups}
\begin{center}
by A. Grothendieck
\end{center}

\subsection*{0. Introduction}

The notion of a \emph{biextension} was introduced by D. Mumford, in the context of formal groups, in order to express conveniently the relations between the formal groups associated with an abelian scheme and with the dual abelian scheme.  In the present exposé we study this notion systematically in the general context of sheaves of groups on a fixed site.  We show in particular that it fits, in a remarkably simple way, into the classical cohomological formalism of \(\otimes^{L}\) and \(R\!\operatorname{Hom}\).  In the next exposé we shall spell out the most interesting special cases, in particular the case of biextensions of group schemes by \(\Gm\), working with sites such as the Zariski site, the fpqc site, or an intermediate site.  One will see, in particular, that if \(A\) and \(B\) are two abelian schemes over a scheme \(S\), then biextensions of \((A,B)\) by \(\Gm\) have a remarkable interpretation as, essentially, the classical divisor correspondences on \(A\) and \(B\).

Thus the cohomological meaning of the notion of divisor correspondence is made quite explicit; this meaning is already implicit in Tate duality for abelian varieties over \(p\)-adic fields.  This interpretation is also convenient for following the fate of such a correspondence when one lets \(A\) and \(B\), supposed defined over the fraction field of a discrete valuation ring, degenerate by means of Néron's theory.  The results of Exposé VIII show, for example, that one obtains a biextension of the pair of neutral components of the Néron models, hence a biextension of the connected special fibres of the Néron models, playing the role of a specialization of the original divisor correspondence.  These results will be used in Exposé IX, together with the monodromy theorem of Exposé I, 3, to prove results on Néron models of abelian varieties, the most important being the semi-stable reduction theorem for the Néron model.

Throughout what follows, unless explicitly stated otherwise, we work in a fixed \(\cU\)-topos, denoted \(\cT\).  Thus, in the terminology of SGA 4 IV, \(\cT\) is equivalent to the category of sheaves on a suitable \(\cU\)-site, which may be taken to be \(\cT\) itself endowed with its canonical topology.  We shall often identify, in terminology, the objects of \(\cT\) with sheaves of sets on \(\cT\).  The Groups, extensions of Groups, and so on which occur below are all relative to this topos \(\cT\).

\section*{1. Complements on extensions of Groups}

\subsection*{1.1. Extensions and bitorsors}

Let \(P\) and \(G\) be two Groups of \(\cT\), and let
\begin{equation}\tag{1.1.1}
  1\longrightarrow G\xrightarrow{i}E\xrightarrow{j}P\longrightarrow 1
\end{equation}
be an extension of \(P\) by \(G\).  Thus \(E\) is a Group, \(j:E\to P\) is an epimorphism of Groups, and \(i\) identifies \(G\) with \(\Ker j\).  We shall reinterpret the datum of such an extension in different terms, useful for defining and studying biextensions.

First, the homomorphism \(i\) makes \(G\) act on \(E\) on the left and on the right by
\[
  g\,x\,g'=i(g)\,x\,i(g')
\]
for \(g,g'\in G(S)\), \(x\in E(S)\), and any object \(S\) of \(\cT\).  If, via \(j\), one regards \(E\) as an object above \(P\), i.e. as an object of the induced topos \(\cT/P\), then this says equivalently that the induced group
\[
  G_P=G\times P
\]
over \(P\) acts on \(E\) on the left and on the right.  Moreover each of these actions makes \(E\) a torsor over \(P\), with structural group \(G_P\), and more precisely a bitorsor under \((G_P,G_P)\).

Recall that this means that \(E\) is a right torsor \(E_d\) for the right action of \(G_P\), and that the left action induces an isomorphism from \(G_P\) to the Group of automorphisms of \(E_d\), i.e. the sheaf of groups of automorphisms of \(E\) commuting with the right action of \(G_P\).  This condition is equivalent to the symmetric condition saying that the left action of \(G_P\) makes \(E\) a left torsor \(E_s\), and the right action induces an isomorphism from the opposite Group \(G_P^\circ\) of \(G_P\) to the Group of automorphisms of \(E_s\).

When \(\cT\) is the punctual topos, so that \(P\) and \(G\) are ordinary groups, the datum of a bitorsor \(E\) under \((G_P,G_P)\) is plainly equivalent to the datum, for every \(p\in P\), of a bitorsor \(E_p\) under \((G,G)\), namely the fibre of \(E\) over \(p\).  In the general case the analogous interpretation is: for every point \(p:S\to P\), where \(S\) is an arbitrary object of \(\cT\), one is given a bitorsor \(E_p\) under the pair of Groups \((G_S,G_S)\) on \(S\), in such a way that the formation of \(E_p\) is compatible with base change.  More explicitly, for a morphism \(S'\to S\) and the composite \(p':S'\to P\), one is given an isomorphism between the bitorsor \(E_{p'}\) and the inverse image of \(E_p\), subject to the usual transitivity condition for composites \(S''\to S'\to S\).

In canonical language this assertion says that the datum of a bitorsor under \((G_P,G_P)\) on \(P\) is equivalent to a Cartesian section, over the category \(\cT/P\), of the fibred category of bitorsors of structural group \((G_S,G_S)\) over a variable base \(S\).  More precisely, the functor ``value at \(P\)'' induces an equivalence from the category of such Cartesian sections to the category of bitorsors under \((G_P,G_P)\).  This is tautological, since \(P\) is the final object of the base category \(\cT/P\).  Nevertheless this interpretation is useful for formulating compatibilities in a language close to the set-theoretic case, using points with values in an object \(S\) rather than ordinary points.

The bitorsor structure just described uses only part of the multiplication on \(E\): namely multiplication of pairs of points of which at least one comes from \(G\).  We now interpret the full multiplicative structure of \(E\) as an additional structure on the bitorsor.  Consider two sections \(p,p'\) of \(P\).  With the preceding notation this gives two bitorsors \(E_p,E_{p'}\) under \((G,G)\), namely the inverse images in \(E\) of the sections \(p,p':e\to P\), where \(e\) denotes the final object of \(\cT\).  The multiplication of \(E\) induces a morphism of objects of \(\cT\)
\[
  \varphi:E_p\times E_{p'}\longrightarrow E_{pp'}.
\]
In terms of the left and right structures of the three bitorsors involved, it satisfies
\begin{equation}\tag{1.1.3.0}
\begin{aligned}
 \varphi(gx,x')&=g\varphi(x,x'),&
 \varphi(x,x'g')&=\varphi(x,x')g',\\
 \varphi(xg,x')&=\varphi(x,gx').
\end{aligned}
\end{equation}
The third relation can be interpreted by means of the contracted product of bitorsors,
\[
  E_p\overset{G}{\times}E_{p'},
\]
and says that \(\varphi\) factors as
\[
  E_p\times E_{p'}\longrightarrow
  E_p\overset{G}{\times}E_{p'}
  \xrightarrow{\;\varphi_{p,p'}\;}E_{pp'}.
\]
The two first relations then mean that
\begin{equation}\tag{1.1.3.1}
  \varphi_{p,p'}:E_p\overset{G}{\times}E_{p'}\xrightarrow{\sim}E_{pp'}
\end{equation}
is an isomorphism of bitorsors, the bitorsor structure on the left being the one for which the left operations of \(G\) act on the first factor and the right operations act on the second factor.  As every homomorphism of bitorsors under a fixed pair of Groups is necessarily an isomorphism, \(\varphi_{p,p'}\) is an isomorphism.

Again, the definition extends when \(p,p'\) are points of \(P\) with values in an arbitrary object \(S\) of \(\cT\), by applying what precedes in the induced topos \(\cT/S\).  The formation of the preceding \(\varphi_{p,p'}\)'s is compatible with arbitrary base change \(S'\to S\).  Equivalently, it gives a homomorphism of Cartesian sections over the category of double arrows \(p,p':S\rightrightarrows P\) with target \(P\).  The whole system may also be recovered from the universal case, \(S=P\times P\) and \(p,p'\) the two projections.  In that case it is an isomorphism of bitorsors
\begin{equation}\tag{1.1.3.2}
  \varphi:\pr_1^*(E)\overset{G}{\times}\pr_2^*(E)
     \xrightarrow{\sim}\pi^*(E),
\end{equation}
where \(\pi:P\times P\to P\) is multiplication in \(P\).  In practice the compatibilities are easiest to state using the data in the form (1.1.3.1), rather than by drawing less intuitive universal diagrams.

Writing the relations (1.1.3.0), which express the composition law of \(E\) through the bitorsor isomorphisms (1.1.3.1), uses associativity only when the product involves three points of \(E\) of which one at least comes from \(G\).  Full associativity is expressed by requiring that, for three points \(p,p',p''\) of \(P\), the following diagram commute:
\begin{equation}\tag{1.1.4.1}
\begin{tikzcd}[column sep=large,row sep=large]
& E_{pp'}E_{p''}\arrow[dr,"\varphi_{pp',p''}"]&\\
E_pE_{p'}E_{p''}\arrow[ur,"\varphi_{p,p'}\times\id"]\arrow[dr,"\id\times\varphi_{p',p''}"']&&E_{pp'p''}\\
& E_pE_{p'p''}\arrow[ur,"\varphi_{p,p'p''}"']&
\end{tikzcd}
\end{equation}
where the contracted product signs have been omitted.  The convention is that the contracted product \(Q.R\) of two bitorsors uses, for the quotient in \(Q\times R\), the right action of the structural group \(G\) on \(Q\) and the left action of \(G\) on \(R\).  This contracted product is associative up to canonical isomorphism, and we henceforth make this identification, writing both \((E_pE_{p'})E_{p''}\) and \(E_p(E_{p'}E_{p''})\) simply as \(E_pE_{p'}E_{p''}\).  Similar evident conventions will be used without further mention.

It is enough to verify this associativity in the universal case \(S=P\times P\times P\), \(p=\pr_1\), \(p'=\pr_2\), \(p''=\pr_3\).  We leave the reader to write the corresponding diagram in whatever notation he chooses.

Conversely, suppose \(P\) and \(G\) are given Groups of \(\cT\), together with a bitorsor
\begin{equation}\tag{1.1.5.1}
  j:E\longrightarrow P
\end{equation}
under \((G_P,G_P)\), and with a system of bitorsor isomorphisms over a variable base \(S\)
\begin{equation}\tag{1.1.5.2}
  \varphi_{p,p'}:E_pE_{p'}\xrightarrow{\sim}E_{pp'}
  \qquad (p,p'\in P(S)),
\end{equation}
compatible with base change \(S'\to S\), equivalently an isomorphism as in (1.1.3.2).  Suppose moreover that the corresponding associativity diagrams (1.1.4.1) commute.  Then there exists on \(E\) a unique structure of extension of \(P\) by \(G\) from which the data just described are obtained by the constructions above.

First take \(p=p'=1\), the unit section of \(P\), in (1.1.5.2).  One obtains an isomorphism
\begin{equation}\tag{1.1.5.3}
  \varphi_{1,1}:E_1\xleftarrow{\sim}E_1\overset{G}{\times}E_1.
\end{equation}
Taking the contracted product of both sides with the inverse bitorsor \(E_1^{-1}\) gives a canonical isomorphism
\begin{equation}\tag{1.1.5.4}
  E_1\xrightarrow{\sim}G_d^s,
\end{equation}
where \(G_d^s\) is the unit bitorsor, with \(G\) acting on itself by left and right translations.  Equivalently one has a section \(1_E\) of \(E_1\), hence of \(E\),
\begin{equation}\tag{1.1.5.5}
  1_E\in\Gamma(E_1)=\operatorname{Hom}(e,E_1)\subset\Gamma(E),
\end{equation}
characterized by
\[
  g1_E=1_Eg\quad(g\in G(S)),
\]
and by
\begin{equation}\tag{1.1.5.6}
  \varphi_{1,1}(1_E,1_E)=1_E.
\end{equation}
The system \(\varphi\) is equivalently a multiplication morphism
\begin{equation}\tag{1.1.5.7}
  \pi_E:E\times E\longrightarrow E,\qquad (x,y)\longmapsto xy,
\end{equation}
such that
\begin{equation}\tag{1.1.5.8}
  j(xy)=j(x)j(y).
\end{equation}
Finally define
\begin{equation}\tag{1.1.5.9}
  i:G\longrightarrow E,\qquad i(g)=g1_E=1_Eg.
\end{equation}
Then \(\pi_E\) makes \(E\) a Group with unit \(1_E\), and \(j:E\to P\) and \(i:G\to E\) are homomorphisms of Groups forming an extension of \(P\) by \(G\).  Associativity of \(\pi_E\) is exactly the commutativity of (1.1.4.1).  The construction by means of the \(\varphi_{p,p'}\)'s also shows that translations on \(E(S)\) by an arbitrary element are injective.  Since \(1_E1_E=1_E\), the element \(1_E\) is a two-sided identity.  The inverse of \(x\in E(S)\), lying over \(p\in P(S)\), is obtained by using the canonical isomorphism
\begin{equation}\tag{1.1.5.10}
  E_{p^{-1}}\xrightarrow{\sim}(E_p)^\circ.
\end{equation}
Thus \(E\) is a Group; moreover \(j\) is multiplicative by (1.1.5.8), and \(i\) is multiplicative by the calculation
\[
  i(g)i(g')=(g1_E)(g'1_E)=(g1_Eg')1_E=g(1_Eg')1_E=g(g'1_E)=i(gg').
\]
Since \(j\) is the structural projection of a torsor, it is an epimorphism of objects of \(\cT\), and \(i\) identifies \(G\) with \(\Ker j=E_1\).

We shall use the preceding arguments as proving that, for a fixed topos \(\cT\), the constructions above give an equivalence between the category of extensions (1.1.1) and the category of systems \((P,G,E,\varphi)\) where \(P\) and \(G\) are Groups of \(\cT\), \(E\) is a bitorsor under \((G_P,G_P)\) over \(P\), and \(\varphi\) is an isomorphism of bitorsors as in (1.1.3.2), equivalently a base-change-compatible system (1.1.3.1), satisfying associativity (1.1.4.1).  We also leave to the reader the evident description of inverse and direct images of extensions under a morphism of Groups
\[
  P'\longrightarrow P,\qquad\text{respectively}\qquad G\longrightarrow G',
\]
and of the compatibility with inverse image \(f^*\) for a morphism of topoi \(f:\cT'\to\cT\).

\subsection*{1.2. The commutative case}

From 1.4 onward we shall be concerned only with extensions of commutative Groups.  Thus \(P\) and \(G\) are assumed commutative.  We spell out the preceding description in terms of bitorsors in the case where \(E\) is a commutative Group.  A first necessary condition, expressing the fact that \(i(G)\) is a central sub-Group of \(E\), is that the bitorsors \(E_p\) be ordinary torsors; equivalently, the left and right actions of \(G\) on \(E_p\) coincide.  Once this condition is imposed, commutativity of \(E\) is expressed by the commutativity, for \(p,p'\in P(S)\), of the square
\begin{equation}\tag{1.2.1}
\begin{tikzcd}[column sep=large,row sep=large]
E_pE_{p'}\arrow[r,"\varphi_{p,p'}"]\arrow[d,"\sym"']&
E_{pp'}\arrow[d,"\id"]\\
E_{p'}E_p\arrow[r,"\varphi_{p',p}"']&
E_{p'p}.
\end{tikzcd}
\end{equation}
The first vertical arrow is the canonical commutativity isomorphism of the contracted product, defined because \(G\) is now commutative; and \(pp'=p'p\) because \(P\) is commutative.  The universal version over \(P\times P\) is equivalent.

\subsection*{1.3. Rigidified and birigidified torsors}

Let \(P\) be an object of \(\cT\) endowed with a section \(e_P\), let \(G\) be a Group of \(\cT\), and let \(E\) be a torsor under \(G_P\).  A \emph{rigidification} of \(E\), relative to \(e_P\), is a trivialization of the torsor
\[
  E_1=e_P^*(E)
\]
under \(G\), i.e. a section of this torsor.  The terminology is inspired by the case where \(G\) is commutative and
\begin{equation}\tag{1.3.1.1}
  \Hompt(P,G)\simeq (e),
\end{equation}
where the left-hand side denotes the set of morphisms transforming \(e_P\) into the unit section \(e_G\) of \(G\).  This condition is equivalent to the statement that the natural homomorphism
\begin{equation}\tag{1.3.1.2}
  \Gamma(G)\xrightarrow{\sim}\operatorname{Hom}(P,G)
\end{equation}
which associates to a section \(g\) of \(G\) the constant morphism \(P\to G\) with value \(g\), is an isomorphism.  Under this condition, every automorphism of a torsor respecting the rigidification is the identity.

Let \(Q\) also be an object of \(\cT\) with a section \(e_Q\), and let \(E\) be a torsor under \(G_{P\times Q}\).  The section \(e_Q\) defines a section \(e_{P\times Q/P}\) of \(P\times Q\) over \(P\); hence one may consider rigidifications \(\alpha\) of \(E\) along \(P\times e_Q\).  Interchanging \(P\) and \(Q\), one similarly has rigidifications \(\beta\) along \(e_P\times Q\).  Two such partial rigidifications are called \emph{compatible} if the sections \(\alpha\) and \(\beta\) coincide on \(e_P\times e_Q\), i.e. if they define the same section of the induced torsor \(E_{1,1}\) under \(G\).  A pair of compatible partial rigidifications is called a \emph{birigidification} of \(E\).

If \(G\) is commutative, a birigidification exists if and only if partial rigidifications exist, and one of the two partial rigidifications can be chosen in advance.  Indeed one then corrects the other by a morphism \(Q\to G\).  Moreover two birigidified torsors under \(G_{P\times Q}\) are isomorphic as torsors if and only if they are isomorphic as birigidified torsors, provided morphisms \(P\times Q\to G\) with prescribed restrictions are controlled as above.  Explicitly, if an isomorphism \(u:E\to E'\) of \(G_{P\times Q}\)-torsors changes the two rigidifications by morphisms \(f:P\to G\) and \(g:Q\to G\), with the same value at the marked point, then the correction
\[
  h(p,q)=f(p)g(q)a^{-1}
\]
gives an isomorphism respecting both rigidifications.

The same terminology extends to bitorsors.  A rigidification of a bitorsor \(E\) under \((G_P,G_P)\) is an isomorphism of \(e_P^*(E)\) with the trivial bitorsor \(G_d^s\), equivalently a central section \(e_E\) of \(e_P^*(E)\), i.e. one satisfying
\[
  g e_E=e_E g.
\]
Birigidifications for bitorsors are defined similarly.

Returning to two Groups \(P\) and \(G\), choose the unit section \(e_P\).  There are natural functors
\begin{equation}\tag{1.3.4.1}
  \EXT(P;G)\xrightarrow{\;i\;}\BITORSRIG(P,G)
  \xrightarrow{\;\delta\;}\BITORSBIRIG(P,P;G),
\end{equation}
where the three expressions denote respectively the category of extensions of Groups of \(P\) by \(G\), the category of bitorsors under \((G_P,G_P)\) rigidified at \(e_P\), and the category of bitorsors under \((G_{P\times P},G_{P\times P})\) birigidified at \((e_P,e_P)\).  The functor \(i\) assigns to an extension its corresponding bitorsor over \(P\), rigidified by the unit of \(E\).  The functor \(\delta\) is given by
\begin{equation}\tag{1.3.4.2}
  \delta(E)=\pi^*(E)\pr_1^*(E)^{-1}\pr_2^*(E)^{-1},
\end{equation}
where \(\pi:P\times P\to P\) is multiplication and \(\pr_1,\pr_2\) are the projections.  There is a canonical birigidification of \(\delta(E)\), obtained by restricting to \(P\times e_P\) and \(e_P\times P\), using the given rigidification of \(E\).

The composite \(\delta i\) is canonically isomorphic to the trivial functor: the multiplication isomorphism \(\varphi\) of (1.1.3.2) trivializes \(\delta i(E)\) compatibly with its canonical birigidification.

\textbf{Proposition 1.3.5.}  Let \(P\) and \(Q\) be two Groups of \(\cT\).
\begin{enumerate}[label=\alph*)]
\item Suppose that every morphism of sheaves of sets \(P\times P\to G\) which is trivial on \(P\times e_P\) and on \(e_P\times P\) is trivial.  Then the functor \(i\) in (1.3.4.1) is fully faithful.  In particular, on a bitorsor \(E\) under \((G_P,G_P)\), birigidified relative to \(e_P\), there is at most one extension structure of \(P\) by \(G\) compatible with the bitorsor structure and admitting \(e_E\) as unit.
\item Suppose moreover that every morphism \(P\times P\times P\to G\) trivial on the three partial products \(P\times P\times e_P\), \(P\times e_P\times P\), and \(e_P\times P\times P\), is trivial.  Let \(E\) be a torsor under \(G_P\).  An extension structure of \(P\) by \(G\) compatible with the torsor structure exists if and only if the birigidified torsor
\[
  \delta(E)=\pi^*(E)\pr_1^*(E)^{-1}\pr_2^*(E)^{-1}
\]
is trivial.  When this is so, every central section \(e_E\) of \(E_1=e_P^*(E)\), i.e. every rigidification of \(E\), determines a unique extension structure on \(E\), compatible with its bitorsor structure and admitting \(e_E\) as unit.  Equivalently, the rigidified bitorsor lies in the essential image of \(i\).  In addition there is a unique trivialization of the birigidified torsor \(\delta(E)\).
\end{enumerate}

In the language of (1.3.4.1), this says that the essential image of the fully faithful functor \(i\) is the kernel of \(\delta\); the diagram is therefore an exact sequence of pointed categories, and in fact of categories endowed with the pseudo-group structures introduced below.

For (a), let \(E,E'\) be two extensions of \(P\) by \(G\), and let \(f:E\to E'\) be a morphism of the underlying rigidified bitorsors.  We show that \(f\) is a morphism of extensions, equivalently a homomorphism of Groups.  It is enough to prove \(f(xy)=f(x)f(y)\).  Since both sides have the same image in \(P\), there is a well-defined morphism
\[
  u:P\times P\longrightarrow G
\]
such that
\[
  f(xy)=u(j(x),j(y))\,f(x)f(y).
\]
Putting \(x=1\) and \(y=1\), and using \(f(1)=1\), gives
\[
  u(p,1)=u(1,p)=1.
\]
By the hypothesis, \(u\) is trivial; hence \(f(xy)=f(x)f(y)\).

For (b), the necessity of the condition is the triviality of \(\delta i\).  Conversely, choose a trivialization \(\varphi\) of \(\delta(E)\).  Pulling it back by the unit section of \(P\times P\) gives a rigidification \(e_E\) of \(E\).  A chosen trivialization of \(\delta(E)\) is not necessarily compatible with this birigidification; its restrictions to \(P\times e_P\) and \(e_P\times P\) are expressed by central morphisms
\[
  f:P\to G,\qquad g:P\to G,
\]
with common value at the unit section.  The central morphism
\[
  h:P\times P\to G,\qquad h(p,q)=f(p)g(q)a^{-1}
\]
has these prescribed restrictions.  Replacing \(\varphi\) by the corrected trivialization gives one compatible with the birigidification.  The ambiguity is a central morphism \(P\times P\to G\) trivial on \(P\times e_P\) and \(e_P\times P\), hence trivial by hypothesis; thus the compatible trivialization is unique.  It remains only to verify associativity.  The two composites in (1.1.4.1) differ by a central morphism
\[
  u:P\times P\times P\to G,
\]
given by
\[
  (xy)z=u(jx,jy,jz)\,x(yz).
\]
Specializing successively \(y=z=1\), \(z=x=1\), and \(x=y=1\), and using the fact that \(1\) is a two-sided identity, shows that \(u\) is trivial on the three partial products; by the hypothesis it is trivial.  This proves (b).

\textbf{Corollary 1.3.6.}  Suppose every morphism of sheaves of sets \(P\times G\to G\) trivial on \(P\times e_G\) and on \(e_P\times G\) is trivial.  Then, if \(G\) is commutative, \(G\) is central in every extension of Groups of \(P\) by \(G\).  If \(P\) is also commutative and the hypothesis of Proposition 1.3.5(a) is satisfied, then every extension of \(P\) by \(G\) is commutative.

The first assertion follows by observing that the action of \(P\) on \(G\) defined by an extension is represented by a morphism \(u:P\times G\to G\) such that the corresponding morphism \((p,g)\mapsto u(p,g)g^{-1}\) is trivial on \(P\times e_G\) and on \(e_P\times G\).  Hence \(u(p,g)=g\).  For the second assertion, compare the extension \(E\) with its opposite \(E^\circ\).  If the extension is central and \(P=P^\circ\), then the uniqueness assertion in Proposition 1.3.5(a) gives \(E=E^\circ\).

A particularly interesting case where all the hypotheses just used are satisfied is
\begin{equation}\tag{1.3.7.1}
  \Hompt(P,G)=e
\end{equation}
(the final sheaf), where the left-hand side denotes morphisms of pointed sheaves from \(P\) to \(G\).  Equivalently,
\begin{equation}\tag{1.3.7.2}
  G\xrightarrow{\sim}\Homens(P,G),
\end{equation}
where the second member denotes morphisms of the underlying sheaves of sets.  This condition implies that, for every object \(Q\) of \(\cT\), the map
\begin{equation}\tag{1.3.7.3}
  \Homens(Q,G)\xrightarrow{\sim}\Homens(P\times Q,G),\qquad
  u\longmapsto u\circ\pr_2
\end{equation}
is bijective.  Hence the hypotheses of 1.3.5 and 1.3.6 are satisfied; the categories in (1.3.4.1) are rigid.  Since the condition is stable under all base change \(S\to e\), descent and the uniqueness part of 1.3.5(a) imply that a rigidified bitorsor \(E\) over \(P\) comes from an extension of \(P\) by \(G\) as soon as it does so locally over \(e\).  In view of 1.3.5(b), this is expressed by the vanishing of the section
\[
  \delta(E)\in\Gamma R^1 h_*(G_{P\times P}),
\]
where
\[
  h:P\times P\to e
\]
is the structural morphism.

Let \(p:P\to S\) be a coherent, that is quasi-compact and quasi-separated, flat morphism of schemes, and suppose
\begin{equation}\tag{1.3.8.1}
  p_*(\cO_P)\xleftarrow{\sim}\cO_S.
\end{equation}
Because \(p\) is coherent, the formation of \(p_*(\cO_P)\) commutes with all flat base change.  It follows that the conditions just considered for the relative scheme \(P/S\) are stable under Cartesian products, hence hold for \(P\times P\), \(P\times P\times P\), and so on, and are stable under flat base change.  If \(G\) is affine over \(S\), then (1.3.8.1) gives (1.3.7.2), since for \(G=\Spec A\),
\[
  \Hom_S(P,G)=\Homalg(A,p_*\cO_P)
     \simeq \Homalg(A,\cO_S).
\]
The same assertion applies to the Cartesian products, and, if \(G\) is flat over \(S\), to the relative scheme \(P_G\) over \(G_G\).  Hence, when \(P\) and \(G\) are group schemes over \(S\), viewed as sheaves on a topology on \((\Sch)/S\) no finer than the canonical topology, every morphism from \(P\), \(P\times P\), \(P\times P\times P\), and so on to \(G\), compatible with the marked sections, is trivial.  Also every \(S\)-morphism \(P\times G\to G\) compatible with marked sections is trivial.  Thus the hypotheses of 1.3.5 and 1.3.6 apply, and if \(G\) is flat the same is true for the corresponding relative statement.

The most interesting case in which (1.3.8.1) holds, and remains true after all base extension, is that of an abelian scheme \(P\) over \(S\).  Thus the conclusions of 1.3.5 and 1.3.6 are valid for every affine \(S\)-group \(G\).  This is the case first considered by J.-P. Serre, and the preceding developments are visibly inspired by it.  Another interesting case in which (1.3.8.1) holds, though not universally, is when \(P\) is the Néron model, over a discrete valuation base \(S\), of an abelian scheme over the fraction field.

\subsection*{1.4. Extensions of a free commutative Group}

Let \(I\) be an object of \(\cT\), and put
\begin{equation}\tag{1.4.1}
  P=\Z[I],
\end{equation}
the free \(\Z\)-Module generated by \(I\).  To an extension \(E\) of \(P\) by a given \(\Z\)-Module \(G\) we associate the corresponding torsor under \(G_P\), and then its restriction to \(I\) via the canonical morphism
\[
  I\longrightarrow P=\Z[I].
\]
This gives a canonical functor
\begin{equation}\tag{1.4.2}
  \EXT(P,G)\longrightarrow \TORS(I,G_I).
\end{equation}
This functor is an equivalence of categories.  Equivalently, for every \(G\) the maps
\begin{align}
  \Hom_{\Z}(P,G)&\longrightarrow H^0(I,G_I)=G(I), \tag{1.4.3}\\
  \operatorname{Ext}^{1}_{\Z}(P,G)&\longrightarrow H^1(I,G_I) \tag{1.4.4}
\end{align}
are isomorphisms.  More generally, for any ring object \(A\) of \(\cT\), not necessarily commutative, and \(P=A[I]\) the free \(A\)-Module generated by \(I\), there are functorial canonical isomorphisms
\begin{equation}\tag{1.4.5}
  \operatorname{Ext}^{i}_{A}(A[I],G)\xrightarrow{\sim}H^i(I,G_I),
\end{equation}
which for \(A=\Z\) and \(i=0,1\) are (1.4.3) and (1.4.4).  Indeed, the first member is the derived functor of
\[
  G\longmapsto\operatorname{Hom}_A(A[I],G),
\]
which is, by the definition of \(A[I]\), isomorphic to \(G\mapsto G(I)=H^0(I,G_I)\).  Since \(G\mapsto G_I\) sends injective Modules to injective Modules of \(\cT/I\), one obtains the asserted result.  The functor (1.4.2) is compatible with inverse images under a morphism of topoi
\[
  f:\cT'\longrightarrow\cT.
\]



% SGA 7-I Expose VII section 2 macro additions
\providecommand{\BIEXT}{\operatorname{BIEXT}}
\providecommand{\Biext}{\operatorname{Biext}}
\providecommand{\GRAB}{\operatorname{GRAB}}
\providecommand{\TORSBIRIG}{\operatorname{TORSBIRIG}}
\providecommand{\Corr}{\operatorname{Corr}}
\providecommand{\sym}{\mathrm{sym}}
\providecommand{\id}{\mathrm{id}}
\providecommand{\cY}{\mathcal Y}
\providecommand{\cC}{\mathcal C}
\providecommand{\e}{\mathrm e}

\section*{Exposé VII. Biextensions of sheaves of groups: section 2 and opening of section 3}

\section*{2. The notion of biextension of abelian sheaves}

\subsection*{2.0. Notation}
Throughout this section, \(P,Q,G\) denote commutative Groups of the topos \(\mathcal T\).  We shall often use the Cartesian square
\begin{equation}\tag{2.0.1}
\begin{tikzcd}
P\times Q \arrow[r] \arrow[d] & P\arrow[d]\\
Q\arrow[r] & e .
\end{tikzcd}
\end{equation}
The Groups \(G_P\), \(G_Q\), and \(G_{P\times Q}\) are the inverse images of \(G\) on \(P\), \(Q\), and \(P\times Q\).  We shall identify \(G_{P\times Q}\) equally with \((G_P)_{P\times Q}\) via \(\pr_1\), and with \((G_Q)_{P\times Q}\) via \(\pr_2\).  In this notation the group structures of \(P,Q,P\times Q\) do not intervene.  We shall almost never need the product group structure on \(P\times Q\).  What will be useful is to consider \(P\times Q\) as the inverse-image \(P\)-Group of \(Q\),
\begin{equation}\tag{2.0.2}
  P\times Q=Q_P,
\end{equation}
and symmetrically as the inverse-image \(Q\)-Group of \(P\),
\begin{equation}\tag{2.0.3}
  P\times Q=P_Q.
\end{equation}
Thus, over \(P\), one has the two Groups \(G_P\) and \(Q_P\); over \(Q\), one has \(G_Q\) and \(P_Q\).  If \(E\) is a torsor on \(P\times Q\) with group \(G_{P\times Q}\), then \(E\) may be regarded either as a torsor on \(P\) of group \((G_P)_{Q_P}\), or as a torsor on \(Q\) of group \((G_Q)_{P_Q}\).  It therefore makes sense to speak of a commutative extension structure of \(Q_P\) by \(G_P\) on \(E\), compatible with its torsor structure, and symmetrically of an extension structure of \(P_Q\) by \(G_Q\) on \(E\).

Applying the dictionary of 1.6 in the induced topoi, a structure of extension of \(Q_P\) by \(G_P\) on \(E\) is described by the following data.  For every object \(S\) of \(\mathcal T\), for \(p,p'\in P(S)\) and \(q\in Q(S)\), one is given an isomorphism of \(G_S\)-torsors
\begin{equation}\tag{2.0.4}
  \lambda_{p,p';q}: E_{p,q}\,E_{p',q}\xrightarrow{\sim}E_{pp',q}.
\end{equation}
Here \(E_{p,q}\) denotes the inverse image of \(E\) under the point \((p,q):S\to P\times Q\).  These isomorphisms are functorial in \(S\), i.e. compatible with pullback along morphisms \(S'\to S\) of \(\mathcal T\), and satisfy the associativity condition expressed by the commutativity of
\begin{equation}\tag{2.0.5}
\begin{tikzcd}[column sep=large,row sep=large]
& E_{pp',q}E_{p'',q}\arrow[dr,"\lambda_{pp',p'';q}"]&\\
E_{p,q}E_{p',q}E_{p'',q}
\arrow[ur,"\lambda_{p,p';q}\times\id"]
\arrow[dr,"\id\times\lambda_{p',p'';q}"']&&E_{pp'p'',q}\\
& E_{p,q}E_{p'p'',q}\arrow[ur,"\lambda_{p,p'p'';q}"']&
\end{tikzcd}
\end{equation}
and the commutativity condition
\begin{equation}\tag{2.0.6}
\begin{tikzcd}[column sep=large,row sep=large]
E_{p,q}E_{p',q}\arrow[r,"\lambda_{p,p';q}"]\arrow[d,"\sym"']&E_{pp',q}\arrow[d,"\id"]\\
E_{p',q}E_{p,q}\arrow[r,"\lambda_{p',p;q}"']&E_{p'p,q}.
\end{tikzcd}
\end{equation}

Symmetrically, a structure of extension of \(P_Q\) by \(G_Q\) on \(E\) is described by isomorphisms
\begin{equation}\tag{2.0.7}
  \mu_{p;q,q'}:E_{p,q}\,E_{p,q'}\xrightarrow{\sim}E_{p,qq'}
\end{equation}
for \(p\in P(S)\), \(q,q'\in Q(S)\), functorial in \(S\), satisfying the analogous associativity and commutativity conditions
\begin{equation}\tag{2.0.8}
\begin{tikzcd}[column sep=large,row sep=large]
& E_{p,qq'}E_{p,q''}\arrow[dr,"\mu_{p;qq',q''}"]&\\
E_{p,q}E_{p,q'}E_{p,q''}
\arrow[ur,"\mu_{p;q,q'}\times\id"]
\arrow[dr,"\id\times\mu_{p;q',q''}"']&&E_{p,qq'q''}\\
& E_{p,q}E_{p,q'q''}\arrow[ur,"\mu_{p;q,q'q''}"']&
\end{tikzcd}
\end{equation}
and
\begin{equation}\tag{2.0.9}
\begin{tikzcd}[column sep=large,row sep=large]
E_{p,q}E_{p,q'}\arrow[r,"\mu_{p;q,q'}"]\arrow[d,"\sym"']&E_{p,qq'}\arrow[d,"\id"]\\
E_{p,q'}E_{p,q}\arrow[r,"\mu_{p;q',q}"']&E_{p,q'q}.
\end{tikzcd}
\end{equation}

\textbf{Definition 2.1.}  Let \(P,Q,G\) be abelian Groups of \(\mathcal T\), and let \(E\) be a torsor under \(G_{P\times Q}\).  Suppose that \(E\) is endowed with a structure of extension of \(Q_P\) by \(G_P\), given by the \(\lambda\)'s above, and with a structure of extension of \(P_Q\) by \(G_Q\), given by the \(\mu\)'s above.  The two structures are said to be \emph{compatible} if, for every object \(S\) of \(\mathcal T\) and every \(p,p'\in P(S)\), \(q,q'\in Q(S)\), the following square of contracted products commutes:
\begin{equation}\tag{2.1.1}
\begin{tikzcd}[column sep=huge,row sep=large]
(E_{p,q}E_{p,q'})E_{p',q'}
\arrow[r,"\mu_{p;q,q'}\times\id"]
\arrow[d,"\id\times\lambda_{p,p';q'}"']
& E_{p,qq'}E_{p',q'}\arrow[d,"\lambda_{p,p';qq'}"]\\
E_{p,q}E_{pp',q'}\arrow[r,"\mu_{pp';q,q'}"']&E_{pp',qq'} .
\end{tikzcd}
\end{equation}
Equivalently, in the partial-composition notation below, this is the identity
\begin{equation}\tag{2.1.1.2}
  \lambda\bigl(\mu(x,y),\mu(z,t)\bigr)
  =
  \mu\bigl(\lambda(x,z),\lambda(y,t)\bigr),
\end{equation}
for \(x,y,z,t\) lying respectively in \(E_{p,q}(S),E_{p,q'}(S),E_{p',q}(S),E_{p',q'}(S)\), with the evident notation.

A \emph{biextension} of \((P,Q)\) by \(G\) is a torsor \(E\) under \(G_{P\times Q}\) endowed with these two extension structures, compatible in the preceding sense.

In more elementary language, a biextension structure on \(E\) is given by two partially defined laws of composition, denoted \(\lambda(x,y)\) and \(\mu(x,y)\).  The first is defined when \(\pr_2 j(x)=\pr_2 j(y)\), the second when \(\pr_1 j(x)=\pr_1 j(y)\), where \(j:E\to P\times Q\) is the structural morphism.  The axioms are the associativity and commutativity of these two partial laws, the relations
\begin{equation}\tag{2.1.1.1}
  \lambda(gx,y)=\lambda(x,gy)=g\lambda(x,y)
\end{equation}
and the analogous relations for \(\mu\), together with the mixed relation (2.1.1.2).

\subsection*{2.2. Unit sections}
There are unit sections for \(E\) considered as an extension of \(Q_P\) by \(G_P\),
\begin{equation}\tag{2.2.1}
  e_{E/P}\in \Gamma(E/P)=\operatorname{Hom}_P(P,E),
\end{equation}
and, symmetrically, for \(E\) considered as an extension of \(P_Q\) by \(G_Q\),
\begin{equation}\tag{2.2.2}
  e_{E/Q}\in \Gamma(E/Q)=\operatorname{Hom}_Q(Q,E).
\end{equation}
Thus the torsor \(E\) over \(P\times Q\) is canonically trivialized over \(1_P\times Q\), by \(e_{E/P}\), and over \(P\times1_Q\), by \(e_{E/Q}\).  These two trivializations agree over \(1_P\times1_Q\).  They therefore define one section
\begin{equation}\tag{2.2.3}
  e_E\in\Gamma(E),
\end{equation}
called, by abuse of language, the unit or neutral section of \(E\).  The equality follows by applying the compatibility relation to the four sections \(e,e',e',e\) over \(1_P\times1_Q\): the two sides are respectively \(e\) and \(e'\).  Consequently the two partial laws \(\lambda\) and \(\mu\) coincide on the fibre over \(1_P\times1_Q\), both being transported from the group law of \(G\).

Restricting the extension of \(Q_P\) by \(G_P\) to the unit section of \(P\) gives an extension of \(Q\) by \(G\) which is canonically split by \(e_{E/P}\).  The symmetric statement holds for the restriction over the unit section of \(Q\).

\subsection*{2.3. Morphisms of biextensions}
Let \(E\) be a biextension of \((P,Q)\) by \(G\), and let \(E'\) be a biextension of \((P',Q')\) by \(G'\).  A morphism of biextensions from \(E\) to \(E'\) is a system \((u,v,w,f)\), where
\begin{equation}\tag{2.3.1}
  u:P\to P',\qquad v:Q\to Q',\qquad w:G\to G'
\end{equation}
are morphisms of Groups and
\begin{equation}\tag{2.3.2}
  f:E\to E'
\end{equation}
is a morphism of the underlying sheaves of sets such that \(f\) is both a homomorphism of the relative extensions over \(P\) and \(P'\), associated with the relative homomorphisms \(u\times v\) and \(u\times w\), and a homomorphism of the relative extensions over \(Q\) and \(Q'\), associated with \(u\times v\) and \(v\times w\).

The maps \(u,v,w\) are determined by \(f\): one recovers \(u\times v\) by passing to the quotient, hence \(u\) and \(v\) by restriction to \(P\times1_Q\) and \(1_P\times Q\), and one recovers \(w\) by restricting \(f\) to \(G\cdot e_E\).  We shall therefore often identify the morphism with \(f\) itself.  Adding source and target to morphisms in the usual canonical way, biextensions form a category.  There is a canonical functor
\begin{equation}\tag{2.3.3}
  \BIEXT(\mathcal T)\longrightarrow \GRAB(\mathcal T)\times\GRAB(\mathcal T)\times\GRAB(\mathcal T),
\end{equation}
where \(\GRAB(\mathcal T)\) denotes abelian Groups of \(\mathcal T\).  It is useful to decompose it as
\begin{equation}\tag{2.3.4}
  \BIEXT(\mathcal T)\longrightarrow\GRAB(\mathcal T)\times\GRAB(\mathcal T),
\end{equation}
by \(E\mapsto(P,Q)\), and
\begin{equation}\tag{2.3.5}
  \BIEXT(\mathcal T)\longrightarrow\GRAB(\mathcal T),
\end{equation}
by \(E\mapsto G\).

\subsection*{2.4. The categories \(\BIEXT(P,Q;G)\) with variables}
The functor (2.3.4) is fibred, while the functor (2.3.5) is cofibred.  Thus, given a biextension \(E\) of \((P,Q)\) by \(G\), and morphisms \(u:P'\to P\), \(v:Q'\to Q\), there is an inverse-image biextension \(E'\) of \((P',Q')\) by \(G\), together with a universal morphism \(E'\to E\) compatible with \(u,v,\id_G\).  Similarly, given \(w:G\to G'\), there is a direct-image biextension \(E'\) of \((P,Q)\) by \(G'\), together with a universal morphism \(E\to E'\) compatible with \(\id_P,\id_Q,w\).  These are the usual constructions for extensions of Groups, paraphrased in the present setting.

For fixed \(P,Q,G\), write
\begin{equation}\tag{2.4.1}
  \BIEXT(P,Q;G)
\end{equation}
for the fibre category over \((P,Q,G)\).  This category is covariant in \(G\) and contravariant in \(P,Q\), and additively dependent on each of its three arguments.  For the argument \(G\), for example, this says that \(\BIEXT(P,Q;0)\) is equivalent to the punctual category, and that for two abelian Groups \(G,G'\) of \(\mathcal T\), the functor
\begin{equation}\tag{2.4.2}
  \BIEXT(P,Q;G\times G')\longrightarrow
  \BIEXT(P,Q;G)\times\BIEXT(P,Q;G')
\end{equation}
induced by the two projections is an equivalence.  In the language of the cited reference, for fixed \(P,Q\) the category of biextensions of \((P,Q)\) by variable Groups \(G\) is an additive cofibred category over \(\GRAB(\mathcal T)\).

\textbf{Remark 2.4.3.}  With the terminology of loc. cit., this cofibred category is not only additive but left exact.  We shall not verify this here, since it will follow below, in 3.5, from more precise facts.  The general theory would then describe this cofibred category, up to equivalence, in terms of a chain complex \(L_\bullet\) in \(\GRAB(\mathcal T)\), depending only on \(P\) and \(Q\).  Below we shall give a direct construction of this complex as \(P\otimes^L Q\), the tensor product being taken in a derived category.

\subsection*{2.5. The structure of \(\BIEXT(P,Q;G)\)}
The formal facts for additive cofibred categories give useful consequences.  For fixed \(G\), two biextensions \(E,E'\) of \((P,Q)\) by \(G\) have a product biextension \(E.E'\).  This construction is the analogue of the product of two extensions of one commutative Group by another, and is compatible with viewing \(E,E'\) either as extensions of \(Q_P\) by \(G_P\), or as extensions of \(P_Q\) by \(G_Q\).  As a \(G_{P\times Q}\)-torsor, \(E.E'\) is the product of the torsors underlying \(E\) and \(E'\).  This law is associative and commutative up to canonical isomorphism and has a two-sided unit, the trivial biextension.  This unit can be described as the trivial torsor
\begin{equation}\tag{2.5.1.1}
  P\times Q\times G,
\end{equation}
with its two relative group structures induced from the product group structures on \(Q\times G\) and \(P\times G\).  Every biextension \(E\) also has an inverse \(E^{-1}\), characterized up to unique isomorphism by an isomorphism \(E.E^{-1}\simeq1\).  Thus \(\BIEXT(P,Q;G)\) is a pseudo-group in the 2-category of categories.

It follows that, for any biextension \(E\), its endomorphism monoid is canonically isomorphic to the endomorphism monoid of the trivial biextension.  This latter monoid is a commutative group, denoted
\begin{equation}\tag{2.5.2.1}
  \Biext^0(P,Q;G).
\end{equation}
Moreover the product law on biextensions makes the set
\begin{equation}\tag{2.5.2.2}
  \Biext^1(P,Q;G)
\end{equation}
of isomorphism classes of biextensions of \((P,Q)\) by \(G\) into a commutative group.

For a homomorphism \(G\to G'\), the induced functor
\begin{equation}\tag{2.5.3.1}
  \BIEXT(P,Q;G)\longrightarrow\BIEXT(P,Q;G')
\end{equation}
is compatible with products and inverses up to canonical isomorphism.  Hence it induces homomorphisms
\begin{equation}\tag{2.5.3.2}
  \Biext^i(P,Q;G)\longrightarrow\Biext^i(P,Q;G')\qquad (i=0,1).
\end{equation}
For \(i=0\), this is the natural map on automorphism groups of a biextension and its image by \(G\to G'\).

The group \(\Biext^0(P,Q;G)\) is explicit.  An automorphism \(u\) of a biextension \(E\) is first an automorphism of the underlying \(G_{P\times Q}\)-torsor, and is therefore given by a morphism
\begin{equation}\tag{2.5.4.0}
  u:P\times Q\longrightarrow G,
\end{equation}
where \(u(x)=u(j(x))\,x\) in torsor notation.  Compatibility with the first extension structure is equivalent to additivity in the second variable; compatibility with the second extension structure is equivalent to additivity in the first variable.  Thus \(u\) is an automorphism of the biextension structure if and only if it is bilinear.  This gives the canonical bijection, in fact isomorphism of groups,
\begin{equation}\tag{2.5.4.1}
  \Biext^0(P,Q;G)\simeq\operatorname{Hom}(P\otimes Q,G),
\end{equation}
functorial in \(P,Q,G\).  Below we shall determine \(\Biext^1(P,Q;G)\) as an \(\Ext^1\)-group.

\subsection*{2.6. Variance in \(P\), \(Q\), and \(G\)}
The contravariance in \(P,Q\) commutes with the covariance in \(G\).  Hence the categories just considered can be regarded as fibres of a tri-additive cofibred category
\begin{equation}\tag{2.6.1}
  \BIEXT'(\mathcal T)\longrightarrow
  \GRAB(\mathcal T)^\circ\times\GRAB(\mathcal T)^\circ\times\GRAB(\mathcal T),
\end{equation}
where the structural functor is \(E\mapsto(P,Q,G)\).  The pseudo-group structures obtained by using the additivity in any one of the three variables are canonically isomorphic.  Consequently the abelian groups
\begin{equation}\tag{2.6.2}
  \Biext^i(P,Q;G)\qquad (i=0,1)
\end{equation}
are trilinear functors: contravariant in \(P,Q\), covariant in \(G\), and valued in abelian groups.  The isomorphism (2.5.4.1) is functorial for this full variance, and the same will hold for the later description of \(\Biext^1\) by an \(\Ext^1\).

\subsection*{2.7. The symmetric biextension}
If \(E\) is a biextension of \((P,Q)\) by \(G\), define its symmetric biextension \({}^sE\), a biextension of \((Q,P)\) by \(G\), as follows.  The underlying sheaf of sets and the \(G\)-action are the same.  The structural morphism is the composite of \(E\to P\times Q\) with the symmetry \(P\times Q\to Q\times P\).  The left and right extension structures of \({}^sE\) are respectively the right and left extension structures of \(E\).  The operation \(E\mapsto{}^sE\) is an involutive automorphism of \(\BIEXT(\mathcal T)\):
\begin{equation}\tag{2.7.1}
  {}^s({}^sE)=E.
\end{equation}
Thus every assertion concerning biextensions has a symmetric assertion obtained by interchanging left and right.  In practice only one of the two symmetric statements will usually be written out.

\subsection*{2.8. Change of topos and localization}
Let
\begin{equation}\tag{2.8.0}
  f:\mathcal T'\longrightarrow\mathcal T
\end{equation}
be a morphism of topoi.  The inverse-image functor \(f^*\) sends a biextension \(E\) of \((P,Q)\) by \(G\) to a biextension \(f^*E\) of \((f^*P,f^*Q)\) by \(f^*G\).  This is another instance of the fact that such a structure is expressible by finite projective limits and arbitrary inductive limits.  Compatibility of the two extension structures on \(f^*E\) follows from the universal form of the commutative diagram (2.1.1).  Hence one obtains a canonical functor
\begin{equation}\tag{2.8.1}
  f^*:\BIEXT(\mathcal T)\longrightarrow\BIEXT(\mathcal T').
\end{equation}
All constructions with biextensions used later are compatible with inverse images by morphisms of topoi; in particular this includes the product \(E.E'\) and inverse \(E^{-1}\) of 2.5.

Taking \(\mathcal T'=\mathcal T/S'\) and \(f=i_{S'}:\mathcal T/S'\to\mathcal T\), one obtains the biextension induced by \(E\) on \(S'\), denoted \(E|S'\) or \(E_{S'}\).  The usual transitivity of restriction implies that the categories \(\BIEXT(\mathcal T/S)\), for variable \(S\in\mathcal T\), are the fibres of a fibred category over \(\mathcal T\).  This fibred category is a stack on \(\mathcal T\) in Giraud's sense: morphisms and objects can be glued.  Thus a biextension on \(S\), up to canonical isomorphism, is obtained from a covering family \(S_i\to S\), biextensions \(E_i\) on \(S_i\), and descent isomorphisms on \(S_i\times_SS_j\) satisfying the usual cocycle condition.  This stack is the stack of biextensions on \(\mathcal T\).  The functors (2.3.3)--(2.3.5), and the operations \(E.E'\) and \(E^{-1}\), are morphisms of stacks or stack operations in this sense.

\subsection*{2.9. Relation with birigidified torsors}
We saw in 2.2 that every biextension \(E\) of \((P,Q)\) by \(G\) has an underlying \(G_{P\times Q}\)-torsor canonically birigidified with respect to the unit sections of \(P\) and \(Q\).  This gives a canonical functor
\begin{equation}\tag{2.9.1}
  \BIEXT(P,Q;G)\longrightarrow\TORSBIRIG(P,Q;G),
\end{equation}
where the target is the category of birigidified torsors on \(P\times Q\), with structural group \(G_{P\times Q}\).

Using the unit sections, for every finite product of factors isomorphic to \(P\) or \(Q\) one has a canonical inclusion of each factor.  We shall therefore say that a morphism from such a product to \(G\) is trivial on a given factor.  Consider the products
\begin{equation}\tag{2.9.2}
  P\times Q,
  \quad P\times Q\times Q,
  \quad P\times Q\times Q\times Q,
  \quad P\times P\times Q,
  \quad P\times P\times P\times Q,
  \quad P\times P\times Q\times Q.
\end{equation}

\textbf{Proposition 2.9.3.}  Suppose that every morphism from one of the products (2.9.2) to \(G\), trivial on each of its factors, is trivial.  Then:
\begin{enumerate}[label=\alph*)]
\item The functor (2.9.1) is fully faithful, and both categories are rigid.  In particular, if \(E\) is a \(G_{P\times Q}\)-torsor birigidified for \((e_P,e_Q)\), then there is at most one biextension structure on \(E\) compatible with its birigidified torsor structure.
\item Such a structure exists on \(E\) if and only if there exists on \(E\) a structure of extension of \(Q_P\) by \(G_P\), compatible with its \(G_P\)-torsor structure and the rigidification on \(P\times e_Q\), and a structure of extension of \(P_Q\) by \(G_Q\), compatible with its \(G_Q\)-torsor structure and the rigidification on \(e_P\times Q\).  In that case both extension structures are unique and are compatible, hence define a biextension.
\end{enumerate}

For (a), the automorphism group of a birigidified torsor on \(P\times Q\) is the group of morphisms \(P\times Q\to G\) trivial on the two partial factors; by the hypothesis it is the identity.  If \(E,E'\) are biextensions and \(f:E\to E'\) is a morphism of birigidified torsors, Proposition 1.3.5(a), applied to the two relative extension structures, shows that \(f\) is a morphism of biextensions.  For (b), Proposition 1.3.5(b) gives the unique relative extension structures; the two composites in the compatibility diagram (2.1.1) differ by a morphism \(P\times P\times Q\times Q\to G\), trivial on each factor, hence trivial by the hypothesis.

\textbf{Corollary 2.9.4.}  If \(P,Q,G\) satisfy
\begin{equation}\tag{2.9.4.1}
  \operatorname{Hom}_{\mathrm{pt}}(P,G)=e,
  \qquad
  \operatorname{Hom}_{\mathrm{pt}}(Q,G)=e,
\end{equation}
then the hypotheses of Proposition 2.9.3 hold.  Moreover, for a birigidified \(G_{P\times Q}\)-torsor \(E\), the existence of a compatible biextension structure is equivalent to the condition that the two pointed morphisms which it defines,
\begin{equation}\tag{2.9.4.2}
  Q\longrightarrow R^1p_*(G_{P}),
  \qquad
  P\longrightarrow R^1q_*(G_{Q}),
\end{equation}
where \(p:P\to e\) and \(q:Q\to e\) are the structural morphisms, are morphisms of Groups.

The first assertion follows from 1.3.7.  For the second, multiplicativity of the first morphism says that the rigidified torsor locally satisfies the condition of 1.3.5(b) over \(P\); by 1.3.7 it then satisfies it globally.  The symmetric argument applies to the second morphism, and Proposition 2.9.3(b) gives the conclusion.

\textbf{Example 2.9.5.}  Let \(S\) be a scheme, let \(P,Q\) be abelian schemes over \(S\), and let \(G=\mathbb G_m\).  As noted in 1.3.8, condition (2.9.4.1) holds.  A birigidified \(\mathbb G_m\)-torsor on \(P\times Q\) is the same as a birigidified invertible Module on \(P\times Q\).  A divisor correspondence on \((P,Q)\) is the isomorphism class of such a birigidified invertible Module.  With the fppf topology, the two morphisms (2.9.4.2) are the morphisms
\begin{equation}\tag{2.9.5.1}
  Q\longrightarrow \Pic(P/S),
  \qquad
  P\longrightarrow \Pic(Q/S)
\end{equation}
defined by the divisor correspondence; by the theorem of the square these are necessarily compatible with the group structures.  Thus the category of biextensions of \((P,Q)\) by \(\mathbb G_m\) is equivalent to the category of birigidified invertible Modules on \(P\times Q\).  This category is rigid, and its set of isomorphism classes is the set of divisor correspondences on \((P,Q)\).  The identification respects the natural group laws.

\textbf{Remark 2.9.6.}  Additivity of the first morphism in (2.9.4.2) is equivalent to saying that the second morphism factors through the subsheaf \(\Ext^1(Q,G)\); the canonical map
\begin{equation}\tag{2.9.6.0}
  \Ext^1(Q,G)\longrightarrow R^1q_*(G_Q)
\end{equation}
is injective by 1.3.5.  Symmetrically, additivity of the second morphism is equivalent to the factorization of the first through \(\Ext^1(P,G)\).  Hence one obtains group isomorphisms
\begin{equation}\tag{2.9.6.1}
  \Biext(P,Q;G)
  \simeq \operatorname{Hom}\bigl(Q,\Ext^1(P,G)\bigr)
  \simeq \operatorname{Hom}\bigl(P,\Ext^1(Q,G)\bigr).
\end{equation}
In the situation of Example 2.9.5 these become the classical isomorphisms
\begin{equation}\tag{2.9.6.2}
  \Corr(P,Q)\simeq \operatorname{Hom}(P,Q')\simeq\operatorname{Hom}(Q,P'),
\end{equation}
where \(P'=\Ext^1(P,\mathbb G_m)\) and \(Q'=\Ext^1(Q,\mathbb G_m)\) are the dual abelian schemes, identified with \(\Pic^0_{P/S}\) and \(\Pic^0_{Q/S}\).

\subsection*{2.10. Various generalizations of biextensions}
We mention these only for completeness.

First, one can define biextensions of \((P,Q)\) by \(G\) without assuming that \(P\) and \(Q\) are commutative, while still assuming \(G\) commutative.  Such an object is a torsor under \(G_{P\times Q}\), endowed with central extension structures of \(Q_P\) by \(G_P\) and of \(P_Q\) by \(G_Q\), compatible with the torsor structure and with one another by diagrams of the form (2.1.1).  Most of the preceding development applies mutatis mutandis.  When \(P\) and \(Q\) are also commutative, this is more general than Definition 2.1, which corresponds to the case in which both relative extension structures are commutative.  In the non-commutative case the category still depends additively on \(G\), but not on \(P\) and \(Q\).

Second, for a fixed integer \(n\geq1\), one defines in the same way an \(n\)-extension of a sequence of \(n\) commutative Groups \((P_1,\ldots,P_n)\) by a commutative Group \(G\): it is a torsor over \(P_1\times\cdots\times P_n\) under the inverse image of \(G\), with \(n\) partial extension structures on the products obtained by omitting one factor, and a compatibility relation of type (2.1.1) for every pair of distinct indices.  When \(n=3\), the \(P_i\) are abelian schemes over \(S\), and \(G=\mathbb G_m\), every such \(n\)-extension is trivial.

Third, biextensions of \(A\)-Modules may be defined.  Let \(A\) be a commutative ring object of \(\mathcal T\), and let \(P,Q,G\) be \(A\)-Modules.  A biextension of \(A\)-Modules of \((P,Q)\) by \(G\) is a biextension in the sense of 2.1, in which the two partial extension structures are enriched to extensions of relative \(A_P\)- and \(A_Q\)-Modules, compatible with the relative Module structures on \(G_P,Q_P\) and \(G_Q,P_Q\).  The two relative Module structures on \(E\) amount to two actions of the multiplicative monoid underlying \(A\) on the sheaf underlying \(E\), which are required to commute.  Each action must be distributive with respect not only to its corresponding partial composition law on \(E\), but also with respect to the other partial composition law.

The dictionary of 1.1.6 also has an \(A\)-linear variant: an extension of the \(A\)-Module \(P\) by \(G\) corresponds to a family of torsors \(E_p\) under \(G_S\), parameterized by points \(p\in P(S)\), which varies \(A\)-linearly in \(p\).  In addition to the isomorphisms \(\varphi_{p,p'}\) of (1.1.3.1), one has isomorphisms
\begin{equation}\tag{2.10.3.1}
  \mu_{p,a}: a_*(E_p)\xrightarrow{\sim} E_{ap}
  \qquad(p\in P(S),\ a\in A(S)),
\end{equation}
where the first member denotes the image of the \(G_S\)-torsor \(E_p\) by extension of scalars \(g\mapsto ag:G_S\to G_S\).  The conditions on \(\varphi\) and \(\mu\) are the evident ones ensuring that they come from an extension of Modules.



% SGA 7-I pages 184--212 macro additions
\providecommand{\Biext}{\operatorname{Biext}}
\providecommand{\GrAb}{\operatorname{GrAb}}
\providecommand{\Coker}{\operatorname{Coker}}
\providecommand{\coker}{\operatorname{coker}}
\providecommand{\Ker}{\operatorname{Ker}}
\providecommand{\cofib}{\operatorname{cofib}}
\providecommand{\Dcat}{\mathrm D}
\providecommand{\Kcat}{\mathrm K}
\providecommand{\cPsi}{\mathcal\Psi}
\providecommand{\cY}{\mathcal Y}
\providecommand{\cC}{\mathcal C}
\providecommand{\cT}{\mathcal T}
\providecommand{\Hh}{\mathrm H}
\providecommand{\Tot}{\operatorname{Tot}}
\providecommand{\id}{\operatorname{id}}


\section*{Expos\'e VII, \S3. Complements on homological algebra}

The principal aim of this number is to define the isomorphism announced above,
\begin{equation}\tag{3.0.1}
  \Biext^1(P,Q;G)\simeq \Ext^1(P\otimes^{L}Q,G).
\end{equation}
For this purpose we need some homological algebra, useful independently, occupying \S\S3.1--3.4.  These developments may be regarded as complements to the theory of additive cofibred categories used earlier.

\subsection*{3.1. The additive cofibred category attached to a chain complex}
Let \(\cC\) be an abelian category and let
\begin{equation}\tag{3.1.1}
  \cdots\longrightarrow L_2\xrightarrow{d_2}L_1\xrightarrow{d_1}L_0
\end{equation}
be a chain complex in \(\cC\), written with lower indices and differential of degree \(-1\).  Thus \(L_i=0\) for \(i<0\).  We define a cofibred category \(\cPsi_L\) over \(\cC\).  An object of the fibre \(\cPsi_L(A)\) is a pair \((X,a)\) consisting of an extension
\begin{equation}\tag{3.1.2}
  0\longrightarrow A\longrightarrow X\xrightarrow{j}L_0\longrightarrow 0
\end{equation}
and a partial trivialization of its inverse image by \(d_1:L_1\to L_0\).  Equivalently it is a lifting
\begin{equation}\tag{3.1.3}
\begin{tikzcd}[column sep=large,row sep=large]
& X \arrow[d,"j"]\arrow[dl] & \\
A \arrow[r] & L_0 & L_1\arrow[l,"d_1"']\arrow[ul,"a"]
\end{tikzcd}
\qquad\text{with}\qquad a\,d_2=0.
\end{equation}
This description depends only on the truncation of \(L_\bullet\) to order \(1\), obtained from \(L_\bullet\) by killing the homology in degrees \(\geq2\).  A morphism \((X,a)\to(X',a')\) is a morphism of extensions over \(L_0\) compatible with the identity on \(A\) and satisfying \(u a=a'\).

The obvious functor
\begin{equation}\tag{3.1.4}
  \cPsi_L\longrightarrow \cC,
  \qquad (X,a)\longmapsto A,
\end{equation}
is cofibred.  For a morphism \(u:A\to B\) and an object \((X,a)\in\cPsi_L(A)\), one takes the push-forward extension \(u_*(X)\) of \(L_0\) by \(B\), equipped with the induced partial splitting.  The canonical map \((X,a)\to(u_*(X),u_*a)\) is cocartesian by the universal property of push-forward of extensions.

The cofibred category \(\cPsi_L\) is additive: \(\cPsi_L(0)\) is equivalent to the punctual category and, for any \(A,B\),
\begin{equation}\tag{3.1.6}
  \cPsi_L(A\oplus B)\longrightarrow \cPsi_L(A)\times\cPsi_L(B)
\end{equation}
is an equivalence.  We shall use freely the consequences of additivity recalled earlier.  The trivial object of \(\cPsi_L(A)\) is obtained from the split extension \(L_0\oplus A\) and the lifting \((d_1,0):L_1\to L_0\oplus A\).  The automorphism group of any object in \(\cPsi_L(A)\) is canonically the automorphism group of the trivial object; explicitly
\begin{equation}\tag{3.1.8}
  \cPsi^0_L(A)
  \simeq \Ker\bigl(\Hom(L_0,A)\to\Hom(L_1,A)\bigr)
  \simeq \Hom(H_0(L_\bullet),A).
\end{equation}
The set of isomorphism classes in \(\cPsi_L(A)\) is likewise a commutative group, denoted \(\cPsi^1_L(A)\), additive in \(A\).

\subsection*{3.2. Determination of \(\cPsi_L^1(A)\)}
Let \((X,a)\in\cPsi_L(A)\).  Consider the cochain complex of length one
\begin{equation}\tag{3.2.1}
  K^\bullet=[X\xrightarrow{j}L_0]
\end{equation}
with \(X\) in degree \(0\).  It is a resolution of \(A\), hence represents \(A\) in \(\Dcat(\cC)\).  The diagram defining \(a\) gives a degree-one morphism of complexes
\begin{equation}\tag{3.2.2}
  u:L_\bullet\longrightarrow K^\bullet[-1],
\end{equation}
which may be displayed as the anticommutative square diagram
\begin{equation}\tag{3.2.3}
\begin{tikzcd}[column sep=large,row sep=large]
 L_2\arrow[r,"d_2"]\arrow[d,"0"']&L_1\arrow[r,"d_1"]\arrow[d,"a"]&L_0\arrow[d,"\id"]\arrow[r]&0\arrow[d]\\
 0\arrow[r]&X\arrow[r,"j"]&L_0\arrow[r]&0 .
\end{tikzcd}
\end{equation}
Thus one obtains, in the derived category, a degree-one morphism
\begin{equation}\tag{3.2.4}
  c(X,a):L_\bullet\longrightarrow A[1],
\end{equation}
and hence a canonical map
\begin{equation}\tag{3.2.4'}
  \cPsi_L^1(A)\longrightarrow \Ext^1(L_\bullet,A)
  =\Hom_{\Dcat(\cC)}(L_\bullet,A[1]).
\end{equation}

\textbf{Theorem 3.2.5.}  The map \((3.2.4')\) is an isomorphism of groups, functorial in \(A\).

We recall the proof.  Additivity and functoriality in \(A\) are immediate from the construction.  For injectivity, suppose that the morphism of complexes \(u\) in (3.2.3) is zero in \(\Dcat(\cC)\).  Then it is already zero in the homotopy category \(\Kcat(\cC)\): there is a morphism \(s:L_0\to X\) such that
\begin{equation}\tag{3.2.5.1}
  js=\id_{L_0},\qquad s d_1=-a.
\end{equation}
Indeed, after composing with a suitable quasi-isomorphism to another resolution of \(A\), one may suppose the target resolution has no terms in degrees \(\geq2\), and the required homotopy factors uniquely through the fibre product defining \(X\).  Such an \(s\) is precisely a splitting of the extension \(X\) of \(L_0\) by \(A\), compatible with the partial splitting \(a\), hence \((X,a)\) is trivial in \(\cPsi^1_L(A)\).

For surjectivity, represent an element of \(\Ext^1(L_\bullet,A)\) by a degree-one morphism to a resolution \(C^\bullet\) of \(A\), which may be taken with \(C^i=0\) for \(i\geq2\):
\begin{equation}\tag{3.2.5.2}
\begin{tikzcd}[column sep=large,row sep=large]
L_2\arrow[r]\arrow[d]&L_1\arrow[r]\arrow[d,"-u^0"]&L_0\arrow[d,"u^1"]\arrow[r]&0\\
0\arrow[r]&C^0\arrow[r]&C^1\arrow[r]&0 .
\end{tikzcd}
\end{equation}
The object \(C^0\) is an extension of \(C^1\) by \(A\); pulling it back by \(L_0\to C^1\) gives an extension \(X\) of \(L_0\) by \(A\).  The anticommutativity of the square forces \(-u^0\) to factor through a map \(a:L_1\to X\) lifting \(d_1\), and the relation on \(L_2\) gives \(a d_2=0\).  Then \((X,a)\in\cPsi_L(A)\), and the original class is exactly \(c(X,a)\).  This proves the theorem.

\subsection*{3.3. Properties of \(\cPsi_L\)}
The cofibred additive category \(\cPsi_L\) is left exact in the sense of the preceding theory; it admits a quasi-maximal object; every object has a typical complex; and the typical pro-complex of \(\cPsi_L\) is canonically isomorphic in \(\Dcat(\cC)\) to the truncation of \(L_\bullet\) used in \S3.1.

Left exactness means the following.  If
\begin{equation}\tag{3.3.2.1}
  0\longrightarrow A'\longrightarrow A\xrightarrow{v}A''
\end{equation}
is exact, then the natural functor
\begin{equation}\tag{3.3.2.2}
  \cPsi_L(A')\longrightarrow \cPsi_L([A\to A''])
\end{equation}
is an equivalence.  Here \(\cPsi_L([A\to A''])\) is the category of pairs consisting of an object over \(A\) and a trivialization of its image over \(A''\).  For \(\cPsi_L\), such an object is a triple \((X,a,p)\), with \((X,a)\) as above and \(p:X\to A''\) satisfying \(p i=v\) and \(p a=0\).  Then \(X'=\Ker(p)\) is an extension of \(L_0\) by \(A'=\Ker v\), and \(a\) factors through \(X'\).

A quasi-maximal object is constructed as follows.  Put
\begin{equation}\tag{3.3.3.1}
  L'_1=\Coker(d_2:L_2\to L_1).
\end{equation}
Let \(E_0=(X_0,a_0)\), where \(X_0=L_0\oplus L'_1\) is the split extension of \(L_0\) by \(L'_1\), and where \(a_0:L_1\to X_0\) has components \(d_1\) and the quotient map \(L_1\to L'_1\).  Every object \((X,a)\) maps, after a monomorphism on the base object, to an object whose underlying extension is split, and every such object is dominated by \(E_0\).

If \(E=(X,a)\) lies above \(A\), the functor
\begin{equation}\tag{3.3.4.1}
  B\longmapsto \Hom(E,B)
\end{equation}
is represented by \(\Coker(a)\), and the canonical map from \(A\) is the composite
\begin{equation}\tag{3.3.4.2}
  A\longrightarrow X\longrightarrow \Coker(a).
\end{equation}
The complex \([A\to\Coker(a)]\) is the typical complex of \(E\).  Applying the general theory to a quasi-maximal object gives the typical complex of \(\cPsi_L\).  In the present construction it is canonically the truncation of \(L_\bullet\) itself.

For every short exact sequence as above, the Ext long exact sequence
\begin{equation}\tag{3.3.6.1}
\begin{aligned}
0&\to\Ext^0(L,A')\to\Ext^0(L,A)\to\Ext^0(L,A'')\\
 &\to\Ext^1(L,A')\to\Ext^1(L,A)\to\Ext^1(L,A'')\to\cdots
\end{aligned}
\end{equation}
corresponds, under (3.1.8) and (3.2.5), to the geometric exact sequence obtained from the left exactness of \(\cPsi_L\).  The boundary morphism differs from the geometric one by the conventional sign described in the source; the editor explicitly notes that these sign questions will not be used in the subsequent expos\'es.

Finally, for variable \(L\), the theory also describes the variance of \(\cPsi_L\).  The class of cocartesian functors \(\cPsi_L\to\cPsi_{L'}\) corresponds canonically to isomorphisms \(L'\to L\) in \(\Dcat(\cC)\).  We shall use only the elementary form developed next.

\subsection*{3.4. Variance with respect to \(L\)}
Let
\begin{equation}\tag{3.4.1}
  f:L'\longrightarrow L
\end{equation}
be a morphism of chain complexes, given in low degrees by
\begin{equation}\tag{3.4.2}
\begin{tikzcd}[column sep=large,row sep=large]
\cdots\arrow[r]&L'_2\arrow[r]\arrow[d,"f_2"']&L'_1\arrow[r]\arrow[d,"f_1"']&L'_0\arrow[r]\arrow[d,"f_0"']&0\\
\cdots\arrow[r]&L_2\arrow[r]&L_1\arrow[r]&L_0\arrow[r]&0.
\end{tikzcd}
\end{equation}
There is a canonical cocartesian functor
\begin{equation}\tag{3.4.3}
  f^*:\cPsi_L\longrightarrow\cPsi_{L'}.
\end{equation}
If \((X,a)\in\cPsi_L(A)\), then \(f^*(X,a)=(X',a')\), where \(X'\) is the inverse image of the extension \(X\) of \(L_0\) by \(A\) along \(f_0:L'_0\to L_0\), and \(a'\) is induced from \(a\) by the commutative square involving \(f_1\).

The functor (3.4.3) gives, for \(i=0,1\), homomorphisms
\begin{equation}\tag{3.4.4}
  f^i:\cPsi_L^i\longrightarrow\cPsi_{L'}^i,
\end{equation}
and the morphism of complexes gives
\begin{equation}\tag{3.4.5}
  f^i:\Ext^i(L,-)\longrightarrow\Ext^i(L',-).
\end{equation}
For \(i=0,1\), these are compatible with the identifications (3.1.8) and (3.2.5):
\begin{equation}\tag{3.4.6}
\begin{tikzcd}[column sep=large,row sep=large]
\cPsi_L^i(A)\arrow[r,"\sim"]\arrow[d]&\Ext^i(L,A)\arrow[d]\\
\cPsi_{L'}^i(A)\arrow[r,"\sim"']&\Ext^i(L',A).
\end{tikzcd}
\end{equation}

\textbf{Proposition 3.4.7.}  The functor (3.4.3) is faithful, respectively fully faithful, respectively an equivalence of categories, if and only if
\begin{equation}\tag{3.4.7.1}
  H_0(L')\longrightarrow H_0(L)
\end{equation}
is an epimorphism, respectively \(H_0(L')\to H_0(L)\) is an isomorphism and \(H_1(L')\to H_1(L)\) is an epimorphism, respectively both \(H_0\) and \(H_1\) maps are isomorphisms.  This is just the preceding compatibility plus the elementary criterion for a homomorphism of additive cofibred categories to be faithful, fully faithful, or an equivalence.

Now take a short exact sequence of complexes
\begin{equation}\tag{3.4.8.1}
  0\longrightarrow L'\longrightarrow L\longrightarrow L''\longrightarrow0.
\end{equation}
The expected transitivity of functors of the form (3.4.3) gives an equivalence of categories
\begin{equation}\tag{3.4.8.2}
  \cPsi_{L'}(A)\xrightarrow{\sim}\cPsi_L(A)\times_{\cPsi_{L''}(A)}\{0\},
\end{equation}
where \(0\) is the trivial object of \(\cPsi_{L''}(A)\).  In concrete terms, an object on the right is an object \((X,a)\in\cPsi_L(A)\) together with a trivialization of its image in \(\cPsi_{L''}(A)\); the usual left exactness for extensions shows that this is the same as an object of \(\cPsi_{L'}(A)\).

The long exact sequence in \(\Ext\) associated with (3.4.8.1)
\begin{equation}\tag{3.4.9.1}
\begin{aligned}
0&\to\Ext^0(L'',A)\to\Ext^0(L,A)\to\Ext^0(L',A)\\
 &\to\Ext^1(L'',A)\to\Ext^1(L,A)\to\Ext^1(L',A)\to\Ext^2(L'',A)\to\cdots
\end{aligned}
\end{equation}
also has a geometric interpretation through \(\cPsi_L\).  The boundary map \(\Ext^0(L',A)\to\Ext^1(L'',A)\) sends a morphism \(u:L'_1\to A\), with \(u d=0\), to the class of the object of \(\cPsi_{L''}(A)\) obtained by lifting the trivial object and twisting its automorphism by \(u\).  The compatibility is reduced in the source to the corresponding statement of the general theory, after passing to the opposite category; the editor again warns that one must be careful with signs, but these signs do not enter the later expos\'es.

\subsection*{3.5. A canonical partial flat resolution of a \(\Z\)-module}
We apply the preceding material to the abelian category \(\GrAb(\cT)\) of abelian groups, or \(\Z\)-modules, in the topos \(\cT\).  Let \(P\) be such a group.  We define functorially an augmented chain complex
\begin{equation}\tag{3.5.1}
  0\longrightarrow L_2(P)\xrightarrow{d_2}L_1(P)\xrightarrow{d_1}L_0(P)\xrightarrow{\epsilon}P\longrightarrow0
\end{equation}
with \(L_i(P)=0\) for \(i\geq3\), where
\begin{equation}\tag{3.5.1'}
  L_0(P)=\Z[P],\qquad
  L_1(P)=\Z[P\times P],\qquad
  L_2(P)=\Z[P\times P\times P]\oplus\Z[P\times P].
\end{equation}
For \(p,q,r\in P(S)\), write \([p]\), \([p,q]\), \([p,q,r]\) for the corresponding generators.  The maps are
\begin{align}\tag{3.5.2}
 \epsilon([p])&=p,\\
 d_1([p,q])&=[p+q]-[p]-[q],\\
 d_2([p,q,r])&=[p+q,r]-[p,q+r]+[p,q]-[q,r],\\
 d_2'([p,q])&=[p,q]-[q,p].
\end{align}
The relations \(\epsilon d_1=0\) and \(d_1 d_2=0\) are just the commutativity and associativity of the addition law on \(P\).  The terms \(\Z[P^n]\) are flat \(\Z\)-modules in the topos.

\textbf{Proposition 3.5.3.}  One has \(H_i(L_\bullet(P))=0\) for \(i>0\), and the augmentation induces \(H_0(L_\bullet(P))\simeq P\).

The proof reduces to the following criterion supplied by Proposition 3.4.7.  It is enough to show that, for every abelian group \(G\) of the topos, the functor
\begin{equation}\tag{3.5.3.1}
  \cPsi_P(G)\longrightarrow \cPsi_{L(P)}(G)
\end{equation}
is an equivalence.  The left-hand side is the category of extensions of the \(\Z\)-module \(P\) by \(G\).  On the right, using \S1.4, extensions of the free modules \(\Z[P]\), \(\Z[P\times P]\), and \(\Z[P^3]\oplus\Z[P^2]\) by \(G\) are described as torsors over the corresponding bases.  Thus an object is a torsor \(E\) under \(G_P\), together with a trivialization of
\begin{equation}\tag{3.5.3.2}
  \pr_1^*E\;\pr_2^*E\;(+)^*E^{-1}
\end{equation}
over \(P\times P\), satisfying the commutativity and associativity conditions encoded by the two summands of \(L_2(P)\).  This is precisely the bitorsor dictionary of \S1.1 for a commutative extension of \(P\) by \(G\).

\textbf{Remark 3.5.6.}  It would be useful to strengthen the preceding partial resolution to a full functorial flat resolution of \(P\), with terms functorially sums of free modules generated by Cartesian powers of \(P\).  Such a resolution would give a spectral sequence whose first term is expressed in terms of cohomology groups of powers \(P^n\), and which would calculate all \(\Ext^i(P,A)\), not only the partial information just obtained.  S. Kleiman constructed such a resolution, according to the proof-corrector's note in the source.

There is also a variant for \(A\)-modules, where \(A\) is a fixed ring of the topos.  The shape of the truncated complex is more complicated because the corresponding extension dictionary is more complicated:
\begin{align}\tag{3.5.9}
L_0(P)&=A[P],\\
L_1(P)&=A[P\times P]\oplus A[P],\\
L_2(P)&=A[P^3]\oplus A[P^2]\oplus A[P^2]\oplus A[P].
\end{align}
The formulas (3.5.2) are supplemented by relations expressing compatibility of scalar multiplication:
\begin{align}\tag{3.5.10}
 d_1([a,p])&=[ap]-a[p],\\
 d_2([a,p,q])&=[a,p+q]-[a,p]-[a,q],\\
 d_2([a,b,p])&=[ab,p]-a[b,p]-[a,bp],
\end{align}
and the remaining formulas say exactly that the internal addition on the torsor and the external action of \(A\) make it an \(A\)-module extension of \(P\) by \(G\).

\subsection*{3.6. Application to biextensions}
Let \(P,Q\) be two \(\Z\)-modules of the topos \(\cT\).  Choose flat resolutions \(L_\bullet(P)\) and \(L_\bullet(Q)\), agreeing in dimensions \(\leq2\) with the partial complexes just constructed.  There are canonical maps
\begin{equation}\tag{3.6.1}
  L_\bullet(P)\to P,
  \qquad
  L_\bullet(Q)\to Q,
\end{equation}
and hence
\begin{equation}\tag{3.6.2}
  L_\bullet(P,Q):=L_\bullet(P)\otimes L_\bullet(Q)
  \longrightarrow P\otimes^{L}Q.
\end{equation}
This map is an isomorphism in degrees \(\leq2\), for the purposes of the preceding construction.

\textbf{Theorem 3.6.4.}  The additive cofibred category, over \(\GrAb(\cT)\), whose fibre over \(G\) is the category of biextensions of \((P,Q)\) by \(G\), is canonically equivalent to the additive cofibred category
\begin{equation}\tag{3.6.4.1}
  \cPsi_{L(P,Q)}.
\end{equation}

Together with Theorem 3.2.5, this gives the announced result.

\textbf{Corollary 3.6.5.}  For \(P,Q,G\in\GrAb(\cT)\), there is a canonical isomorphism, functorial in \(G\),
\begin{equation}\tag{3.6.5.1}
  \Biext^1(P,Q;G)\xrightarrow{\sim}\Ext^1(P\otimes^{L}Q,G).
\end{equation}

The proof of Theorem 3.6.4 is obtained by writing explicitly the first three terms of \(L(P,Q)=L(P)\otimes L(Q)\).  They are
\begin{align}\tag{3.6.6}
L_0(P,Q)&=\Z[P\times Q],\\
L_1(P,Q)&=\Z[P\times Q\times Q]\oplus\Z[P\times P\times Q],\\
L_2(P,Q)&=\Z[P\times Q^3]\oplus\Z[P\times Q^2]\oplus\Z[P^2\times Q^2]\\
&\quad\oplus\Z[P^3\times Q]\oplus\Z[P^2\times Q].
\end{align}
Using the dictionary of \S1.4, an extension of these free modules by \(G\) is described as a system of torsors over the corresponding Cartesian products.  Thus an object of \(\cPsi_{L(P,Q)}(G)\) is a torsor \(E\) on \(P\times Q\), together with two partial composition laws: one in the \(Q\)-direction and one in the \(P\)-direction.  The five components of \(L_2(P,Q)\) express the associativity in each variable, commutativity in each variable, and compatibility of the two partial laws.  These are exactly the axioms of a biextension of \((P,Q)\) by \(G\).  The verification of compatibility of morphisms is formal and is left, as in the source, to the reader.

\textbf{Remark 3.6.7.}  Each of the generalizations of biextensions mentioned in \S2.10 gives a corresponding variant of Theorem 3.6.4 and Corollary 3.6.5.  The case of any number of factors is formal; the case of \(A\)-modules follows from the variant of the construction of \(L_\bullet(P)\) in Remark 3.5.6.  The non-commutative case is more interesting homologically: to a Group \(P\), not necessarily commutative, one attaches a two-truncated complex \(K_\bullet(P)\), analogous to \(L_\bullet(P)\) but with \(K_2(P)=\Z[P\times P\times P]\) only.  Then \(H_0(K_\bullet(P))\) is the abelianization \(P/[P,P]\), identified with \(H_1(P,\Z)\), and \(H_1(K_\bullet(P))\) is identified with \(H_2(P,\Z)\).  Thus \(K_\bullet(P)\) is the truncation to order two of the usual chain complex computing group homology.  The same argument gives the corresponding cofibred-category statement.



% SGA 7-I pages 213--236 macro additions
\providecommand{\Biext}{\operatorname{Biext}}
\providecommand{\Ext}{\operatorname{Ext}}
\providecommand{\Torop}{\operatorname{Tor}}
\providecommand{\GRAB}{\operatorname{GRAB}}
\providecommand{\RHom}{\operatorname{RHom}}
\providecommand{\Ker}{\operatorname{Ker}}
\providecommand{\coker}{\operatorname{coker}}
\providecommand{\cK}{\mathcal K}
\providecommand{\cL}{\mathcal L}
\providecommand{\cM}{\mathcal M}
\providecommand{\cT}{\mathcal T}
\providecommand{\dfn}{\mathrel{\mathrm{dfn}}}


\section*{Exposé VII, continuation: compatibilities and trivializations}

\subsection*{3.7. Various compatibilities}

The preceding discussion first gives, for variable commutative Groups, an equivalence between central extensions of a commutative Group \(P\) by variable commutative Groups \(G\), and the cofibred category \(\Psi_{K(P)}\), where \(K(P)\) denotes the complex attached to \(P\).  Similarly, for two Groups \(P,Q\), the cofibred category of biextensions of \((P,Q)\) is equivalent to the cofibred category
\[
        \Psi_{L(P,Q)},\qquad L(P,Q)=K(P)\otimes K(Q).
\]
These remarks suggest that, for \(P\) commutative, one may perhaps find a canonical resolution of \(P\) of the kind desired in 3.5.4(a), by a judicious modification of the chain complex \(C(P,\mathbb Z)\).

In the present section we examine the behaviour of the equivalence of cofibred categories 3.6.4 when \(P\) and \(Q\) vary.  The variance of the cofibred category of biextensions of \((P,Q)\) by variable \(G\), with respect to the variables \(P,Q\), was defined in principle in 2.4, and may also be formulated by saying that one has a cofibred functor (2.6.1).  Notice that the variance in \(P,Q\) is contravariant.  The variance of the cofibred category \(\Psi_{L(P,Q)}\), with respect to \(P,Q\), is equally clear, because \(L(P,Q)=L(P)\otimes L(Q)\) depends functorially on the pair \((P,Q)\), whereas \(\Psi_L\) depends contravariantly on \(L\) (3.4).  Thus a first compatibility says that, for every triple of morphisms
\begin{equation}\tag{3.7.1}
        P'\longrightarrow P,
        \qquad Q'\longrightarrow Q,
        \qquad G\longrightarrow G',
\end{equation}
one has a commutative square of functors, up to the canonical isomorphism,
\begin{equation}\tag{3.7.2}
\begin{tikzcd}[column sep=large,row sep=large]
 \Psi_{L(P,Q)}(G) \arrow[r] \arrow[d] & \Psi_{L(P',Q')}(G') \arrow[d] \\
 \Biext(P,Q;G) \arrow[r] & \Biext(P',Q';G').
\end{tikzcd}
\end{equation}
We shall admit this compatibility, since writing down the commutativity isomorphism is no less tedious than evident; it has the same character as the complete explicit construction involved in the equivalence of cofibred categories 3.6.4.

It follows in particular that the isomorphism (3.6.5.1) is functorial not only in \(G\), but also in \(P\) and \(Q\).  Similarly, and more trivially, the isomorphism defined through (3.1.8), as (3.6.5.1) was defined through (3.2.4), gives
\begin{equation}\tag{3.7.4}
 \Biext^0(P,Q;G)\xrightarrow{\sim}
 \Ext^0(P\otimes^L Q,G)\simeq \Hom(P\otimes Q,G).
\end{equation}
This is the same isomorphism already obtained in (2.5.4.1), where it was noted in (2.6) that it is functorial in the three arguments.

The interest of (3.6.5) is that it expresses the geometrically defined group \(\Biext^1(P,Q;G)\) by the simple homological invariant \(\Ext^1(P\otimes^L Q,G)\), which belongs to the usual sequence of invariants \(\Ext^i(P\otimes^L Q,G)\).  The latter sequence is a cohomological functor in each of the three arguments.  Thus a short exact sequence
\begin{equation}\tag{3.7.5.1}
        0\longrightarrow G'\longrightarrow G\longrightarrow G''\longrightarrow 0
\end{equation}
or, respectively,
\begin{equation}\tag{3.7.5.2}
        0\longrightarrow P'\longrightarrow P\longrightarrow P''\longrightarrow 0,
        \quad\hbox{or}\quad
        0\longrightarrow Q'\longrightarrow Q\longrightarrow Q''\longrightarrow 0,
\end{equation}
gives rise to an infinite exact sequence in the \(\Ext^i\)'s.  Its first six terms, involving only \(\Ext^0\)'s and \(\Ext^1\)'s, may be interpreted in terms of biextensions by 3.6.5 and 3.7.4.  We have already indicated the geometric interpretation of the exact sequence deduced from (3.7.5.1), and in particular of the boundary operator
\begin{equation}\tag{3.7.5.3}
 \partial:\Biext^0(P,Q;G'')\longrightarrow \Biext^1(P,Q;G'),
\end{equation}
in the framework of left exact cofibred categories (3.3.6).  One may likewise ask for a geometric interpretation of the six-term exact sequence attached to either of the short exact sequences (3.7.5.2), say the first; equivalently, by 3.7.3, one seeks an interpretation of the boundary homomorphism
\begin{equation}\tag{3.7.5.4}
 \partial:\Biext^0(P',Q;G)\longrightarrow \Biext^1(P'',Q;G).
\end{equation}

\textbf{Proposition 3.7.6.}  The cofibred category, over
\(\GRAB(\cT)^{\circ}\times\GRAB(\cT)^{\circ}\times\GRAB(\cT)\), of biextensions of \((P,Q)\) by \(G\), with \(P,Q,G\) variable, is left exact not only in the argument \(G\), but also in each of the arguments \(P\) and \(Q\).  Moreover the boundary homomorphism (3.7.5.4) deduced from the exact sequence of the \(\Ext^i\)'s attached to (3.7.5.2), through 3.6.5 and 3.7.4, is the boundary homomorphism of the theory of left exact cofibred categories over
\(\mathcal C=\GRAB(\cT)^{\circ}\) recalled in 3.3.6, applied to the cofibred category of biextensions of \((P,Q)\) by \(G\), where \(Q\) and \(G\) are fixed and \(P\) is variable, up to the sign convention of 3.4.9.

We sketch the proof.  By 3.6.4 and the compatibility (3.7.2), biextensions of \((P,Q)\) by \(G\) are interpreted as objects of \(\Psi_{L(P,Q)}(G)\).  Consider the morphism
\[
        L(P,Q)=L(P)\otimes L(Q)\longrightarrow P\otimes L(Q)
\]
induced by the augmentation \(L(P)\to P\).  From 3.5.3 and the flatness of \(L(Q)\) this morphism induces isomorphisms on \(H_0\) and \(H_1\), hence by (3.4.7) an equivalence
\[
        \Psi_{P\otimes L(Q)}\xrightarrow{\sim}\Psi_{L(P,Q)}.
\]
Thus biextensions of \((P,Q)\) by \(G\) may be interpreted as objects of \(\Psi_{P\otimes L(Q)}(G)\).  The exact sequence (3.7.5.2) then gives an exact sequence of chain complexes
\begin{equation}\tag{3.7.6.1}
        0\longrightarrow P'\otimes L(Q)
        \longrightarrow P\otimes L(Q)
        \longrightarrow P''\otimes L(Q)
        \longrightarrow 0.
\end{equation}
The asserted left exactness follows from 3.4.8 applied to this exact sequence.  The direct verification is immediate.  Finally, the homomorphism (3.7.5.4) is the boundary \(\Ext^0\to\Ext^1\) obtained from this exact sequence, and the last assertion is therefore a special case of 3.4.9.

\subsection*{3.8. Biextensions of \((P,Q)\) trivialized on extensions of \(P\) and \(Q\)}

Consider two exact sequences of abelian Groups
\begin{equation}\tag{3.8.1}
        0\longrightarrow L_1\xrightarrow{i}L_0\xrightarrow{u}P\longrightarrow 0,
        \qquad
        0\longrightarrow M_1\xrightarrow{j}M_0\xrightarrow{v}Q\longrightarrow 0.
\end{equation}
We wish to describe the category of biextensions \(E\) of \((P,Q)\) by an abelian Group \(G\) whose inverse image \((u,v)^*(E)\) on \((L_0,M_0)\) is trivial.  Thus one is given a homomorphism \(s\) making the triangle
\begin{equation}\tag{3.8.2}
\begin{tikzcd}[column sep=large,row sep=large]
& E \arrow[d] \arrow[dr,"s"]&\\
& P\times Q & L_0\times M_0 \arrow[l,"u\times v"]
\end{tikzcd}
\end{equation}
commute, and satisfying the additivity condition in each of the two arguments.  Such an \(s\) identifies \((u,v)^*(E)\) with the trivial biextension \(E_0=L_0\times M_0\times G\) of \((L_0,M_0)\) by \(G\).

By the left exactness in \(P,Q\) of the cofibred category of biextensions of \((P,Q)\) by a fixed \(G\) (3.7.6), the datum of \(E\) together with the trivialization \(s\) of its inverse image is essentially equivalent to the datum of two partial trivializations of the trivial biextension \(E_0\) of \((L_0,M_0)\) by \(G\), on \((L_1,M_0)\) and \((L_0,M_1)\), respectively, coinciding on \((L_1,M_1)\).  Since the biextensions induced by \(E_0\) on \((L_1,M_0)\), \((L_0,M_1)\), and \((L_1,M_1)\) are trivial, such partial trivializations are equivalent to homomorphisms, or pairings,
\begin{equation}\tag{3.8.3}
        f:L_1\otimes M_0\longrightarrow G,
        \qquad
        g:L_0\otimes M_1\longrightarrow G,
\end{equation}
and the compatibility condition says simply that these two homomorphisms coincide on \(L_1\otimes M_1\):
\begin{equation}\tag{3.8.4}
        f(x_1\otimes y_1)=g(x_1\otimes y_1).
\end{equation}
In terms of the homomorphism \(s\) of (3.8.2), the forms \(f,g\) are characterized by
\begin{equation}\tag{3.8.5}
\begin{aligned}
 s(x_1,y_0)&=f(x_1\otimes y_0)e_{E/P(y_0)},\\
 s(x_0,y_1)&=g(x_0\otimes y_1)e_{E/P(x_0)}.
\end{aligned}
\end{equation}
If another trivialization \(s'\) is used, then
\begin{equation}\tag{3.8.6}
        s'(x_0,y_0)=s(x_0,y_0)+h(x_0\otimes y_0)
\end{equation}
for an arbitrary homomorphism \(h:L_0\otimes M_0\to G\), and the two forms change to
\begin{equation}\tag{3.8.7}
\begin{aligned}
 f'(x_1\otimes y_0)&=f(x_1\otimes y_0)+h(x_1\otimes y_0),\\
 g'(x_0\otimes y_1)&=g(x_0\otimes y_1)+h(x_0\otimes y_1).
\end{aligned}
\end{equation}
Thus one obtains the following description.

\textbf{Proposition 3.8.8.}  The subgroup of \(\Biext^1(P,Q;G)\) formed by the classes of biextensions whose inverse image by \((u,v)\) is a trivial biextension of \((L_0,M_0)\) - equivalently, the kernel of
\[
        \Biext^1(P,Q;G)\longrightarrow \Biext^1(L_0,M_0;G)
\]
- is canonically isomorphic, by the preceding description, to the quotient of the group of pairs of forms \((f,g)\) as in (3.8.3), satisfying (3.8.4), by the subgroup consisting of those pairs obtained by restricting a form \(h:L_0\otimes M_0\to G\).

We now interpret this quotient homologically, taking into account the canonical isomorphism
\(\Biext^1(P,Q;G)\simeq\Ext^1(P\otimes^L Q,G)\) of 3.6.5.  Regard the data (3.8.1) as two resolution complexes of \(P\) and \(Q\):
\begin{equation}\tag{3.8.9}
        L_\bullet:0\longrightarrow L_1\longrightarrow L_0\longrightarrow 0,
        \qquad
        M_\bullet:0\longrightarrow M_1\longrightarrow M_0\longrightarrow 0.
\end{equation}
Their tensor product is the complex
\[
        0\longrightarrow L_1\otimes M_1
        \xrightarrow{d_2}L_1\otimes M_0+L_0\otimes M_1
        \xrightarrow{d_1}L_0\otimes M_0
        \longrightarrow0.
\]
There is a canonical morphism in the derived category
\[
        P\otimes^L Q\longrightarrow L_\bullet\otimes M_\bullet,
\]
hence a canonical homomorphism
\begin{equation}\tag{3.8.10}
        H^1\bigl(\Hom^\bullet(L_\bullet\otimes M_\bullet,G)\bigr)
        \longrightarrow \Ext^1(P\otimes^L Q,G).
\end{equation}
The group of 1-cocycles of the complex \(\Hom^\bullet(L_\bullet\otimes M_\bullet,G)\) identifies with the group of pairs \((f,g)\) as in (3.8.3) satisfying (3.8.4), by assigning to such a pair the homomorphism
\begin{equation}\tag{3.8.11}
        F:L_1\otimes M_0+L_0\otimes M_1\longrightarrow G,
        \qquad
        F(x_1\otimes y_0+x_0\otimes y_1)=f(x_1\otimes y_0)+g(x_0\otimes y_1).
\end{equation}
Moreover, the coboundary of a 0-cochain \(h:L_0\otimes M_0\to G\) is the pair of its restrictions to \(L_1\otimes M_0\) and \(L_0\otimes M_1\).  Therefore the quotient described in Proposition 3.8.8 is canonically the first member of (3.8.10).

\textbf{Proposition 3.8.12.}  Let \(\xi\in\Ker(\Biext^1(P,Q;G)\to\Biext^1(L_0,M_0;G))\) be described by a pair \((f,g)\) as in 3.8.8.  Under the isomorphism (3.6.5.1), the image of \(\xi\) in \(\Ext^1(P\otimes^L Q,G)\) is the image by (3.8.10) of the class represented by the cocycle (3.8.11).

The verification of this compatibility is left to the reader.

We shall use this especially with
\[
        L_0=\Z[P],\qquad M_0=\Z[Q],
\]
where \(L_1\) and \(M_1\) are the kernels of the canonical homomorphisms \(\Z[P]\to P\) and \(\Z[Q]\to Q\).  There are canonical isomorphisms
\[
        L_0\otimes^L M_0\simeq L_0\otimes M_0\simeq \Z[P\times Q]
\]
since \(L_0,M_0\) are flat, and therefore, by (1.4.5), canonical isomorphisms
\[
        \Ext^i(L_0\otimes^L M_0,G)
        \simeq H^i(P\times Q,G_{P\times Q}).
\]
Consequently, the functor sending a biextension of \((L_0,M_0)\) by \(G\) to the torsor which it induces over \(P\times Q\) is an equivalence of categories.  Thus, for a biextension \(E\) of \((P,Q)\) by \(G\), giving a trivialization of its inverse image on \((L_0,M_0)\) is the same as giving a section of \(E\) as an object over \(P\times Q\), i.e. a trivialization of \(E\) as a torsor under \(G_{P\times Q}\).

Biextensions endowed with this supplementary structure are therefore described by pairs \((f,g)\) satisfying the compatibility (3.8.4), modulo the relation described in 3.8.8.  With the notation of 3.5, one may view \(L_1\) as the cokernel of \(L_2(P)\to L_1(P)\), and similarly for \(M_1\).  This expresses \((f,g)\) as a system of homomorphisms
\[
\begin{cases}
 f':L_1(P)\otimes L_0(Q)=\Z(P\times P\times Q)\longrightarrow G,\\
 g':L_0(P)\otimes L_1(Q)=\Z(P\times Q\times Q)\longrightarrow G,
\end{cases}
\]
i.e. as homomorphisms of sheaves of sets
\[
        P\times P\times Q\longrightarrow G,
        \qquad
        P\times Q\times Q\longrightarrow G,
\]
satisfying five compatibility conditions: four expressing the factorization through \(f\) and \(g\), and the condition (3.8.4).  These \(f'\) and \(g'\) are precisely the partial multiplication data defining a biextension, and the five conditions are the five axioms of biextensions considered in 2.1.

\subsection*{Bibliography for Exposé VII}
\begin{enumerate}[label={[\arabic*]}]
\item M. Demazure, J. Giraud and M. Raynaud, \emph{Schémas abéliens}, Séminaire Fac. Sc. Orsay 1967/68, to appear.
\item J. Giraud, \emph{Cohomologie non abélienne}, thesis, Paris, 1967, to appear.
\item A. Grothendieck, Sur quelques points d'Algèbre homologique, \emph{Tohoku Math. Journal} 9 (1956), 119--221.
\item A. Grothendieck, \emph{Catégories cofibrées additives et Complexe cotangent relatif}, Lecture Notes in Mathematics 79, Springer, 1968.
\item J.-L. Verdier, \emph{Catégories dérivées des catégories abéliennes}, to appear.
\item D. Mumford, \emph{Biextensions of formal groups}, Tata Institute of Fundamental Research, 1968, to appear.
\end{enumerate}

\newpage
\section*{Exposé VIII. Complements on biextensions; general properties of biextensions of group schemes}
\begin{center}
by A. Grothendieck
\end{center}

\subsection*{Summary}
\begin{enumerate}[label=\arabic*.]
\setcounter{enumi}{-1}
\item Introduction.
\item Various special cases on an arbitrary topos.
\item Pairings defined by a biextension.
\item Biextensions of group schemes \((P,Q)\) by \(\Gm\): generalities.
\item Extensions and biextensions of smooth connected group schemes over a field.
\item Extensions and biextensions by constant torsion-free group schemes.
\item Extensions and biextensions by the Néron model of \(\Gm\).
\item Canonical prolongations of extensions and biextensions by \(\Gm\).
\end{enumerate}

\subsection*{0. Introduction}

We continue here the general study of biextensions begun in the preceding exposé.  We first give complements in the case of an arbitrary base topos, developing in particular the formalism of the pairings attached to biextensions, whose importance in the theory of abelian schemes is well known.  We then begin the study of biextensions of general group schemes, the most interesting case for us being biextensions by the multiplicative group \(\Gm\).  As a technical intermediate step we shall also have to study biextensions by other types of group schemes.

The most important result for later use is 7.1.  It gives conditions under which, for \(P\) and \(Q\) smooth over the generic point \(\eta\) of a trait, a biextension of \((P_\eta,Q_\eta)\) by \(\Gm\) extends to a biextension of \((P,Q)\) by \(\Gm\), determined up to unique isomorphism.  This will be convenient in the next exposé, which concerns Néron models of abelian varieties.  It is likely that these developments will be useful elsewhere as well, although the later part of the seminar will use them only sparingly.

When working with group schemes over a scheme \(S\), we shall tacitly use the fppf topos of \(S\), attached to the \(\cU\)-site of schemes locally of finite presentation over \(S\), endowed with the fppf topology.  Most statements would remain valid, and essentially equivalent, for other topologies whose underlying category contains the schemes locally of finite presentation over \(S\), such as the Zariski, étale, or fpqc topology.  The reason is descent: the category of extensions of a group scheme \(P\), locally of finite presentation over \(S\), by \(\Gm\), and similarly the notion of biextension, do not depend, up to equivalence, on this choice.  Technically, the fppf topology, or a finer topology such as fpqc, is preferred because the homomorphism \(n\cdot\id\) of \(\Gm\) is an fppf epimorphism for every \(n>0\), whereas it is not in general a Zariski or étale epimorphism.

Finally, most statements on extensions or biextensions by \(\Gm\) remain true if \(\Gm\) is replaced by a torus \(G\).  Since such a \(G\) is locally isomorphic, for the fppf and even for the étale topology, to \(\Gm^r\), this is essentially a trivial generalization.

\section*{1. Various special cases on an arbitrary topos}

As in Exposé VII, let \(\cT\) be a topos, and let \(P,Q,G\) be three abelian Groups of \(\cT\).  We use the canonical isomorphism
\begin{equation}\tag{1.1.1}
        \Biext^1(P,Q;G)\simeq \Ext^1(P\otimes^L Q,G).
\end{equation}
Recall also the Cartan isomorphisms
\begin{equation}\tag{1.1.2}
 \Ext^1(P\otimes^L Q,G)
 \simeq \Ext^1(P,\RHom(Q,G))
 \simeq \Ext^1(Q,\RHom(P,G)),
\end{equation}
functorial in the three arguments.  From (1.1.1) one obtains the exact sequence
\begin{equation}\tag{1.1.3}
\begin{gathered}
0\to \Ext^1(P\otimes Q,G)\to \Biext^1(P,Q;G)
\to \Hom(\Torop_1(P,Q),G)\to{}\\
\Ext^2(P\otimes Q,G)\to \Ext^2(P\otimes^L Q,G)\to\cdots,
\end{gathered}
\end{equation}
because the homology objects of \(P\otimes^L Q\) are the \(\Torop_i(P,Q)\), zero for \(i\ge1\).  Similarly, since the homology objects of \(\RHom(Q,G)\) are the \(\Ext^i(Q,G)\), the first isomorphism in (1.1.2) gives a five-term exact sequence
\begin{equation}\tag{1.1.4}
\begin{gathered}
0\to \Ext^1(P,\Hom(Q,G))\to \Biext^1(P,Q;G)
\to \Hom(P,\Ext^1(Q,G))\to{}\\
\Ext^2(P,\Hom(Q,G))\to \Ext^2(P,\RHom(Q,G)).
\end{gathered}
\end{equation}
This sequence is associated with the usual spectral sequence
\[
        \Ext^*(P,\RHom(Q,G))\Longleftarrow
        \Ext^p(P,\Ext^q(Q,G)).
\]
A fully faithful functor may be used to describe the first arrow in (1.1.3), respectively in (1.1.4), namely
\begin{equation}\tag{1.1.5}
        \operatorname{EXT}(P\otimes Q,G)\to\operatorname{BIEXT}(P,Q;G),
\end{equation}
or
\begin{equation}\tag{1.1.6}
        \operatorname{EXT}(P,\Hom(Q,G))\to\operatorname{BIEXT}(P,Q;G).
\end{equation}
The symmetric Cartan isomorphism yields the analogous functor
\begin{equation}\tag{1.1.7}
        \operatorname{EXT}(Q,\Hom(P,G))\to\operatorname{BIEXT}(P,Q;G).
\end{equation}

If \(\Torop_1(P,Q)=0\), then (1.1.3) gives a canonical isomorphism
\begin{equation}\tag{1.2.2}
        \Biext^1(P,Q;G)\simeq \Ext^1(P\otimes Q,G).
\end{equation}
This is the case, for example, if \(P\) or \(Q\) is a flat \(\mathbb Z\)-Module; more precisely, under the same hypothesis the functor (1.1.5) is an equivalence of categories.

If \(P\otimes Q=0\), then (1.1.3) gives a canonical isomorphism
\begin{equation}\tag{1.3.2}
        \Biext^1(P,Q;G)\simeq \Hom(\Torop_1(P,Q),G).
\end{equation}
A particularly useful case occurs when there is an integer \(n>0\) such that
\begin{equation}\tag{1.3.3}
        nP=P,
        \qquad
        Q=\liminj_i {}_{n^i}Q,
\end{equation}
that is, \(P\) is divisible by \(n\) and \(Q\) is a filtered limit of subgroups killed by powers of \(n\).  Put
\[
        T_n(P)=({}_{n^i}P)_{i\ge0},
        \qquad
        {}_{n^\infty}P=\liminj_i {}_{n^i}P.
\]
Then
\begin{equation}\tag{1.3.4}
        P\otimes Q=0,
        \qquad
        \Torop_1(P,Q)\simeq T_n(P)\otimes Q
        \simeq {}_{n^\infty}P\otimes T_n(Q),
\end{equation}
and hence
\begin{equation}\tag{1.3.5}
        \Biext^1(P,Q;G)
        \simeq \Hom(T_n(P),D(Q))
        \simeq \Hom(Q,D(T_n(P))).
\end{equation}
Here \(D\) denotes \(\Hom(-,G)\).  The transition map in the inductive system appearing in (1.3.4) is obtained by identifying the first term with \({}_{n^j}P\otimes{}_{n^i}Q\), using \(n^{j-i}\id\) on the first factor and the inclusion on the second.  For a group \(M\) killed by \(n^i\) one defines
\[
        T_n(P)\otimes M\dfn {}_{n^i}P\otimes M,
\]
and \(D(Q)=(\Hom({}_{n^i}Q,G))_{i\ge0}\).  The remaining identifications are those of morphisms of projective systems.

The calculation of \(\Torop_1(P,Q)\) used above is reduced to the case in which \(Q\) is killed by \(n^i\).  It suffices to use the exact sequence
\[
        0\to P\xrightarrow{m}P\to P/mP\to0
\]
and the standard connecting morphism to get, when \(mP=P\) and \(mQ=0\),
\begin{equation}\tag{1.3.6}
        \Torop_1(P,Q)\xleftarrow{\sim}{}_mP\otimes_{\mathbb Z/m\mathbb Z} Q.
\end{equation}
The choice of this isomorphism amounts only to a sign convention.

If
\begin{equation}\tag{1.4.1}
        \Hom(Q,G)=0,
\end{equation}
then (1.1.4) gives a canonical isomorphism
\begin{equation}\tag{1.4.2}
        \Biext^1(P,Q;G)\simeq\Hom(P,\Ext^1(Q,G)).
\end{equation}
This case occurs, for example, when \(Q\) is an abelian scheme and \(G=\Gm\).  Then the functor in \(P\), \(P\mapsto\Biext^1(P,Q;G)\), is representable by
\[
        Q^*=\Ext^1(Q,G),
\]
and one obtains a canonical biextension \(E^Q\) of \((Q^*,Q)\) by \(G\), uniquely determined up to unique isomorphism.  This biextension, or the symmetric biextension \({}^sE^Q\) of \((Q,Q^*)\) by \(G\), may be called the Weil biextension of \((Q^*,Q)\), or of \((Q,Q^*)\), by \(G\).

If
\begin{equation}\tag{1.5.1}
        \Hom(P,\Ext^1(Q,G))=0,
\end{equation}
then (1.1.4) yields the canonical isomorphism
\begin{equation}\tag{1.5.2}
        \Biext^1(P,Q;G)\simeq \Ext^1(P,\Hom(Q,G)).
\end{equation}
More precisely, the fully faithful functor (1.1.6) is then an equivalence of categories.  If both
\begin{equation}\tag{1.6.1}
        \Hom(P,G)=0,
        \qquad
        \Hom(P,\Ext^1(Q,G))=0
\end{equation}
hold, then the two isomorphisms of 1.4 combine into a duality formula
\begin{equation}\tag{1.6.2}
        \Ext^1(P,\Hom(Q,G))
        \simeq \Hom(Q,\Ext^1(P,G))
        \simeq \Biext^1(P,Q;G).
\end{equation}
The first condition in (1.6.1) implies \(\Hom(P\otimes Q,G)=0\).  Hence the category of commutative extensions of \(P\) by \(\Hom(Q,G)\) is rigid, and the same is true of \(\operatorname{BIEXT}(P,Q;G)\).  Therefore their structure is known once the group of isomorphism classes is known; by (1.6.2) this group is canonically \(\Hom(Q,\Ext^1(P,G))\).

The map (1.6.2) is actually defined without assuming (1.6.1).  It may be described directly, without passing through \(\Biext^1(P,Q;G)\), in terms of the natural pairing
\[
        Q\times Q'\longrightarrow G,
        \qquad Q'=\Hom(Q,G),
\]
which is equivalently a homomorphism \(Q\to\Hom(Q',G)\).  It gives the composite
\begin{equation}\tag{1.6.3}
        \Ext^1(P,Q')\times Q
        \longrightarrow \Ext^1(P,Q')\times\Hom(Q',G)
        \longrightarrow \Ext^1(P,G),
\end{equation}
recovering (1.6.2).

\section*{2. Pairings defined by a biextension}

Let \(P,Q\) be two abelian Groups of \(\cT\), and let \(n>0\) be an integer.  We write \({}_nP\) for the kernel of multiplication by \(n\) on \(P\).  We define a canonical homomorphism
\begin{equation}\tag{2.1.1}
        \alpha=\alpha^n_{P,Q}:{}_nP\otimes{}_nQ\longrightarrow \Torop_1(P,Q).
\end{equation}
It is the composite
\begin{equation}\tag{2.1.2}
        {}_nP\otimes{}_nQ
        \longrightarrow \Torop_1({}_nP,{}_{n}Q)
        \longrightarrow \Torop_1(P,Q),
\end{equation}
where the first arrow is defined first in the case in which \(P\) and \(Q\) are killed by \(n\).  Suppose for instance that \({}_nQ=Q\).  Taking a flat resolution \(L_\bullet\) of \(P\), one has canonical isomorphisms
\[
        P\otimes^L_{\mathbb Z}Q
        \simeq L_\bullet\otimes Q
        \simeq (L_\bullet\otimes \mathbb Z/n\mathbb Z)\otimes_{\mathbb Z/n\mathbb Z}Q,
\]
and this gives a homological Kunneth spectral sequence
\begin{equation}\tag{2.1.4}
        \Torop_*^{\mathbb Z}(P,Q)
        \Longleftarrow
        E^2_{p,q}=\Torop_p^{\mathbb Z/n\mathbb Z}
        (\Torop_q^{\mathbb Z}(P,\mathbb Z/n\mathbb Z),Q).
\end{equation}
The usual flat resolution
\[
        0\longrightarrow\mathbb Z_\cT\xrightarrow{n}\mathbb Z_\cT\longrightarrow0
\]
of \((\mathbb Z/n\mathbb Z)_\cT\) gives
\[
        \Torop_0^{\mathbb Z}(P,\mathbb Z/n\mathbb Z)=P_n=P/nP,
        \qquad
        \Torop_1^{\mathbb Z}(P,\mathbb Z/n\mathbb Z)={}_nP,
\]
and higher \(\Torop_i\)'s vanish.  Therefore (2.1.4) reduces essentially to the exact sequence
\begin{equation}\tag{2.1.5}
0\to \Torop_2^{\mathbb Z/n\mathbb Z}(P_n,Q)\to {}_nP\otimes Q
\xrightarrow{\alpha}\Torop_1^{\mathbb Z}(P,Q)\to\Torop_1^{\mathbb Z/n\mathbb Z}(P_n,Q)\to\cdots .
\end{equation}
This defines the desired map and shows that it is functorial in \(P\) and \(Q\).  In particular it is an isomorphism when \({}_nP\) is a flat \(\mathbb Z/n\mathbb Z\)-Module (for example when \(nP=P\) and \({}_nP\) is flat), or when \(Q\) is flat over \(\mathbb Z/n\mathbb Z\).

The map \(\alpha\) may be described more explicitly.  To know \(\alpha(p,q)\) for sections \(p\in{}_nP\) and \(q\in{}_nQ\), one reduces by functoriality to the case \(Q=(\mathbb Z/n\mathbb Z)_\cT\) and \(q=1\).  In that case \(\Torop_1^{\mathbb Z}(P,Q)\simeq{}_nP\), and \(\alpha\) is the identity.  Equivalently, one further reduces to \(P=Q=(\mathbb Z/n\mathbb Z)_\cT\), where
\begin{equation}\tag{2.1.8.1}
        (\mathbb Z/n\mathbb Z)_\cT
        \xrightarrow{\sim}\Torop_1^{\mathbb Z}\bigl((\mathbb Z/n\mathbb Z)_\cT,(\mathbb Z/n\mathbb Z)_\cT\bigr)
\end{equation}
is induced by the canonical flat resolution
\begin{equation}\tag{2.1.8.2}
        0\longrightarrow\mathbb Z_\cT\xrightarrow{n}\mathbb Z_\cT\longrightarrow0
\end{equation}
of the first factor.

The definition of \(\alpha\) is not symmetric.  By exchanging \(P\) and \(Q\) in the construction, one obtains an analogous map
\begin{equation}\tag{2.1.9.1}
        \beta^n_{P,Q}:{}_nP\otimes{}_nQ\longrightarrow\Torop_1(P,Q).
\end{equation}
The two maps may also be characterized by the commutative square
\begin{equation}\tag{2.1.9.2}
\begin{tikzcd}[column sep=large,row sep=large]
{}_nP\otimes{}_nQ \arrow[r,"\sim"] \arrow[d,"\beta^n_{P,Q}"']
  & {}_nQ\otimes{}_nP \arrow[d,"\alpha^n_{Q,P}"]\\
\Torop_1(P,Q) \arrow[r,"\sim"'] & \Torop_1(Q,P),
\end{tikzcd}
\end{equation}
where the horizontal arrows are the symmetry isomorphisms.  Their relation is
\begin{equation}\tag{2.1.9.3}
        \beta^n_{P,Q}=-\alpha^n_{P,Q}.
\end{equation}

\textbf{Lemma 2.1.10.}  Let \(P,Q\) be two commutative Groups of \(\cT\), and let \(n>0\).  Then the square
\begin{equation}\tag{2.1.10.1}
\begin{tikzcd}[column sep=large,row sep=large]
{}_nP\otimes{}_nQ \arrow[r,"\sim"] \arrow[d,"\alpha^n_{P,Q}"']
  & {}_nQ\otimes{}_nP \arrow[d,"\alpha^n_{Q,P}"]\\
\Torop_1(P,Q) \arrow[r,"\sim"'] & \Torop_1(Q,P)
\end{tikzcd}
\end{equation}
is anti-commutative, the horizontal arrows being the symmetry isomorphisms.

\emph{Proof.}  The procedure of 2.1.8 reduces the question to the case \(P=Q=(\mathbb Z/n\mathbb Z)_\cT\).  It is enough then to work in the punctual topos, i.e. in ordinary commutative groups.  Take
\[
        L_\bullet:0\longrightarrow\mathbb Z\xrightarrow{n}\mathbb Z\longrightarrow0
\]
as a flat resolution of \(P=Q=\mathbb Z/n\mathbb Z\).  Using \(L_\bullet\otimes L_\bullet\), the isomorphism defining \(\alpha\) is represented by the cycle
\[
        -1\otimes l_1+l_1\otimes1,
\]
whereas the isomorphism defining \(\beta\) is represented by
\[
        1\otimes l_1-l_1\otimes1.
\]
This proves the lemma.

\end{document}
